\coursepart{1}{(MATH203) Applied Linear Algebra}
\ifCourseSolutions\else
{\noindent\small\itshape Problems are arranged chronologically by year and semester, with the midterm preceding the final examination within the same semester, while preserving the original problem order within each source.\par}
\vspace{0.22cm}
\fi

% 2018 Fall POSTECH Applied Linear Algebra Midterm, Problem 1
\begin{problem}[2018 Fall \POSTECH{} Applied Linear Algebra]
In this problem, $V$ is a vector space. Determine whether each of the following statements is true or false.
\begin{problemchoices}
\item There exists a linearly independent subset $S\subseteq V$ that contains the zero vector $\mathbf 0$.
\item If $\mathbf A\in\R^{n\times n}$ is invertible, then the columns of $\mathbf A$ form a basis for $\R^n$.
\item Every projection matrix is invertible.
\item $\Null(\mathbf A)^\perp=\Col(\mathbf A)$.
\item Let $\{\mathbf v_1,\dots,\mathbf v_n\}$ be a basis for $V$, and let $T\colon V\to W$ be a linear transformation. If
\[
T(\mathbf v_i)=\mathbf w_i,
\qquad 1\le i\le n,
\]
for given vectors $\mathbf w_1,\dots,\mathbf w_n\in W$, then $T$ is uniquely determined.
\end{problemchoices}
\end{problem}

% 2018 Fall POSTECH Applied Linear Algebra Midterm, Problem 2
\begin{problem}[2018 Fall \POSTECH{} Applied Linear Algebra]
For each statement, prove it or disprove it by providing a counterexample.
\begin{problemchoices}
\item Consider a linear system
\[
\mathbf A\mathbf x=\mathbf b,
\qquad
\mathbf A\in\R^{m\times n},\quad
\mathbf x\in\R^n,\quad
\mathbf b\in\R^m.
\]
If there are more variables than equations, $m<n$, then the system always has a solution.
\item If $\mathbf A\in\R^{n\times n}$ satisfies
\[
\mathbf A^2=\mathbf A
\qquad\text{and}\qquad
\rk(\mathbf A)=n,
\]
then $\mathbf A=\mathbf I_n$.
\item The null space of $\mathbf A$ is equal to the null space of $\mathbf A^{\mathsf T}\mathbf A$.
\end{problemchoices}
\end{problem}

% 2018 Fall POSTECH Applied Linear Algebra Midterm, Problem 3
\begin{problem}[2018 Fall \POSTECH{} Applied Linear Algebra]
Let
\[
\mathbf A=
\begin{pmatrix}
1&1&2\\
0&1&0\\
2&1&1
\end{pmatrix}.
\]
\begin{problemsubparts}
\item Find $\mathbf A^{-1}$.
\item Compute $\bigl(\mathbf A\mathbf A^{\mathsf T}\bigr)^{-1}$.
\end{problemsubparts}
\end{problem}

% 2018 Fall POSTECH Applied Linear Algebra Midterm, Problem 4
\begin{problem}[2018 Fall \POSTECH{} Applied Linear Algebra]
Let
\[
\mathbf A=
\begin{pmatrix}
1&1&0\\
0&1&1
\end{pmatrix}.
\]
\begin{problemsubparts}
\item Find a right inverse $\mathbf B$ such that
\[
\mathbf A\mathbf B=\mathbf I_2.
\]
\item Using the matrix $\mathbf B$ obtained in part \textup{(a)}, compute $\bigl(\mathbf B\mathbf A\bigr)^{100}$.
\end{problemsubparts}
\end{problem}

% 2018 Fall POSTECH Applied Linear Algebra Midterm, Problem 5
\begin{problem}[2018 Fall \POSTECH{} Applied Linear Algebra]
Let
\[
\mathbf A=
\begin{pmatrix}
1&1\\
4&4\\
4&4
\end{pmatrix},
\qquad
\mathbf b=
\begin{pmatrix}
1\\1\\1
\end{pmatrix}.
\]
Find the orthogonal projection $\mathbf b_c$ of $\mathbf b$ onto $\Col(\mathbf A)$.
\end{problem}

% 2018 Fall POSTECH Applied Linear Algebra Midterm, Problem 6
\begin{problem}[2018 Fall \POSTECH{} Applied Linear Algebra]
Let
\[
\mathbf A=
\begin{pmatrix}
1&1&1\\
1&2&2\\
1&2&3
\end{pmatrix},
\qquad
\mathbf B=
\begin{pmatrix}
0&1&1&1\\
-1&0&-1&0\\
1&-1&0&0
\end{pmatrix},
\qquad
\mathbf C=
\begin{pmatrix}
1&2&2\\
2&4&5
\end{pmatrix}.
\]
\begin{problemsubparts}
\item Find an $LU$ decomposition $\mathbf A=\mathbf L\mathbf U$, where $\mathbf L$ is lower triangular and $\mathbf U$ is upper triangular.
\item Find a basis for $\Col(\mathbf B)$.
\item Find $\Null(\mathbf B)$.
\item Write the complete solution in the form $\mathbf x=\mathbf x_p+\mathbf x_n$ for
\[
\mathbf C
\begin{pmatrix}
x\\y\\z
\end{pmatrix}
=
\begin{pmatrix}
1\\4
\end{pmatrix},
\]
where $\mathbf x_p$ is a particular solution and $\mathbf x_n\in\Null(\mathbf C)$.
\end{problemsubparts}
\end{problem}

% 2018 Fall POSTECH Applied Linear Algebra Midterm, Problem 7
\begin{problem}[2018 Fall \POSTECH{} Applied Linear Algebra]
Let $V=M_2(\R)$, the vector space of $2\times2$ real matrices with the usual addition and scalar multiplication.
\begin{problemsubparts}
\item Show that
\[
\mathbf v_1=
\begin{pmatrix}
1&0\\0&1
\end{pmatrix},
\quad
\mathbf v_2=
\begin{pmatrix}
1&0\\0&-1
\end{pmatrix},
\quad
\mathbf v_3=
\begin{pmatrix}
0&1\\1&0
\end{pmatrix},
\quad
\mathbf v_4=
\begin{pmatrix}
0&1\\-1&0
\end{pmatrix}
\]
form a basis for $V$.
\item Find the matrix representation of the linear transformation
\[
T\colon V\to V,
\qquad
\begin{pmatrix}
a&b\\c&d
\end{pmatrix}
\mapsto
\begin{pmatrix}
a+d&c\\b&a-d
\end{pmatrix},
\]
with respect to the ordered basis $(\mathbf v_1,\mathbf v_2,\mathbf v_3,\mathbf v_4)$.
\end{problemsubparts}
\end{problem}

% 2018 Fall POSTECH Applied Linear Algebra Midterm, Problem 8
\begin{problem}[2018 Fall \POSTECH{} Applied Linear Algebra]
Let
\[
P_k=\left\{a_k t^k+a_{k-1}t^{k-1}+\cdots+a_1t+a_0\mid a_0,\dots,a_k\in\R\right\},
\]
with the usual polynomial addition and scalar multiplication. Use the ordered basis
\[
(1,t,t^2,\dots,t^k)
\]
for $P_k$.
\begin{problemsubparts}
\item Find the matrix representation $\mathbf A$ of the differentiation map
\[
D\colon P_2\to P_1,
\qquad
D(f)=\frac{\dd f}{\dd t}.
\]
\item Find the matrix representation $\mathbf B$ of the integration map
\[
T\colon P_1\to P_2,
\qquad
(Tf)(t)=\int_0^t f(s)\,\dd s.
\]
\end{problemsubparts}
\end{problem}

% 2018 Fall POSTECH Applied Linear Algebra Final, Problem 1
\begin{problem}[2018 Fall \POSTECH{} Applied Linear Algebra]
Determine whether each of the following statements is true or false.
\begin{problemchoices}
\item $\mathbf A$ and $\mathbf A^{\mathsf T}\mathbf A$ have the same rank for every rectangular matrix $\mathbf A$.
\item The only possible eigenvalues of a symmetric orthogonal matrix are $1$ and $-1$.
\item For an $n\times n$ matrix $\mathbf A$ and a scalar $k\in\R$,
\[
\det(k\mathbf A)=k\det(\mathbf A).
\]
\item For an orthogonal matrix $\mathbf Q\in\R^{n\times n}$ and $\mathbf x\in\R^n$,
\[
\|\mathbf Q\mathbf x\|=\|\mathbf x\|.
\]
\item If there is a basis consisting of eigenvectors of a matrix, then the matrix is diagonalizable.
\item Every diagonalizable matrix has an orthonormal basis consisting of eigenvectors.
\end{problemchoices}
\end{problem}

% 2018 Fall POSTECH Applied Linear Algebra Final, Problem 2
\begin{problem}[2018 Fall \POSTECH{} Applied Linear Algebra]
Compute the determinants of the following matrices.
\begin{problemsubparts}
\item Use cofactor expansion to find the determinant of
\[
\mathbf A=
\begin{pmatrix}
1&0&2\\
0&2&-1\\
-1&1&3
\end{pmatrix}.
\]
\item Transform the matrix into upper triangular form using elementary row operations, and compute the determinant of
\[
\mathbf B=
\begin{pmatrix}
1&0&2\\
-1&2&-1\\
2&4&0
\end{pmatrix}.
\]
\end{problemsubparts}
\end{problem}

% 2018 Fall POSTECH Applied Linear Algebra Final, Problem 3
\begin{problem}[2018 Fall \POSTECH{} Applied Linear Algebra]
Let $\mathbf A$ and $\mathbf B$ be real square matrices. Prove that if $\mathbf A$ and $\mathbf B$ are orthogonally similar, then they have the same singular values.
\end{problem}

% 2018 Fall POSTECH Applied Linear Algebra Final, Problem 4
\begin{problem}[2018 Fall \POSTECH{} Applied Linear Algebra]
Write down as many pairwise nonsimilar $4\times4$ matrices as possible whose only eigenvalue is $2$.
\end{problem}

% 2018 Fall POSTECH Applied Linear Algebra Final, Problem 5
\begin{problem}[2018 Fall \POSTECH{} Applied Linear Algebra]
Consider $\R^4$ with the standard inner product and let
\[
S=\{\mathbf v_1,\mathbf v_2,\mathbf v_3,\mathbf v_4\}\subseteq\R^4,
\]
where
\[
\mathbf v_1=
\begin{pmatrix}
1\\0\\1\\0
\end{pmatrix},
\qquad
\mathbf v_2=
\begin{pmatrix}
0\\1\\0\\-1
\end{pmatrix},
\qquad
\mathbf v_3=
\begin{pmatrix}
2\\3\\2\\-3
\end{pmatrix},
\qquad
\mathbf v_4=
\begin{pmatrix}
1\\2\\1\\2
\end{pmatrix}.
\]
\begin{problemsubparts}
\item Find a subset of $S$ that is a basis for $\Span(S)$ and includes $\{\mathbf v_1,\mathbf v_2\}$.
\item Apply the Gram--Schmidt process to the subset found in part \textup{(a)} and obtain an orthonormal basis for $\Span(S)$.
\end{problemsubparts}
\end{problem}

% 2018 Fall POSTECH Applied Linear Algebra Final, Problem 6
\begin{problem}[2018 Fall \POSTECH{} Applied Linear Algebra]
Identify and sketch the curve
\[
x^2+xy+y^2=1.
\]
Justify your identification with an explicit equation. Indicate all $x$-intercepts and $y$-intercepts, and label all coordinate axes.
\end{problem}

% 2018 Fall POSTECH Applied Linear Algebra Final, Problem 7
\begin{problem}[2018 Fall \POSTECH{} Applied Linear Algebra]
Solve the initial value problem
\[
\mathbf u'(t)=
\begin{pmatrix}
1&2\\
0&1
\end{pmatrix}\mathbf u(t),
\qquad
\mathbf u(0)=
\begin{pmatrix}
0\\1
\end{pmatrix}.
\]
\end{problem}

% 2018 Fall POSTECH Applied Linear Algebra Final, Problem 8
\begin{problem}[2018 Fall \POSTECH{} Applied Linear Algebra]
Use a similarity transformation to find the matrix $\mathbf A\in\R^{3\times3}$ representing reflection across the plane
\[
x+y+z=0.
\]
\end{problem}

% 2018 Fall POSTECH Applied Linear Algebra Final, Problem 9
\begin{problem}[2018 Fall \POSTECH{} Applied Linear Algebra]
Find the singular value decomposition of
\[
\mathbf A=
\begin{pmatrix}
4&0&3\\
0&0&5
\end{pmatrix}.
\]
\end{problem}

% 2019 Spring POSTECH Applied Linear Algebra Final, Problem 1
\begin{problem}[2019 Spring \POSTECH{} Applied Linear Algebra]
Let
\[
\mathbf A=
\begin{pmatrix}
1&0&1\\
0&1&0\\
1&0&3
\end{pmatrix}.
\]
\begin{problemsubparts}
\item Find an orthogonal matrix $\mathbf P$ such that $\mathbf P^{\mathsf T}\mathbf A\mathbf P$ is diagonal.
\item Find the characteristic polynomial of $\mathbf A$.
\item Show that $\mathbf A^{-1}$ can be expressed as
\[
a\mathbf A^2+b\mathbf A+c\mathbf I_3
\]
for some $a,b,c\in\R$, and find $a,b,c$.
\item Find $e^{\mathbf A}$.
\end{problemsubparts}
\end{problem}

% 2019 Spring POSTECH Applied Linear Algebra Final, Problem 2
\begin{problem}[2019 Spring \POSTECH{} Applied Linear Algebra]
Let
\[
\mathbf A=
\begin{pmatrix}
0&1&2\\
0&0&1\\
0&0&1
\end{pmatrix}.
\]
\begin{problemsubparts}
\item Find all eigenvalues of $\mathbf A$.
\item Find all eigenvectors of $\mathbf A$.
\item Determine whether $\mathbf A$ is diagonalizable. Explain why.
\item Find the Jordan form of $\mathbf A$.
\item Find a Jordan basis of $\mathbf A$.
\item Find $\det(e^{\mathbf A})$.
\item Calculate $e^{\mathbf A}$.
\item Solve
\[
\mathbf y'=\mathbf A\mathbf y,
\qquad
\mathbf y(0)=
\begin{pmatrix}
2\\1\\4
\end{pmatrix}.
\]
\end{problemsubparts}
\end{problem}

% 2019 Spring POSTECH Applied Linear Algebra Final, Problem 3
\begin{problem}[2019 Spring \POSTECH{} Applied Linear Algebra]
Let
\[
\mathbf D=
\begin{pmatrix}
\mathbf A&\mathbf B\\
\mathbf 0&\mathbf C
\end{pmatrix},
\]
where $\mathbf A$ and $\mathbf C$ are square matrices and $\mathbf 0$ is a zero matrix of compatible size. Prove that
\[
\det(\mathbf D)=\det(\mathbf A)\det(\mathbf C).
\]
\end{problem}

% 2019 Spring POSTECH Applied Linear Algebra Final, Problem 4
\begin{problem}[2019 Spring \POSTECH{} Applied Linear Algebra]
Let $(g_k)$ be a sequence satisfying
\[
g_{k+2}=g_{k+1}-g_k,
\qquad
g_1=1,
\qquad
g_2=2.
\]
\begin{problemsubparts}
\item Find $g_{300}$.
\item Determine whether the sequence $(g_k)$ converges. In either case, explain why.
\end{problemsubparts}
\end{problem}

% 2019 Spring POSTECH Applied Linear Algebra Final, Problem 5
\begin{problem}[2019 Spring \POSTECH{} Applied Linear Algebra]
Calculate
\[
\det
\begin{pmatrix}
0&1&1&1\\
1&0&1&1\\
1&1&0&1\\
1&1&1&0
\end{pmatrix}.
\]
\end{problem}

% 2019 Spring POSTECH Applied Linear Algebra Final, Problem 6
\begin{problem}[2019 Spring \POSTECH{} Applied Linear Algebra]
Let
\[
f(x)=
\det
\begin{pmatrix}
1&a&a^2&a^3\\
1&b&b^2&b^3\\
1&c&c^2&c^3\\
1&x&x^2&x^3
\end{pmatrix}.
\]
Find all zeros of $f(x)$ and verify your answer.
\end{problem}

% 2019 Spring POSTECH Applied Linear Algebra Final, Problem 7
\begin{problem}[2019 Spring \POSTECH{} Applied Linear Algebra]
Let $\mathbf M\in\R^{n\times n}$ be invertible, and let $\mathbf v_1,\dots,\mathbf v_l\in\R^n$ be linearly independent. Prove that
\[
\mathbf M\mathbf v_1,\dots,\mathbf M\mathbf v_l
\]
are also linearly independent.
\end{problem}

% 2019 Spring POSTECH Applied Linear Algebra Final, Problem 8
\begin{problem}[2019 Spring \POSTECH{} Applied Linear Algebra]
Prove that the two matrices
\[
\begin{pmatrix}
\mathbf A_1&\mathbf 0&\mathbf 0\\
\mathbf 0&\mathbf A_2&\mathbf 0\\
\mathbf 0&\mathbf 0&\mathbf A_3
\end{pmatrix}
\qquad\text{and}\qquad
\begin{pmatrix}
\mathbf A_2&\mathbf 0&\mathbf 0\\
\mathbf 0&\mathbf A_3&\mathbf 0\\
\mathbf 0&\mathbf 0&\mathbf A_1
\end{pmatrix}
\]
are similar, where $\mathbf A_1,\mathbf A_2,\mathbf A_3$ are square matrices of sizes $m,n,l$, respectively.
\end{problem}

% 2019 Spring POSTECH Applied Linear Algebra Final, Problem 9
\begin{problem}[2019 Spring \POSTECH{} Applied Linear Algebra]
Let
\[
\mathbf A=
\begin{pmatrix}
1&0&-1+i\\
0&1&1-i\\
0&0&i
\end{pmatrix}.
\]
\begin{problemsubparts}
\item Find the eigenvalues of $\mathbf A$.
\item Find the eigenvectors of $\mathbf A$.
\item Find a matrix $\mathbf S$ such that $\mathbf S^{-1}\mathbf A\mathbf S$ is diagonal.
\end{problemsubparts}
\end{problem}

% 2019 Spring POSTECH Applied Linear Algebra Final, Problem 10
\begin{problem}[2019 Spring \POSTECH{} Applied Linear Algebra]
Let $\mathbf A$ be an $n\times n$ real symmetric matrix, and let $\lambda$ be an eigenvalue of $\mathbf A$ with algebraic multiplicity $m$. Show that
\[
\dim\Null(\mathbf A-\lambda\mathbf I_n)=m.
\]
\end{problem}

% 2019 Spring POSTECH Applied Linear Algebra Final, Problem 11
\begin{problem}[2019 Spring \POSTECH{} Applied Linear Algebra]
\begin{problemchoices}
\item State and prove Schur's Lemma.
\item State and prove the complex spectral theorem.
\end{problemchoices}
\end{problem}

% 2019 Fall POSTECH Applied Linear Algebra Midterm, Problem 1
\begin{problem}[2019 Fall \POSTECH{} Applied Linear Algebra]
Let $\mathbf A=(a_{ij})$ be the $n\times n$ matrix defined by
\[
a_{ij}=
\begin{cases}
1, & i-j=1\text{ and }i<n,\\
0, & \text{otherwise}.
\end{cases}
\]
Find the sum of all entries of
\[
\sum_{k=1}^{2n}\mathbf A^k.
\]
\end{problem}

% 2019 Fall POSTECH Applied Linear Algebra Midterm, Problem 2
\begin{problem}[2019 Fall \POSTECH{} Applied Linear Algebra]
Let
\[
\mathbf A=
\begin{pmatrix}
a&0&b\\
a&a&4\\
0&a&2
\end{pmatrix}.
\]
For what values of $a$ and $b$ does the system
\[
\mathbf A\mathbf x=
\begin{pmatrix}
2\\4\\b
\end{pmatrix}
\]
have
\begin{problemchoices}
\item a unique solution,
\item a one-parameter family of solutions,
\item a two-parameter family of solutions,
\item no solution?
\end{problemchoices}
\end{problem}

% 2019 Fall POSTECH Applied Linear Algebra Midterm, Problem 3
\begin{problem}[2019 Fall \POSTECH{} Applied Linear Algebra]
Find the inverse of
\[
\mathbf B=
\begin{pmatrix}
1&2&-3&0\\
1&-2&0&4\\
-1&2&0&4\\
1&2&3&0
\end{pmatrix}.
\]
\end{problem}

% 2019 Fall POSTECH Applied Linear Algebra Midterm, Problem 4
\begin{problem}[2019 Fall \POSTECH{} Applied Linear Algebra]
Let $\{\mathbf e_1,\dots,\mathbf e_n\}$ and $\{\mathbf f_1,\dots,\mathbf f_n\}$ be two bases of a vector space $V$. If
\[
\mathbf e_i=\sum_{j=1}^n a_{ij}\mathbf f_j,
\qquad i=1,\dots,n,
\]
show that the matrix $\mathbf A=(a_{ij})$ is invertible.
\end{problem}

% 2019 Fall POSTECH Applied Linear Algebra Midterm, Problem 5
\begin{problem}[2019 Fall \POSTECH{} Applied Linear Algebra]
Let
\[
\mathbf A=
\begin{pmatrix}
1&1&5&1&1\\
1&1&3&0&-1\\
2&2&12&1&0
\end{pmatrix}.
\]
\begin{problemsubparts}
\item Find a basis for $\Row(\mathbf A)$ and a basis for $\Col(\mathbf A)$.
\item Consider the restriction
\[
\mathbf A\big|_{\Row(\mathbf A)}\colon \Row(\mathbf A)\to\Col(\mathbf A).
\]
Find its matrix representation with respect to the bases obtained in part \textup{(a)}, and verify that this matrix is invertible.
\end{problemsubparts}
\end{problem}

% 2019 Fall POSTECH Applied Linear Algebra Midterm, Problem 6
\begin{problem}[2019 Fall \POSTECH{} Applied Linear Algebra]
Let $V\subseteq\R^3$ be a two-dimensional subspace.
\begin{problemsubparts}
\item Show that the reflection of any vector $\mathbf x\in\R^3$ across $V$ is
\[
(2\mathbf P-\mathbf I_3)\mathbf x,
\]
where $\mathbf P$ is the orthogonal projection onto $V$. \textit{Hint.} Consider an orthonormal basis of $V$ and extend it to an orthonormal basis of $\R^3$.
\item Let
\[
\mathbf e_1=
\begin{pmatrix}
1\\0\\0
\end{pmatrix},
\qquad
\mathbf e_2=
\begin{pmatrix}
0\\1\\0
\end{pmatrix},
\qquad
\mathbf e_3=
\begin{pmatrix}
0\\0\\1
\end{pmatrix}
\]
be the standard basis of $\R^3$. Let $V$ be spanned by
\[
\begin{pmatrix}
1\\0\\1
\end{pmatrix}
\qquad\text{and}\qquad
\begin{pmatrix}
0\\1\\2
\end{pmatrix}.
\]
Find the matrix representation of reflection across $V$ with respect to the standard basis.
\end{problemsubparts}
\end{problem}

% 2019 Fall POSTECH Applied Linear Algebra Midterm, Problem 7
\begin{problem}[2019 Fall \POSTECH{} Applied Linear Algebra]
Let
\[
\mathbf A=
\begin{pmatrix}
0&0&0&0&2\\
0&0&2&0&-1\\
1&2&-1&0&2\\
0&1&0&2&0\\
0&0&0&2&1
\end{pmatrix}.
\]
\begin{problemsubparts}
\item Find the matrices $\mathbf P$, $\mathbf U$, and $\mathbf L$ in a decomposition
\[
\mathbf P\mathbf A=\mathbf L\mathbf U.
\]
\item Solve $\mathbf A\mathbf x=\mathbf b$ for $\mathbf x$, where
\[
\mathbf b=
\begin{pmatrix}
-2\\3\\5\\4\\1
\end{pmatrix}.
\]
\end{problemsubparts}
\end{problem}

% 2019 Fall POSTECH Applied Linear Algebra Midterm, Problem 8
\begin{problem}[2019 Fall \POSTECH{} Applied Linear Algebra]
Find a nonzero matrix $\mathbf A$ whose null space is spanned by
\[
\begin{pmatrix}
1\\1\\0
\end{pmatrix}
\qquad\text{and}\qquad
\begin{pmatrix}
0\\1\\1
\end{pmatrix},
\]
and whose left null space is spanned by
\[
\begin{pmatrix}
2\\3\\-1
\end{pmatrix}.
\]
\end{problem}

% 2019 Fall POSTECH Applied Linear Algebra Midterm, Problem 9
\begin{problem}[2019 Fall \POSTECH{} Applied Linear Algebra]
Fit a quadratic function
\[
f(t)=c_0+c_1t+c_2t^2
\]
to the four data points
\[
(-1,8),\qquad(0,8),\qquad(1,4),\qquad(2,16).
\]
Explain why your answer is the best fit.
\end{problem}

% 2019 Fall POSTECH Applied Linear Algebra Midterm, Problem 10
\begin{problem}[2019 Fall \POSTECH{} Applied Linear Algebra]
Decompose
\[
\mathbf A=
\begin{pmatrix}
1&-1&0&0\\
-1&2&-1&0\\
0&-1&2&-1\\
0&0&-1&2
\end{pmatrix}
\]
as
\[
\mathbf A=\mathbf Q\mathbf R,
\]
where $\mathbf R$ is an invertible upper triangular matrix and $\mathbf Q$ is an orthogonal matrix.
\end{problem}

% 2019 Fall POSTECH Applied Linear Algebra Midterm, Problem 11
\begin{problem}[2019 Fall \POSTECH{} Applied Linear Algebra]
Let $P_4$ be the vector space of real polynomials of degree at most $4$ in the variable $t$. The map
\[
T\colon P_4\to P_4,
\qquad
T(p(t))=p(t+1),
\]
is linear.
\begin{problemsubparts}
\item Find the matrix $\mathbf A$ representing $T$ with respect to the ordered basis
\[
(1,t+1,t^2+1,t^3+1,t^4+1).
\]
\item Find $\mathbf A^{-1}$.
\end{problemsubparts}
\end{problem}

% 2019 Fall POSTECH Applied Linear Algebra Final, Problem 1
\begin{problem}[2019 Fall \POSTECH{} Applied Linear Algebra]
Fit a quadratic function
\[
f(t)=c_0+c_1t+c_2t^2
\]
to the four data points
\[
(-1,8),\qquad(0,8),\qquad(1,4),\qquad(2,16).
\]
Explain why your answer is the best fit.
\end{problem}

% 2019 Fall POSTECH Applied Linear Algebra Final, Problem 2
\begin{problem}[2019 Fall \POSTECH{} Applied Linear Algebra]
Decompose
\[
\mathbf A=
\begin{pmatrix}
1&-1&0&0\\
-1&2&-1&0\\
0&-1&2&-1\\
0&0&-1&2
\end{pmatrix}
\]
as
\[
\mathbf A=\mathbf Q\mathbf R,
\]
where $\mathbf R$ is an invertible upper triangular matrix and $\mathbf Q$ is an orthogonal matrix.
\end{problem}

% 2019 Fall POSTECH Applied Linear Algebra Final, Problem 3
\begin{problem}[2019 Fall \POSTECH{} Applied Linear Algebra]
For a positive integer $n$, let $S_n$ be the set of linear combinations of
\[
\sin x,\dots,\sin nx,\cos x,\dots,\cos nx,
\]
that is,
\[
S_n=\left\{\sum_{k=1}^n\bigl(a_k\sin kx+b_k\cos kx\bigr)\ \middle|\ a_k,b_k\in\R\right\}.
\]
Then the map
\[
D\colon S_n\to S_n,
\qquad
D(f)=\frac{\dd f}{\dd x},
\]
is linear.
\begin{problemsubparts}
\item Find the matrix representation of $D$ with respect to the ordered basis
\[
(\sin x,\cos x,\sin 2x,\cos 2x,\dots,\sin nx,\cos nx).
\]
\item Find the determinant and trace of $D$.
\end{problemsubparts}
\end{problem}

% 2019 Fall POSTECH Applied Linear Algebra Final, Problem 4
\begin{problem}[2019 Fall \POSTECH{} Applied Linear Algebra]
Let $T$ be a linear operator on $\R^5$ represented by
\[
\mathbf A=
\begin{pmatrix}
1&0&-2&0&2\\
1&1&-3&0&3\\
1&2&-4&0&3\\
1&0&-1&1&1\\
1&2&-3&0&2
\end{pmatrix}.
\]
\begin{problemsubparts}
\item Show that the characteristic polynomial has the form
\[
p_{\mathbf A}(\lambda)=-(\lambda-1)^3(\lambda-a)^2,
\]
and find $a$.
\item Find bases for the eigenspaces $E_1(\mathbf A)$ and $E_a(\mathbf A)$. Here, $E_1(\mathbf A)$ (respectively, $E_a(\mathbf A)$) denotes the subspace of $\R^5$ consisting of all eigenvectors of $\mathbf A$ corresponding to the eigenvalue $1$ (respectively, $a$), together with the zero vector.
\item Find the Jordan form of $\mathbf A$.
\end{problemsubparts}
\end{problem}

% 2019 Fall POSTECH Applied Linear Algebra Final, Problem 5
\begin{problem}[2019 Fall \POSTECH{} Applied Linear Algebra]
Factorize
\[
\mathbf A=
\begin{pmatrix}
1&i\\
-i&1
\end{pmatrix}
\]
as
\[
\mathbf A=\mathbf U\boldsymbol\Lambda\mathbf U^{-1},
\]
where $\mathbf U$ is unitary and $\boldsymbol\Lambda$ is diagonal.
\end{problem}

% 2019 Fall POSTECH Applied Linear Algebra Final, Problem 6
\begin{problem}[2019 Fall \POSTECH{} Applied Linear Algebra]
Suppose three cities $A$, $B$, and $C$ are centers for trucks. Every month, half of the trucks in $A$ and half of those in $B$ go to $C$, while the other halves stay where they are. Every month, half of the trucks in $C$ go to $A$ and the other half go to $B$.
\begin{problemsubparts}
\item Find the Markov transition matrix acting on the initial-state vector
\[
\begin{pmatrix}
a_0\\b_0\\c_0
\end{pmatrix},
\]
where $a_0,b_0,c_0$ denote the initial numbers of trucks in cities $A,B,C$, respectively.
\item If there are always $450$ trucks, find the long-run distribution of trucks.
\end{problemsubparts}
\end{problem}

% 2019 Fall POSTECH Applied Linear Algebra Final, Problem 7
\begin{problem}[2019 Fall \POSTECH{} Applied Linear Algebra]
Let
\[
\mathbf A=
\begin{pmatrix}
a&1&1\\
1&a&1\\
1&1&a
\end{pmatrix}.
\]
For what range of $a$ is $\mathbf A$ positive definite?
\end{problem}

% 2019 Fall POSTECH Applied Linear Algebra Final, Problem 8
\begin{problem}[2019 Fall \POSTECH{} Applied Linear Algebra]
Find the singular value decomposition of
\[
\mathbf A=
\begin{pmatrix}
3&2&2\\
2&3&-2
\end{pmatrix}.
\]
\end{problem}

% 2019 Fall POSTECH Applied Linear Algebra Final, Problem 9
\begin{problem}[2019 Fall \POSTECH{} Applied Linear Algebra]
Let $\mathbf A$ be an $n\times n$ matrix.
\begin{problemsubparts}
\item Show that $\mathbf A=\mathbf 0$ if
\[
\det(\mathbf A-\lambda\mathbf I_n)=(-\lambda)^n
\]
and $\mathbf A$ is Hermitian.
\item Show that $\mathbf A^n=\mathbf 0$ if
\[
\det(\mathbf A-\lambda\mathbf I_n)=(-\lambda)^n.
\]
\textit{Hint.} Consider the Jordan form.
\end{problemsubparts}
\end{problem}

% 2019 Fall POSTECH Applied Linear Algebra Final, Problem 10
\begin{problem}[2019 Fall \POSTECH{} Applied Linear Algebra]
Let $\lambda_0$ be an eigenvalue of a square matrix $\mathbf A$, and let $m$ be the algebraic multiplicity of $\lambda_0$ in the characteristic polynomial of $\mathbf A$. Show that
\[
\dim E_{\lambda_0}(\mathbf A)\le m.
\]
\end{problem}

% 2020 Spring POSTECH Applied Linear Algebra Quiz, Problem 1
\begin{problem}[2020 Spring \POSTECH{} Applied Linear Algebra]
Determine whether the following statement is true or false: the product of triangular matrices is a triangular matrix.
\end{problem}

% 2020 Spring POSTECH Applied Linear Algebra Quiz, Problem 2
\begin{problem}[2020 Spring \POSTECH{} Applied Linear Algebra]
Let $\mathbf D$ be an $n\times n$ diagonal matrix. Determine whether the following statement is true or false:
\[
\mathbf A\mathbf D=\mathbf D\mathbf A
\]
for every $n\times n$ matrix $\mathbf A$.
\end{problem}

% 2020 Spring POSTECH Applied Linear Algebra Quiz, Problem 3
\begin{problem}[2020 Spring \POSTECH{} Applied Linear Algebra]
Determine whether the following statement is true or false: every nonsingular symmetric square matrix $\mathbf A$ admits an $LDL^{\mathsf T}$ factorization, where $\mathbf L$ is lower triangular and $\mathbf D$ is diagonal.
\end{problem}

% 2020 Spring POSTECH Applied Linear Algebra Quiz, Problem 4
\begin{problem}[2020 Spring \POSTECH{} Applied Linear Algebra]
Let $\mathbf A$ be a square matrix. Determine whether the following statement is true or false:
\[
\dim\Col(\mathbf A^2)=\dim\Col(\mathbf A).
\]
\end{problem}

% 2020 Spring POSTECH Applied Linear Algebra Quiz, Problem 5
\begin{problem}[2020 Spring \POSTECH{} Applied Linear Algebra]
Let $\mathbf A$ be a rectangular matrix. Determine whether each of the following statements is true or false.
\begin{problemchoices}
\item $\Null(\mathbf A^{\mathsf T}\mathbf A)=\Null(\mathbf A)$.
\item $\Null(\mathbf A\mathbf A^{\mathsf T})=\Null(\mathbf A)$.
\item If $\mathbf A$ is square, then $\Null(\mathbf A^2)=\Null(\mathbf A)$.
\end{problemchoices}
\end{problem}

% 2020 Spring POSTECH Applied Linear Algebra Quiz, Problem 6
\begin{problem}[2020 Spring \POSTECH{} Applied Linear Algebra]
Determine whether the following statement is true or false: if distinct nonzero vectors $\mathbf v_1,\mathbf v_2,\mathbf v_3,\mathbf v_4$ span a subspace $V$ of $\R^5$, then there exist mutually orthonormal vectors $\mathbf w_1,\mathbf w_2,\mathbf w_3,\mathbf w_4$ that span $V$.
\end{problem}

% 2020 Spring POSTECH Applied Linear Algebra Quiz, Problem 7
\begin{problem}[2020 Spring \POSTECH{} Applied Linear Algebra]
Let
\[
\begin{aligned}
\mathbf v_1&=\begin{pmatrix}-1\\1\\0\\0\end{pmatrix},&
\mathbf v_2&=\begin{pmatrix}1\\0\\1\\0\end{pmatrix},&
\mathbf v_3&=\begin{pmatrix}-1\\0\\0\\1\end{pmatrix},\\[0.4em]
\mathbf v_4&=\begin{pmatrix}0\\-1\\1\\0\end{pmatrix},&
\mathbf v_5&=\begin{pmatrix}0\\1\\0\\1\end{pmatrix},&
\mathbf v_6&=\begin{pmatrix}0\\0\\1\\1\end{pmatrix}.
\end{aligned}
\]
Select all true statements.
\begin{problemchoices}
\item $\mathbf v_1,\mathbf v_2,\mathbf v_3,\mathbf v_4$ span $\R^4$.
\item $\mathbf v_1,\mathbf v_3,\mathbf v_4,\mathbf v_5,\mathbf v_6$ span $\R^4$.
\item $\mathbf v_4,\mathbf v_5,\mathbf v_6$ are linearly independent over $\R$.
\item $\mathbf v_3,\mathbf v_4,\mathbf v_5$ are linearly independent over $\R$.
\item $\mathbf v_3,\mathbf v_4,\mathbf v_5$ span $\R^3$.
\item None of the above.
\end{problemchoices}
\end{problem}

% 2020 Spring POSTECH Applied Linear Algebra Quiz, Problem 8
\begin{problem}[2020 Spring \POSTECH{} Applied Linear Algebra]
Let
\[
\mathbf A=
\begin{pmatrix}
3&1&-1\\
1&3&1\\
1&1&0
\end{pmatrix}
=
\begin{pmatrix}
1&0&0\\
\frac13&1&0\\
\frac13&\frac14&1
\end{pmatrix}
\begin{pmatrix}
3&1&-1\\
0&\frac83&\frac43\\
0&0&0
\end{pmatrix}.
\]
Select all true statements.
\begin{problemchoices}
\item The vectors
\[
\begin{pmatrix}3\\1\\-1\end{pmatrix}
\quad\text{and}\quad
\begin{pmatrix}0\\8\\4\end{pmatrix}
\]
span $\Row(\mathbf A)$.
\item The vectors
\[
\begin{pmatrix}3\\0\\0\end{pmatrix}
\quad\text{and}\quad
\begin{pmatrix}-1\\\frac43\\0\end{pmatrix}
\]
span $\Col(\mathbf A)$.
\item The left null space $\Null(\mathbf A^{\mathsf T})$ is spanned by
\[
\begin{pmatrix}-1\\-1\\4\end{pmatrix}
\quad\text{and}\quad
\begin{pmatrix}0\\0\\0\end{pmatrix}.
\]
\item $\mathbf A$ has a left inverse.
\item $\mathbf A$ has a right inverse.
\item None of the above.
\end{problemchoices}
\end{problem}

% 2020 Spring POSTECH Applied Linear Algebra Quiz, Problem 9
\begin{problem}[2020 Spring \POSTECH{} Applied Linear Algebra]
For $a,b,c,d\in\R$, let
\[
\mathbf A=
\begin{pmatrix}
a&b&0&0\\
c&d&0&0\\
0&0&a&b\\
0&0&c&d
\end{pmatrix}.
\]
Select all true statements.
\begin{problemchoices}
\item $\mathbf A$ is invertible if $abcd\ne0$.
\item $\mathbf A$ has a right inverse if $cd=0$.
\item $\mathbf A$ has both a right inverse and a left inverse if $ad-bc\ne0$.
\item $\mathbf A$ has a left inverse if $ab\ne0$.
\item None of the above.
\end{problemchoices}
\end{problem}

% 2020 Spring POSTECH Applied Linear Algebra Quiz, Problem 10
\begin{problem}[2020 Spring \POSTECH{} Applied Linear Algebra]
Let $P_n$ be the vector space of real polynomials of degree at most $n$. Let $E_n$ be the set of even polynomials in $P_n$, that is, polynomials $p$ satisfying $p(x)=p(-x)$, and let $O_n$ be the set of odd polynomials in $P_n$, that is, polynomials $q$ satisfying $q(x)=-q(-x)$. Determine whether each of the following statements is true or false.
\begin{problemchoices}
\item $E_n$ is a subspace of $P_n$.
\item $O_n$ is a subspace of $P_n$.
\item The map
\[
P_{n+1}\to P_n,
\qquad
f(x)\mapsto \frac{\dd f}{\dd x},
\]
is linear.
\item The map
\[
P_n\to P_{2n+2},
\qquad
f(x)\mapsto \int_1^{x^2}f(u)\,\dd u,
\]
is linear.
\item $\dim P_n=n$.
\end{problemchoices}
\end{problem}

% 2020 Spring POSTECH Applied Linear Algebra Quiz, Problem 11
\begin{problem}[2020 Spring \POSTECH{} Applied Linear Algebra]
Let $V$ be a one-dimensional vector space. Determine whether the following statement is true or false: if $f\colon V\to\R$ satisfies
\[
f(a\mathbf v)=a f(\mathbf v)
\]
for every $a\in\R$ and $\mathbf v\in V$, then $f$ is linear.
\end{problem}

% 2020 Spring POSTECH Applied Linear Algebra Quiz, Problem 12
\begin{problem}[2020 Spring \POSTECH{} Applied Linear Algebra]
Let $V=M_3(\R)$, the vector space of all $3\times3$ real matrices. Determine whether each of the following statements is true or false.
\begin{problemchoices}
\item $\dim V=9$.
\item The subset of invertible $3\times3$ matrices is a subspace of $V$.
\item The subset of $3\times3$ matrices whose entries sum to zero is a subspace of $V$.
\item The set of all linear transformations $V\to V$ is a vector space.
\item The matrix representation of a linear transformation $V\to V$ is a $3\times3$ matrix.
\end{problemchoices}
\end{problem}

% 2020 Spring POSTECH Applied Linear Algebra Quiz, Problem 13
\begin{problem}[2020 Spring \POSTECH{} Applied Linear Algebra]
Let $V$ be a vector space and let $\mathbf x,\mathbf y,\mathbf z\in V$. Determine whether each of the following statements is true or false.
\begin{problemchoices}
\item If $\mathbf x,\mathbf y,\mathbf z$ are linearly independent, then $\mathbf x,\mathbf y$ are linearly independent.
\item If $\mathbf x,\mathbf y,\mathbf z$ are linearly dependent, then
\[
\Span(\mathbf x,\mathbf y,\mathbf z)=\Span(\mathbf x,\mathbf y).
\]
\item If $\mathbf x,\mathbf y,\mathbf z$ are linearly independent, then
\[
\dim\Span(\mathbf x,\mathbf y,\mathbf z)=3.
\]
\item If $\dim\Span(\mathbf x,\mathbf y,\mathbf z)=3$, then $\mathbf x,\mathbf y,\mathbf z$ are linearly independent.
\item $\{\mathbf x,\mathbf y,\mathbf z\}$ is a basis for $\Span(\mathbf x,\mathbf y,\mathbf z)$.
\end{problemchoices}
\end{problem}

% 2020 Spring POSTECH Applied Linear Algebra Quiz, Problem 14
\begin{problem}[2020 Spring \POSTECH{} Applied Linear Algebra]
Let $\mathbf A\in\R^{m\times n}$. Determine whether each of the following statements is true or false.
\begin{problemchoices}
\item $\mathbf A\mathbf x=\mathbf b$ has a solution if and only if $\mathbf b\in\Col(\mathbf A)$.
\item $\mathbf A\mathbf x=\mathbf b$ has a unique solution if and only if $\rk(\mathbf A)=n$.
\item If $\rk(\mathbf A)=m$, then $\mathbf A\mathbf x=\mathbf b$ has a solution for every $\mathbf b\in\R^m$.
\item If $m<n$, then $\mathbf A\mathbf x=\mathbf b$ has infinitely many solutions.
\item If $m>n$, then there exists $\mathbf b\in\R^m$ for which $\mathbf A\mathbf x=\mathbf b$ has no solution.
\end{problemchoices}
\end{problem}

% 2020 Spring POSTECH Applied Linear Algebra Quiz, Problem 15
\begin{problem}[2020 Spring \POSTECH{} Applied Linear Algebra]
Let $\mathbf U$ be an echelon form of a matrix $\mathbf A$. Determine whether each of the following statements is true or false.
\begin{problemchoices}
\item $\Row(\mathbf A)=\Row(\mathbf U)$.
\item $\Col(\mathbf A)=\Col(\mathbf U)$.
\item $\Null(\mathbf A)=\Null(\mathbf U)$.
\item $\Row(\mathbf A)=\Col(\mathbf A)$.
\item $\rk(\mathbf A)=\rk(\mathbf U)$.
\end{problemchoices}
\end{problem}

% 2020 Spring POSTECH Applied Linear Algebra Quiz, Problem 16
\begin{problem}[2020 Spring \POSTECH{} Applied Linear Algebra]
Let $U\subseteq\R^n$ be a subspace. Determine whether each of the following statements is true or false.
\begin{problemchoices}
\item Every $\mathbf x\in\R^n$ can be written uniquely as
\[
\mathbf x=\mathbf u+\mathbf v,
\qquad
\mathbf u\in U,
\quad
\mathbf v\in U^\perp.
\]
\item If $\mathbf u\in U$ and $\mathbf v\in U^\perp$, then $\mathbf u$ and $\mathbf v$ are linearly independent.
\item $U\cap U^\perp=\{\mathbf 0\}$.
\item $\dim U=\dim U^\perp$.
\item If $\mathbf A\in\R^{n\times n}$ satisfies $\Col(\mathbf A)=U$, then
\[
\mathbf A(\mathbf A^{\mathsf T}\mathbf A)^{-1}\mathbf A^{\mathsf T}
\]
is the projection matrix onto $U$.
\end{problemchoices}
\end{problem}

% 2020 Spring POSTECH Applied Linear Algebra Quiz, Problem 17
\begin{problem}[2020 Spring \POSTECH{} Applied Linear Algebra]
Let $\mathbf x,\mathbf y,\mathbf z$ be linearly independent vectors in $\R^5$. Applying the Gram--Schmidt process to $\mathbf x,\mathbf y,\mathbf z$ produces orthonormal vectors $\mathbf u_1,\mathbf u_2,\mathbf u_3\in\R^5$. Let $\mathbf A$ be the $5\times3$ matrix with columns $\mathbf x,\mathbf y,\mathbf z$, and let $\mathbf Q$ be the $5\times3$ matrix with columns $\mathbf u_1,\mathbf u_2,\mathbf u_3$. Select all true statements.
\begin{problemchoices}
\item $\mathbf Q^{\mathsf T}\mathbf Q=\mathbf I_3$.
\item $\mathbf Q\mathbf Q^{\mathsf T}=\mathbf I_5$.
\item $\mathbf Q$ has orthonormal rows.
\item $\Col(\mathbf A)=\Col(\mathbf Q)$.
\item
\[
\mathbf u_1\mathbf u_1^{\mathsf T}+\mathbf u_2\mathbf u_2^{\mathsf T}+\mathbf u_3\mathbf u_3^{\mathsf T}
\]
is the projection matrix onto $\Col(\mathbf Q)$.
\item None of the above.
\end{problemchoices}
\end{problem}

% 2020 Spring POSTECH Applied Linear Algebra Quiz, Problem 18
\begin{problem}[2020 Spring \POSTECH{} Applied Linear Algebra]
Suppose that $\mathbf A\mathbf x=\mathbf b$ is inconsistent, and let $\mathbf P$ be the orthogonal projection onto $\Col(\mathbf A)$. Determine whether each of the following statements is true or false.
\begin{problemchoices}
\item There is always a least-squares solution.
\item The least-squares solution is unique.
\item Every least-squares solution satisfies
\[
\mathbf A^{\mathsf T}\mathbf A\mathbf x=\mathbf A^{\mathsf T}\mathbf b.
\]
\item Every solution of
\[
\mathbf A\mathbf x=\mathbf P\mathbf b
\]
is a least-squares solution.
\item
\[
\mathbf A(\mathbf A^{\mathsf T}\mathbf A)^{-1}\mathbf A^{\mathsf T}\mathbf b
\]
is a least-squares solution.
\end{problemchoices}
\end{problem}

% 2020 Spring POSTECH Applied Linear Algebra Quiz, Problem 19
\begin{problem}[2020 Spring \POSTECH{} Applied Linear Algebra]
Let $P_2$ be the vector space of real polynomials of degree at most $2$. Select the matrix representation of the linear transformation
\[
T\colon P_2\to\R^2,
\qquad
p(x)\mapsto\bigl(p(2),p'(2)\bigr),
\]
with respect to the ordered bases $(1,x,x^2)$ of $P_2$ and $(\mathbf e_1,\mathbf e_2)$ of $\R^2$.
\begin{problemchoices}
\item $\displaystyle\begin{pmatrix}2&0\\0&2\end{pmatrix}$
\item $\displaystyle\begin{pmatrix}2&0\\0&2\\0&0\end{pmatrix}$
\item $\displaystyle\begin{pmatrix}2&0&0\\0&2&0\end{pmatrix}$
\item $\displaystyle\begin{pmatrix}1&2&4\\0&1&4\end{pmatrix}$
\item $\displaystyle\begin{pmatrix}1&0\\2&1\\4&4\end{pmatrix}$
\end{problemchoices}
\end{problem}

% 2020 Spring POSTECH Applied Linear Algebra Quiz, Problem 20
\begin{problem}[2020 Spring \POSTECH{} Applied Linear Algebra]
Let $V$ be a vector space. Determine whether the following statement is true or false: the matrix representation of the identity transformation
\[
\id_V\colon V\to V,
\qquad
\mathbf v\mapsto\mathbf v,
\]
with respect to any choice of bases of $V$ is an identity matrix.
\end{problem}

% 2020 Spring POSTECH Applied Linear Algebra Final, Problem 1
\begin{problem}[2020 Spring \POSTECH{} Applied Linear Algebra]
Let $\mathbf A$ and $\mathbf B$ be $n\times n$ matrices. Determine whether the following statement is true or false:
\[
\det(\mathbf A+\mathbf B)=\det(\mathbf A)+\det(\mathbf B).
\]
\end{problem}

% 2020 Spring POSTECH Applied Linear Algebra Final, Problem 2
\begin{problem}[2020 Spring \POSTECH{} Applied Linear Algebra]
Let $\mathbf A$ be an $n\times n$ matrix. Determine whether the following statement is true or false:
\[
\det(2\mathbf A)=2\det(\mathbf A).
\]
\end{problem}

% 2020 Spring POSTECH Applied Linear Algebra Final, Problem 3
\begin{problem}[2020 Spring \POSTECH{} Applied Linear Algebra]
Let $\mathbf A$ and $\mathbf B$ be $n\times n$ matrices. Determine whether the following statement is true or false:
\[
\det(\mathbf A\mathbf B)=\det(\mathbf A)\det(\mathbf B).
\]
\end{problem}

% 2020 Spring POSTECH Applied Linear Algebra Final, Problem 4
\begin{problem}[2020 Spring \POSTECH{} Applied Linear Algebra]
For
\[
\det(\mathbf A-\lambda\mathbf I_3)
=
\det
\begin{pmatrix}
a_{11}-\lambda&a_{12}&a_{13}\\
a_{21}&a_{22}-\lambda&a_{23}\\
a_{31}&a_{32}&a_{33}-\lambda
\end{pmatrix},
\]
what is the coefficient of $\lambda^2$?
\begin{problemchoices}
\item $2$
\item $a_{11}+a_{22}+a_{33}$
\item $-a_{11}a_{22}a_{33}$
\item $\det(\mathbf A)$
\item None of the above.
\end{problemchoices}
\end{problem}

% 2020 Spring POSTECH Applied Linear Algebra Final, Problem 5
\begin{problem}[2020 Spring \POSTECH{} Applied Linear Algebra]
Let $\mathbf A$ be an invertible matrix and let $\mathbf C$ be its cofactor matrix. Determine whether the following statement is true or false:
\[
\mathbf A^{-1}=\frac{1}{\det(\mathbf A)}\mathbf C^{\mathsf T}.
\]
\end{problem}

% 2020 Spring POSTECH Applied Linear Algebra Final, Problem 6
\begin{problem}[2020 Spring \POSTECH{} Applied Linear Algebra]
A parallelepiped in $\R^3$ has edge vectors
\[
\begin{pmatrix}1\\1\\0\end{pmatrix},
\qquad
\begin{pmatrix}-1\\2\\1\end{pmatrix},
\qquad
\begin{pmatrix}0\\1\\-3\end{pmatrix}.
\]
What is its volume?
\begin{problemchoices}
\item $\dfrac{10}{3}$
\item $5$
\item $7.5$
\item $10$
\item None of the above.
\end{problemchoices}
\end{problem}

% 2020 Spring POSTECH Applied Linear Algebra Final, Problem 7
\begin{problem}[2020 Spring \POSTECH{} Applied Linear Algebra]
Determine whether the following statement is true or false: the determinant of a permutation matrix is always $1$.
\end{problem}

% 2020 Spring POSTECH Applied Linear Algebra Final, Problem 8
\begin{problem}[2020 Spring \POSTECH{} Applied Linear Algebra]
Let $\mathbf A$ be an $n\times n$ matrix and $\mathbf b\in\R^n$. Determine whether the following statement is true or false: Cramer's rule gives a solution of $\mathbf A\mathbf x=\mathbf b$ for every such $\mathbf A$ and $\mathbf b$.
\end{problem}

% 2020 Spring POSTECH Applied Linear Algebra Final, Problem 9
\begin{problem}[2020 Spring \POSTECH{} Applied Linear Algebra]
Let $\mathbf A$ be an $n\times n$ matrix. Determine whether the following statement is true or false: if $\mathbf A$ has $n$ distinct eigenvalues, then $\mathbf A$ is diagonalizable.
\end{problem}

% 2020 Spring POSTECH Applied Linear Algebra Final, Problem 10
\begin{problem}[2020 Spring \POSTECH{} Applied Linear Algebra]
Let $\mathbf A$ be an $n\times n$ matrix. Determine whether the following statement is true or false: if $\mathbf A$ has fewer than $n$ distinct eigenvalues, then $\mathbf A$ is not diagonalizable.
\end{problem}

% 2020 Spring POSTECH Applied Linear Algebra Final, Problem 11
\begin{problem}[2020 Spring \POSTECH{} Applied Linear Algebra]
Let $\mathbf A$ be an $n\times n$ matrix. Determine whether the following statement is true or false: eigenvectors of $\mathbf A$ associated with distinct eigenvalues are mutually orthogonal.
\end{problem}

% 2020 Spring POSTECH Applied Linear Algebra Final, Problem 12
\begin{problem}[2020 Spring \POSTECH{} Applied Linear Algebra]
Let $\mathbf A$ be an $n\times n$ complex matrix. Determine whether the following statement is true or false: $\mathbf A$ is unitarily diagonalizable if and only if $\C^n$ has an orthonormal basis consisting of eigenvectors of $\mathbf A$.
\end{problem}

% 2020 Spring POSTECH Applied Linear Algebra Final, Problem 13
\begin{problem}[2020 Spring \POSTECH{} Applied Linear Algebra]
Let $\mathbf A$ be an $n\times n$ matrix. Determine whether the following statement is true or false: the solution of the initial value problem
\[
\mathbf u'(t)=\mathbf A\mathbf u(t),
\qquad
\mathbf u(0)=\mathbf u_0,
\]
is
\[
\mathbf u(t)=e^{t\mathbf A}\mathbf u_0,
\]
regardless of whether $\mathbf A$ is diagonalizable.
\end{problem}

% 2020 Spring POSTECH Applied Linear Algebra Final, Problem 14
\begin{problem}[2020 Spring \POSTECH{} Applied Linear Algebra]
Determine whether the following statement is true or false: every real symmetric matrix is orthogonally diagonalizable.
\end{problem}

% 2020 Spring POSTECH Applied Linear Algebra Final, Problem 15
\begin{problem}[2020 Spring \POSTECH{} Applied Linear Algebra]
Determine whether the following statement is true or false: every Hermitian matrix is unitarily diagonalizable.
\end{problem}

% 2020 Spring POSTECH Applied Linear Algebra Final, Problem 16
\begin{problem}[2020 Spring \POSTECH{} Applied Linear Algebra]
Determine whether the following statement is true or false: every orthogonally diagonalizable matrix is real symmetric.
\end{problem}

% 2020 Spring POSTECH Applied Linear Algebra Final, Problem 17
\begin{problem}[2020 Spring \POSTECH{} Applied Linear Algebra]
Let $\mathbf A$ and $\mathbf B$ be $n\times n$ real symmetric matrices. Determine whether the following statement is true or false: if $\mathbf A$ and $\mathbf B$ are congruent, then they have the same eigenvalues.
\end{problem}

% 2020 Spring POSTECH Applied Linear Algebra Final, Problem 18
\begin{problem}[2020 Spring \POSTECH{} Applied Linear Algebra]
Let $\mathbf A$ and $\mathbf B$ be $n\times n$ real symmetric matrices. Determine whether the following statement is true or false: if $\mathbf A$ and $\mathbf B$ are congruent, then they have the same number of positive eigenvalues.
\end{problem}

% 2020 Spring POSTECH Applied Linear Algebra Final, Problem 19
\begin{problem}[2020 Spring \POSTECH{} Applied Linear Algebra]
Determine whether the following statement is true or false: every real symmetric matrix is congruent to a diagonal matrix.
\end{problem}

% 2020 Spring POSTECH Applied Linear Algebra Final, Problem 20
\begin{problem}[2020 Spring \POSTECH{} Applied Linear Algebra]
Let $\mathbf A$ be a real symmetric matrix. Determine whether the following statement is true or false: if $\mathbf A$ is positive definite, then every leading principal submatrix of $\mathbf A$ has positive determinant.
\end{problem}

% 2020 Spring POSTECH Applied Linear Algebra Final, Problem 21
\begin{problem}[2020 Spring \POSTECH{} Applied Linear Algebra]
Let $\mathbf A\in\R^{m\times n}$. Determine whether the following statement is true or false: every real matrix has a singular value decomposition.
\end{problem}

% 2020 Spring POSTECH Applied Linear Algebra Final, Problem 22
\begin{problem}[2020 Spring \POSTECH{} Applied Linear Algebra]
Let $\mathbf A\in\R^{m\times n}$. Determine whether the following statement is true or false: the singular values of $\mathbf A$ are the nonnegative square roots of the eigenvalues of $\mathbf A^{\mathsf T}\mathbf A$.
\end{problem}

% 2020 Spring POSTECH Applied Linear Algebra Final, Problem 23
\begin{problem}[2020 Spring \POSTECH{} Applied Linear Algebra]
Let $\mathbf A\in\R^{m\times n}$. Determine whether the following statement is true or false: if $\mathbf A$ has a left inverse, then its pseudoinverse is a left inverse of $\mathbf A$.
\end{problem}

% 2020 Spring POSTECH Applied Linear Algebra Final, Problem 24
\begin{problem}[2020 Spring \POSTECH{} Applied Linear Algebra]
Determine whether the following statement is true or false: every matrix has a pseudoinverse.
\end{problem}

% 2020 Spring POSTECH Applied Linear Algebra Final, Problem 25
\begin{problem}[2020 Spring \POSTECH{} Applied Linear Algebra]
Let $\mathbf A\in\R^{m\times n}$. Determine whether the following statement is true or false: $\mathbf A^{\mathsf T}\mathbf A$ and $\mathbf A\mathbf A^{\mathsf T}$ have the same eigenvalues.
\end{problem}

% 2020 Spring POSTECH Applied Linear Algebra Final, Problem 26
\begin{problem}[2020 Spring \POSTECH{} Applied Linear Algebra]
Let $\mathbf A$ be a $2\times2$ upper triangular matrix with diagonal entries $1$ and $2$. Determine whether the following statement is true or false: $\mathbf A$ is diagonalizable.
\end{problem}

% 2020 Spring POSTECH Applied Linear Algebra Final, Problem 27
\begin{problem}[2020 Spring \POSTECH{} Applied Linear Algebra]
Determine whether the following matrix is diagonalizable:
\[
\mathbf A=
\begin{pmatrix}
1&-1\\
1&-1
\end{pmatrix}.
\]
\end{problem}

% 2020 Spring POSTECH Applied Linear Algebra Final, Problem 28
\begin{problem}[2020 Spring \POSTECH{} Applied Linear Algebra]
Determine whether the following statement is true or false:
\[
\begin{pmatrix}
1&0\\
0&2
\end{pmatrix}^{1000}
=
\begin{pmatrix}
1&0\\
0&2^{1000}
\end{pmatrix}.
\]
\end{problem}

% 2020 Spring POSTECH Applied Linear Algebra Final, Problem 29
\begin{problem}[2020 Spring \POSTECH{} Applied Linear Algebra]
Let
\[
\mathbf A=
\begin{pmatrix}
3&2\\
1&1
\end{pmatrix}
\begin{pmatrix}
1&0\\
0&4
\end{pmatrix}
\begin{pmatrix}
3&2\\
1&1
\end{pmatrix}^{-1}.
\]
Determine whether the following statement is true or false: the vectors
\[
\begin{pmatrix}3\\1\end{pmatrix}
\quad\text{and}\quad
\begin{pmatrix}2\\1\end{pmatrix}
\]
are eigenvectors of $\mathbf A$.
\end{problem}

% 2020 Spring POSTECH Applied Linear Algebra Final, Problem 30
\begin{problem}[2020 Spring \POSTECH{} Applied Linear Algebra]
Let $\mathbf A$ be a $2\times2$ matrix whose eigenvalues are $1$ and $2$. Determine whether the following statement is true or false: $\mathbf A$ is invertible.
\end{problem}

% 2020 Spring POSTECH Applied Linear Algebra Final, Problem 31
\begin{problem}[2020 Spring \POSTECH{} Applied Linear Algebra]
Let $\mathbf A$ be a $2\times2$ matrix whose only eigenvalue is $2$. Determine whether the following statement is true or false: $\mathbf A$ is not invertible.
\end{problem}

% 2020 Spring POSTECH Applied Linear Algebra Final, Problem 32
\begin{problem}[2020 Spring \POSTECH{} Applied Linear Algebra]
Let $\mathbf A$ be a $2\times2$ matrix. Determine whether the following statement is true or false: if $\mathbf A$ is not diagonalizable, then it has only one eigenvalue, counted with algebraic multiplicity two.
\end{problem}

% 2020 Spring POSTECH Applied Linear Algebra Final, Problem 33
\begin{problem}[2020 Spring \POSTECH{} Applied Linear Algebra]
Let $\mathbf A$ be a $2\times2$ matrix with repeated eigenvalue $3$. Determine whether the following statement is true or false: $\mathbf A$ is invertible.
\end{problem}

% 2020 Spring POSTECH Applied Linear Algebra Final, Problem 34
\begin{problem}[2020 Spring \POSTECH{} Applied Linear Algebra]
Let $\mathbf A$ and $\mathbf B$ be $2\times2$ matrices. Determine whether the following statement is true or false:
\[
e^{\mathbf A+\mathbf B}=e^{\mathbf A}e^{\mathbf B}.
\]
\end{problem}

% 2020 Spring POSTECH Applied Linear Algebra Final, Problem 35
\begin{problem}[2020 Spring \POSTECH{} Applied Linear Algebra]
Let $\mathbf A$ be an $n\times n$ complex matrix. Determine whether the following statement is true or false: either
\[
\det(\mathbf A)=\det(\mathbf A^{\mathsf H})
\]
or
\[
\det(\mathbf A)=-\det(\mathbf A^{\mathsf H}).
\]
\end{problem}

% 2020 Spring POSTECH Applied Linear Algebra Final, Problem 36
\begin{problem}[2020 Spring \POSTECH{} Applied Linear Algebra]
Let $\mathbf A$ and $\mathbf B$ be $n\times n$ orthogonal matrices. Determine whether the following statement is true or false: $\mathbf A\mathbf B$ is orthogonal.
\end{problem}

% 2020 Spring POSTECH Applied Linear Algebra Final, Problem 37
\begin{problem}[2020 Spring \POSTECH{} Applied Linear Algebra]
Let $\mathbf A$ be a $5\times5$ matrix with rank $4$. Determine whether the following statement is true or false:
\[
\det(\mathbf A)=0.
\]
\end{problem}

% 2020 Spring POSTECH Applied Linear Algebra Final, Problem 38
\begin{problem}[2020 Spring \POSTECH{} Applied Linear Algebra]
Let
\[
\mathbf A=
\begin{pmatrix}
4&1\\
2&3
\end{pmatrix}.
\]
Determine whether the following statement is true or false: $\mathbf A-\alpha\mathbf I_2$ is singular only when $\alpha=2$ or $\alpha=5$.
\end{problem}

% 2020 Spring POSTECH Applied Linear Algebra Final, Problem 39
\begin{problem}[2020 Spring \POSTECH{} Applied Linear Algebra]
Determine whether the following statement is true or false: every full-rank $2\times2$ matrix is diagonalizable.
\end{problem}

% 2020 Spring POSTECH Applied Linear Algebra Final, Problem 40
\begin{problem}[2020 Spring \POSTECH{} Applied Linear Algebra]
Determine whether the following statement is true or false: it is possible to choose $x\in\R$ such that $\det(\mathbf A)\ne0$ for
\[
\mathbf A=
\begin{pmatrix}
x&0&0\\
x&0&0\\
x&x&x
\end{pmatrix}.
\]
\end{problem}

% 2020 Spring POSTECH Applied Linear Algebra Final, Problem 41
\begin{problem}[2020 Spring \POSTECH{} Applied Linear Algebra]
Determine whether the following statement is true or false: the only possible eigenvalues of a real symmetric orthogonal matrix are $1$ and $-1$.
\end{problem}

% 2020 Spring POSTECH Applied Linear Algebra Final, Problem 42
\begin{problem}[2020 Spring \POSTECH{} Applied Linear Algebra]
Determine whether the following statement is true or false: similar matrices have the same eigenvalues, counting algebraic multiplicity.
\end{problem}

% 2020 Spring POSTECH Applied Linear Algebra Final, Problem 43
\begin{problem}[2020 Spring \POSTECH{} Applied Linear Algebra]
Determine whether the following statement is true or false: if two $n\times n$ matrices have the same eigenvalues, counting algebraic multiplicity, then they are similar.
\end{problem}

% 2020 Spring POSTECH Applied Linear Algebra Final, Problem 44
\begin{problem}[2020 Spring \POSTECH{} Applied Linear Algebra]
Let $\mathbf A$ be a symmetric matrix. Determine whether the following statement is true or false: $\mathbf A$ is positive definite if and only if there exists a matrix $\mathbf R$ such that
\[
\mathbf A=\mathbf R^{\mathsf T}\mathbf R.
\]
\end{problem}

% 2020 Spring POSTECH Applied Linear Algebra Final, Problem 45
\begin{problem}[2020 Spring \POSTECH{} Applied Linear Algebra]
Let $\mathbf A$ be a $4\times4$ real symmetric matrix with $\det(\mathbf A)=30$. Suppose that $\mathbf A$ is not positive definite, all eigenvalues of $\mathbf A$ are integers, and two eigenvalues are $2$ and $3$. What are the other two eigenvalues?
\begin{problemchoices}
\item $1,5$
\item $-1,-5$
\item $1,-5$
\item $2,3$
\item None of the above.
\end{problemchoices}
\end{problem}

% 2020 Spring POSTECH Applied Linear Algebra Final, Problem 46
\begin{problem}[2020 Spring \POSTECH{} Applied Linear Algebra]
Determine whether the following statement is true or false: every symmetric matrix with positive determinant is positive definite.
\end{problem}

% 2020 Spring POSTECH Applied Linear Algebra Final, Problem 47
\begin{problem}[2020 Spring \POSTECH{} Applied Linear Algebra]
Let $\mathbf A$ be a square matrix. Determine whether the following statement is true or false: if every cofactor of $\mathbf A$ is zero, then $\mathbf A$ is not invertible.
\end{problem}

% 2020 Spring POSTECH Applied Linear Algebra Final, Problem 48
\begin{problem}[2020 Spring \POSTECH{} Applied Linear Algebra]
Let $\mathbf A$ be a $3\times3$ matrix with exactly two linearly independent eigenvectors $\mathbf v_1,\mathbf v_2$, corresponding to eigenvalues $\lambda_1,\lambda_2$, respectively.
\begin{problemsubparts}
\item Determine whether the following statement is true or false: $\lambda_1$ and $\lambda_2$ must be distinct.
\item Select a Jordan form to which $\mathbf A$ is similar.
\begin{problemchoices}
\item $\displaystyle\begin{pmatrix}\lambda_1&0&0\\0&\lambda_2&0\\0&0&\lambda_2\end{pmatrix}$
\item $\displaystyle\begin{pmatrix}\lambda_1&1&0\\0&\lambda_2&0\\0&0&\lambda_2\end{pmatrix}$
\item $\displaystyle\begin{pmatrix}\lambda_1&0&0\\0&\lambda_2&1\\0&0&\lambda_2\end{pmatrix}$
\item $\displaystyle\begin{pmatrix}\lambda_1&1&0\\0&\lambda_2&1\\0&0&\lambda_2\end{pmatrix}$
\item None of the above.
\end{problemchoices}
\end{problemsubparts}
\end{problem}

% 2020 Spring POSTECH Applied Linear Algebra Final, Problem 49
\begin{problem}[2020 Spring \POSTECH{} Applied Linear Algebra]
Determine whether the following statement is true or false: every complex square matrix has at least one eigenvector.
\end{problem}

% 2020 Fall POSTECH Applied Linear Algebra Midterm Part I, Problem 1
\begin{problem}[2020 Fall \POSTECH{} Applied Linear Algebra]
Let
\[
\mathbf A=
\begin{pmatrix}
1&2&4\\
2&3&5\\
4&5&6
\end{pmatrix},
\qquad
\mathbf b=
\begin{pmatrix}
1\\1\\1
\end{pmatrix}.
\]
\begin{problemsubparts}
\item Solve $\mathbf A\mathbf x=\mathbf b$.
\item Find an $LDU$ decomposition
\[
\mathbf A=\mathbf L\mathbf D\mathbf U,
\]
where $\mathbf L$ is lower triangular with diagonal entries $1$, $\mathbf D$ is diagonal, and $\mathbf U$ is upper triangular with diagonal entries $1$.
\item Find $\mathbf A^{-1}$.
\end{problemsubparts}
\end{problem}

% 2020 Fall POSTECH Applied Linear Algebra Midterm Part I, Problem 2
\begin{problem}[2020 Fall \POSTECH{} Applied Linear Algebra]
For
\[
\mathbf A=
\begin{pmatrix}
1&0&0\\
2&1&0\\
3&4&1
\end{pmatrix}
\begin{pmatrix}
1&2&3&4\\
0&1&2&3\\
0&0&1&2
\end{pmatrix},
\]
find bases for the row space $\Row(\mathbf A)$, equivalently $\Col(\mathbf A^{\mathsf T})$, and the null space $\Null(\mathbf A)$.
\end{problem}

% 2020 Fall POSTECH Applied Linear Algebra Midterm Part I, Problem 3
\begin{problem}[2020 Fall \POSTECH{} Applied Linear Algebra]
Let $V$ be the subspace of $\R^4$ spanned by the orthonormal vectors
\[
\begin{pmatrix}\frac12\\-\frac12\\\frac12\\-\frac12\end{pmatrix},
\qquad
\begin{pmatrix}0\\\frac1{\sqrt2}\\\frac1{\sqrt2}\\0\end{pmatrix},
\qquad
\begin{pmatrix}\frac1{\sqrt2}\\0\\0\\\frac1{\sqrt2}\end{pmatrix}.
\]
Find the orthogonal projection of
\[
\mathbf b=
\begin{pmatrix}
1\\1\\1\\-1
\end{pmatrix}
\]
onto $V$.
\end{problem}

% 2020 Fall POSTECH Applied Linear Algebra Midterm Part II, Problem 1
\begin{problem}[2020 Fall \POSTECH{} Applied Linear Algebra]
Let $V$ be a subspace of $\R^n$ and let $W=V^\perp$. Show that
\[
V+W=\R^n.
\]
\end{problem}

% 2020 Fall POSTECH Applied Linear Algebra Midterm Part II, Problem 2
\begin{problem}[2020 Fall \POSTECH{} Applied Linear Algebra]
Let $\mathfrak{sl}_2$ be the set of $2\times2$ real matrices whose diagonal entries sum to zero.
\begin{problemsubparts}
\item Show that $\mathfrak{sl}_2$ is a vector space.
\item Find a basis for $\mathfrak{sl}_2$, and prove that your chosen matrices form a basis.
\item Find $\dim\mathfrak{sl}_2$.
\end{problemsubparts}
\end{problem}

% 2020 Fall POSTECH Applied Linear Algebra Midterm Part II, Problem 3
\begin{problem}[2020 Fall \POSTECH{} Applied Linear Algebra]
Let
\[
\mathbf H=\mathbf I_n-2\mathbf u\mathbf u^{\mathsf T},
\]
where $\mathbf u\in\R^n$ is a unit vector.
\begin{problemsubparts}
\item Show that
\[
\mathbf H\mathbf u=-\mathbf u
\]
and
\[
\mathbf H\mathbf w=\mathbf w
\]
for every $\mathbf w\perp\mathbf u$.
\item Show that $\|\mathbf H\mathbf x\|=\|\mathbf x\|$ for every $\mathbf x$ of compatible dimension.
\item For $n=3$, let
\[
\mathbf x=
\begin{pmatrix}0\\1\\1\end{pmatrix},
\qquad
\mathbf b=
\begin{pmatrix}\sqrt2\\0\\0\end{pmatrix}.
\]
Find a unit vector $\mathbf u$ such that $\mathbf x$ solves
\[
\mathbf H\mathbf x=\mathbf b.
\]
\end{problemsubparts}
\end{problem}

% 2020 Fall POSTECH Applied Linear Algebra Final Part I, Problem 1
\begin{problem}[2020 Fall \POSTECH{} Applied Linear Algebra]
Let
\[
\mathbf A=
\begin{pmatrix}
0&-1&-2\\
-2&-1&2\\
3&-1&-5
\end{pmatrix}.
\]
\begin{problemsubparts}
\item Find the eigenvalues of $\mathbf A$.
\item Find the eigenvectors of $\mathbf A$.
\item Solve
\[
\mathbf y'(t)=\mathbf A\mathbf y(t),
\qquad
\mathbf y(0)=
\begin{pmatrix}
1\\2\\1
\end{pmatrix},
\]
and express each component of $\mathbf y(t)$ explicitly.
\end{problemsubparts}
\end{problem}

% 2020 Fall POSTECH Applied Linear Algebra Final Part I, Problem 2
\begin{problem}[2020 Fall \POSTECH{} Applied Linear Algebra]
Let
\[
V=\{\mathbf v_1,\dots,\mathbf v_n\}\subseteq\R^n
\]
be a set of nonzero mutually orthogonal vectors.
\begin{problemsubparts}
\item Show that $V$ is a basis for $\R^n$.
\item Let $\mathbf A$ be the $n\times n$ matrix whose columns are $\mathbf v_1,\dots,\mathbf v_n$. Find $\det(\mathbf A)$ in terms of $\mathbf v_1,\dots,\mathbf v_n$.
\item Let $\mathbf B$ be any $n\times n$ matrix. Show that
\[
|\det(\mathbf B)|=\operatorname{vol}(P),
\]
where $P$ is the parallelepiped spanned by the columns of $\mathbf B$.
\end{problemsubparts}
\end{problem}

% 2020 Fall POSTECH Applied Linear Algebra Final Part I, Problem 3
\begin{problem}[2020 Fall \POSTECH{} Applied Linear Algebra]
\begin{problemsubparts}
\item Find a $6\times6$ matrix $\mathbf A$ satisfying
\[
\det(\mathbf A-\lambda\mathbf I_6)
=(\lambda-2)(\lambda-3)^2(\lambda-4)^3,
\]
and
\[
\dim E_2(\mathbf A)=1,
\qquad
\dim E_3(\mathbf A)=1,
\qquad
\dim E_4(\mathbf A)=2.
\]
\item Compute $e^{t\mathbf A}$ for the matrix $\mathbf A$ found in part \textup{(a)}.
\end{problemsubparts}
\end{problem}

% 2020 Fall POSTECH Applied Linear Algebra Final Part I, Problem 4
\begin{problem}[2020 Fall \POSTECH{} Applied Linear Algebra]
Consider the Markov matrix
\[
\mathbf A=
\begin{pmatrix}
0.5&0.5&0.25\\
0&0.5&0.25\\
0.5&0&0.5
\end{pmatrix}.
\]
\begin{problemsubparts}
\item Determine whether the sequence
\[
\mathbf A^k
\begin{pmatrix}
0.25\\0.25\\0.5
\end{pmatrix}
\]
converges as $k\to\infty$.
\item Give an example of a Markov matrix $\mathbf M$ and a column vector $\mathbf u_0$ for which the sequence $\mathbf M^k\mathbf u_0$ does not converge as $k\to\infty$.
\end{problemsubparts}
\end{problem}

% 2020 Fall POSTECH Applied Linear Algebra Final Part I, Problem 5
\begin{problem}[2020 Fall \POSTECH{} Applied Linear Algebra]
For
\[
f(x,y,z)=2x^2+5y^2+8z^2+4xy+6yz,
\]
let $\mathbf A$ be a symmetric matrix satisfying
\[
f(x,y,z)=\mathbf x^{\mathsf T}\mathbf A\mathbf x,
\qquad
\mathbf x=
\begin{pmatrix}
x\\y\\z
\end{pmatrix}.
\]
\begin{problemsubparts}
\item Find $\mathbf A$.
\item Show that $\mathbf A$ is positive definite using determinants.
\item Show that $\mathbf A$ is positive definite using pivots.
\item Show that $\mathbf A$ is positive definite using eigenvalues. \textit{Hint.} $10$ is an eigenvalue.
\end{problemsubparts}
\end{problem}

% 2020 Fall POSTECH Applied Linear Algebra Final Part II, Problem 1
\begin{problem}[2020 Fall \POSTECH{} Applied Linear Algebra]
Suppose a vector $\mathbf p_0\in\R^3$ rotates at a constant angular rate of $1$ rad/s about the unit axis
\[
\mathbf w=
\begin{pmatrix}
w_1\\w_2\\w_3
\end{pmatrix},
\qquad
\|\mathbf w\|=1.
\]
Let $\mathbf p(t)$ denote its position at time $t$. Its velocity is
\[
\mathbf p'(t)=\mathbf w\times\mathbf p(t).
\]
Since $\mathbf w\times\mathbf p(t)=\boldsymbol\Omega\mathbf p(t)$, where
\[
\boldsymbol\Omega=
\begin{pmatrix}
0&-w_3&w_2\\
w_3&0&-w_1\\
-w_2&w_1&0
\end{pmatrix},
\]
we have
\[
\mathbf p'(t)=\boldsymbol\Omega\mathbf p(t),
\qquad
\mathbf p(0)=\mathbf p_0.
\]
\begin{problemsubparts}
\item Find the eigenvalues of $\boldsymbol\Omega$.
\item Show that
\[
e^{t\boldsymbol\Omega}
=\mathbf I_3+(\sin t)\boldsymbol\Omega+(1-\cos t)\boldsymbol\Omega^2.
\]
\textit{Hint.} Use the power-series expansions of $\sin t$ and $\cos t$.
\item Show that the position vector preserves its length:
\[
\|\mathbf p(t)\|=\|\mathbf p_0\|
\]
for all $t$.
\item Suppose
\[
\mathbf w=
\begin{pmatrix}0\\1\\0\end{pmatrix},
\qquad
\mathbf p_0=
\begin{pmatrix}0\\1\\1\end{pmatrix}.
\]
Find $\mathbf p(\pi/2)$.
\end{problemsubparts}
\end{problem}

% 2020 Fall POSTECH Applied Linear Algebra Final Part II, Problem 2
\begin{problem}[2020 Fall \POSTECH{} Applied Linear Algebra]
Let
\[
\mathbf A=\mathbf I_n-\mathbf u\mathbf v^{\mathsf T},
\qquad
\mathbf u,\mathbf v\in\R^n.
\]
\begin{problemsubparts}
\item Show that if $\mathbf A$ is invertible, then its inverse has the form
\[
\mathbf A^{-1}=\mathbf I_n+\alpha\mathbf u\mathbf v^{\mathsf T}
\]
for some scalar $\alpha$, and find $\alpha$.
\item Determine for which $\mathbf u$ and $\mathbf v$ the matrix $\mathbf A$ is not invertible.
\item If $\mathbf A$ is not invertible, find $\Null(\mathbf A)$.
\end{problemsubparts}
\end{problem}

% 2020 Fall POSTECH Applied Linear Algebra Final Part II, Problem 3
\begin{problem}[2020 Fall \POSTECH{} Applied Linear Algebra]
Let
\[
\mathbf A=
\begin{pmatrix}
2&1&0\\
1&2&1\\
0&1&2
\end{pmatrix}.
\]
\begin{problemsubparts}
\item Find a diagonal matrix $\boldsymbol\Lambda$ and an orthogonal matrix $\mathbf Q$ such that
\[
\mathbf A=\mathbf Q\boldsymbol\Lambda\mathbf Q^{\mathsf T}.
\]
\item Compute $\mathbf A^{2020}$, describing each entry exactly.
\item Let $T\colon \R^3\to\R^3$ be the linear transformation represented by $\mathbf A$. Find an orthonormal basis of $\R^3$ with respect to which the matrix representation of $T$ is diagonal.
\end{problemsubparts}
\end{problem}

% 2021 Spring POSTECH Applied Linear Algebra Midterm, Source 1, Problem 1
\begin{problem}[2021 Spring \POSTECH{} Applied Linear Algebra]
Find an $LDU$ decomposition
\[
\mathbf A=
\begin{pmatrix}
1&2&3\\
0&1&0\\
1&4&4
\end{pmatrix}
=\mathbf L\mathbf D\mathbf U,
\]
where $\mathbf L$ is lower triangular with diagonal entries $1$, $\mathbf D$ is diagonal, and $\mathbf U$ is upper triangular with diagonal entries $1$.
\end{problem}

% 2021 Spring POSTECH Applied Linear Algebra Midterm, Source 1, Problem 2
\begin{problem}[2021 Spring \POSTECH{} Applied Linear Algebra]
Find an $LDU$ decomposition
\[
\mathbf A=
\begin{pmatrix}
1&0&1&2\\
0&1&1&1\\
1&1&3&1\\
2&1&1&4
\end{pmatrix}
=\mathbf L\mathbf D\mathbf U,
\]
where $\mathbf L$ is lower triangular with diagonal entries $1$, $\mathbf D$ is diagonal, and $\mathbf U$ is upper triangular with diagonal entries $1$.
\end{problem}

% 2021 Spring POSTECH Applied Linear Algebra Midterm, Source 1, Problem 3
\begin{problem}[2021 Spring \POSTECH{} Applied Linear Algebra]
Let
\[
\mathbf A=
\begin{pmatrix}
1&2&0\\
2&3&4\\
3&4&5
\end{pmatrix}.
\]
Find a $QR$ decomposition $\mathbf A=\mathbf Q\mathbf R$, where $\mathbf Q$ is orthogonal and $\mathbf R$ is invertible and upper triangular.
\end{problem}

% 2021 Spring POSTECH Applied Linear Algebra Midterm, Source 1, Problem 4
\begin{problem}[2021 Spring \POSTECH{} Applied Linear Algebra]
Let $\mathbf A\in\R^{m\times n}$ and assume that $\mathbf A^{\mathsf T}\mathbf A$ is invertible. Define
\[
\mathbf P=\mathbf A(\mathbf A^{\mathsf T}\mathbf A)^{-1}\mathbf A^{\mathsf T}.
\]
For $\mathbf b\in\R^m$, show that $\mathbf b-\mathbf P\mathbf b$ is orthogonal to every vector in $\Col(\mathbf A)$.
\end{problem}

% 2021 Spring POSTECH Applied Linear Algebra Midterm, Source 1, Problem 5
\begin{problem}[2021 Spring \POSTECH{} Applied Linear Algebra]
Determine whether each of the following statements is true or false for every matrix $\mathbf A\in\R^{m\times n}$. You need not justify your answers.
\begin{problemchoices}
\item $\mathbf A$ and $-\mathbf A$ have the same reduced row echelon form.
\item If $\mathbf A$ is invertible, then $\mathbf A^{\mathsf T}$ is invertible.
\item If $\mathbf A\mathbf x=\mathbf b$ has a solution for every $\mathbf b\in\R^m$, then $n\ge m$.
\item If $\mathbf A\mathbf x=\mathbf b$ has a unique solution for some $\mathbf b\in\R^m$, then $n\le m$.
\item If the columns of a square matrix $\mathbf A$ are orthonormal, then its rows are orthonormal.
\item If $\mathbf A\mathbf x=\mathbf b$ has infinitely many solutions for some $\mathbf b\in\R^m$, then $\mathbf A\mathbf x=\mathbf 0$ has a nonzero solution.
\item $\rk(\mathbf A^{\mathsf T}\mathbf A)=\rk(\mathbf A)$.
\end{problemchoices}
\end{problem}

% 2021 Spring POSTECH Applied Linear Algebra Midterm, Source 1, Problem 6
\begin{problem}[2021 Spring \POSTECH{} Applied Linear Algebra]
Find the inverse of
\[
\mathbf B=
\begin{pmatrix}
1&0&0&0\\
4&1&0&0\\
3&4&1&0\\
2&3&4&1
\end{pmatrix}.
\]
\end{problem}

% 2021 Spring POSTECH Applied Linear Algebra Midterm, Source 1, Problem 7
\begin{problem}[2021 Spring \POSTECH{} Applied Linear Algebra]
Let $P_4$ be the vector space of real polynomials of degree at most $4$, and define the linear map
\[
T\colon P_4\to P_4,
\qquad
T(p(t))=t^2\frac{\dd^2p(t)}{\dd t^2}.
\]
\begin{problemsubparts}
\item Find the matrix $\mathbf A$ representing $T$ with respect to the ordered basis
\[
(1,t,t^2,t^3,t^4).
\]
\item Find the matrix $\mathbf B$ representing $T$ with respect to the ordered basis
\[
(1,t+1,t^2+1,t^3+1,t^4+1).
\]
\end{problemsubparts}
\end{problem}

% 2021 Spring POSTECH Applied Linear Algebra Midterm, Source 1, Problem 8
\begin{problem}[2021 Spring \POSTECH{} Applied Linear Algebra]
Let
\[
\mathbf A=
\begin{pmatrix}
1&4&1&1&1\\
1&4&2&1&1\\
0&0&1&0&0
\end{pmatrix}.
\]
\begin{problemsubparts}
\item Find a basis for $\Col(\mathbf A)$.
\item Find a basis for $\Col(\mathbf A)^\perp$.
\item Find an orthonormal basis for $\Null(\mathbf A)$.
\end{problemsubparts}
\end{problem}

% 2021 Spring POSTECH Applied Linear Algebra Midterm, Source 1, Problem 9
\begin{problem}[2021 Spring \POSTECH{} Applied Linear Algebra]
State precisely the definition of each of the following concepts.
\begin{problemchoices}
\item Linear independence.
\item Orthogonal matrix.
\item Vector space.
\item Projection matrix.
\item Linear transformation.
\end{problemchoices}
\end{problem}

% 2021 Spring POSTECH Applied Linear Algebra Midterm, Source 2, Part I, Problem 1
\begin{problem}[2021 Spring \POSTECH{} Applied Linear Algebra]
Let
\[
\mathbf A=
\begin{pmatrix}
1&2&1\\
2&-1&0\\
2&4&3
\end{pmatrix}.
\]
\begin{problemsubparts}
\item Find $\mathbf A^{-1}$. Write your answer as a single explicit $3\times3$ matrix.
\item Find an $LDU$ decomposition of $\mathbf A$.
\end{problemsubparts}
\end{problem}

% 2021 Spring POSTECH Applied Linear Algebra Midterm, Source 2, Part I, Problem 2
\begin{problem}[2021 Spring \POSTECH{} Applied Linear Algebra]
Let
\[
V=\Span\!\left(
\begin{pmatrix}1\\2\\1\\2\\-2\end{pmatrix},
\begin{pmatrix}2\\2\\1\\1\\0\end{pmatrix},
\begin{pmatrix}3\\0\\0\\1\\1\end{pmatrix}
\right)\subseteq\R^5.
\]
\begin{problemsubparts}
\item Find $V^\perp$.
\item Find all $\mathbf x\in\R^5$ satisfying $\mathbf A\mathbf x=\mathbf b$, where
\[
\mathbf A=
\begin{pmatrix}
1&2&1&2&-2\\
2&2&1&1&0\\
3&0&0&1&1
\end{pmatrix},
\qquad
\mathbf b=
\begin{pmatrix}
1\\2\\3
\end{pmatrix}.
\]
\end{problemsubparts}
\end{problem}

% 2021 Spring POSTECH Applied Linear Algebra Midterm, Source 2, Part I, Problem 3
\begin{problem}[2021 Spring \POSTECH{} Applied Linear Algebra]
Let $T\colon \R^3\to\R^3$ be reflection across the plane
\[
z=2y.
\]
Find the matrix representing $T$ with respect to the standard basis. Write your answer as a single explicit $3\times3$ matrix.
\end{problem}

% 2021 Spring POSTECH Applied Linear Algebra Midterm, Source 2, Part II, Problem 1
\begin{problem}[2021 Spring \POSTECH{} Applied Linear Algebra]
Let
\[
\mathbf A=
\begin{pmatrix}
2&1&0\\
2&2&1\\
0&-1&-1
\end{pmatrix},
\qquad
\mathbf b=
\begin{pmatrix}
1\\2\\3
\end{pmatrix}.
\]
Find the minimum value of
\[
\|\mathbf A\mathbf x-\mathbf b\|^2.
\]
\end{problem}

% 2021 Spring POSTECH Applied Linear Algebra Midterm, Source 2, Part II, Problem 2
\begin{problem}[2021 Spring \POSTECH{} Applied Linear Algebra]
Let $V$ be a vector space. Suppose that
\[
\{\mathbf v_1,\dots,\mathbf v_n\}
\qquad\text{and}\qquad
\{\mathbf w_1,\dots,\mathbf w_m\}
\]
are two bases of $V$. Show that $m=n$.
\end{problem}

% 2021 Spring POSTECH Applied Linear Algebra Midterm, Source 2, Part II, Problem 3
\begin{problem}[2021 Spring \POSTECH{} Applied Linear Algebra]
Let $\mathbf A$ and $\mathbf B$ be $n\times n$ matrices. Show that if either $\mathbf A$ or $\mathbf B$ is not invertible, then $\mathbf A\mathbf B$ is not invertible.
\end{problem}

% 2021 Spring POSTECH Applied Linear Algebra Midterm, Source 2, Part II, Problem 4
\begin{problem}[2021 Spring \POSTECH{} Applied Linear Algebra]
Let $\mathbf Q$ be an $n\times n$ matrix. Show that if
\[
\|\mathbf Q\mathbf x\|=\|\mathbf x\|
\]
for every $\mathbf x\in\R^n$, then the rows of $\mathbf Q$ are orthonormal.
\end{problem}

% 2021 Spring POSTECH Applied Linear Algebra Midterm, Source 3, Problem 1
\begin{problem}[2021 Spring \POSTECH{} Applied Linear Algebra]
Let
\[
\mathbf A=
\begin{pmatrix}
2&4&6&4\\
2&5&7&6\\
2&3&5&2
\end{pmatrix}.
\]
\begin{problemsubparts}
\item Find an echelon form of $\mathbf A$.
\item Find $\rk(\mathbf A)$.
\item Find a basis for $\Col(\mathbf A)$.
\item Find a basis for $\Col(\mathbf A^{\mathsf T})$.
\item Find an $LDU$ factorization of $\mathbf A$.
\item Solve $\mathbf A\mathbf x=\mathbf b$, where
\[
\mathbf b=
\begin{pmatrix}
16\\20\\12
\end{pmatrix}.
\]
\end{problemsubparts}
\end{problem}

% 2021 Spring POSTECH Applied Linear Algebra Midterm, Source 3, Problem 2
\begin{problem}[2021 Spring \POSTECH{} Applied Linear Algebra]
Let
\[
\mathbf A=
\begin{pmatrix}
1&2&3\\
2&4&6\\
2&5&7\\
3&9&13
\end{pmatrix}.
\]
Show that the columns of $\mathbf A$ are linearly independent.
\end{problem}

% 2021 Spring POSTECH Applied Linear Algebra Midterm, Source 3, Problem 3
\begin{problem}[2021 Spring \POSTECH{} Applied Linear Algebra]
Let
\[
\mathbf A=
\begin{pmatrix}
2&-1&0\\
-1&2&-1\\
0&-1&2
\end{pmatrix}.
\]
Using the Gauss--Jordan method, find $\mathbf A^{-1}$.
\end{problem}

% 2021 Spring POSTECH Applied Linear Algebra Midterm, Source 3, Problem 4
\begin{problem}[2021 Spring \POSTECH{} Applied Linear Algebra]
Let $V,W\subseteq\R^n$ be subspaces.
\begin{problemsubparts}
\item Prove that $V+W$ is a subspace of $\R^n$.
\item Prove that $V\cap W$ is a subspace of $\R^n$.
\item Prove that
\[
\dim(V+W)+\dim(V\cap W)=\dim V+\dim W.
\]
\end{problemsubparts}
\end{problem}



% 2022 Spring KAIST Introduction to Linear Algebra Midterm, Problem 1
\begin{problem}[2022 Spring \KAIST{} Introduction to Linear Algebra]
Describe all possible $2\times4$ reduced row echelon forms whose rows are linearly independent.
\end{problem}

% 2022 Spring KAIST Introduction to Linear Algebra Midterm, Problem 2
\begin{problem}[2022 Spring \KAIST{} Introduction to Linear Algebra]
\begin{problemsubparts}
\item Let
\[
\mathbf A(\theta)=
\begin{pmatrix}
\cos\theta & \csc\theta & 0 & -\sin\theta\\
\cos^2\theta & \cot\theta & 2.78 & \sin^2\theta\\
0 & 3.14 & 0 & 0\\
\sin\theta & \tan\theta & 0 & \cos\theta
\end{pmatrix}.
\]
Compute
\[
\det\bigl(\mathbf A(1.76\pi)\bigr)-\det\bigl(\mathbf A(2\pi/3)\bigr).
\]
\item Suppose $\mathbf A\in\R^{3\times3}$ and $\det(\mathbf A)=2$. Compute
\[
\det(\mathbf A^{-1})+\det(3\mathbf A^{-1})+\det\bigl((5\mathbf A)^{-1}\bigr)+\det(\mathbf A^{\mathsf T}).
\]
\end{problemsubparts}
\end{problem}

% 2022 Spring KAIST Introduction to Linear Algebra Midterm, Problem 3
\begin{problem}[2022 Spring \KAIST{} Introduction to Linear Algebra]
There are three main colleges of science and engineering at \KAIST{}: the College of Natural Sciences (NS), the College of Life Science and Bioengineering (LSB), and the College of Engineering (E). Suppose students pursuing transdisciplinary study choose one of the following three course sets:
\begin{problemchoices}
\item Type I: $30$ credits from NS, $21$ credits from LSB, and $27$ credits from E;
\item Type II: $33$ credits from NS, $21$ credits from LSB, and $30$ credits from E;
\item Type III: $30$ credits from NS, $15$ credits from LSB, and $39$ credits from E.
\end{problemchoices}
Among graduates in 2022 who chose one of these three sets, the total numbers of credits taken were $4950$ from NS, $2880$ from LSB, and $5430$ from E. How many students chose Type I?
\end{problem}

% 2022 Spring KAIST Introduction to Linear Algebra Midterm, Problem 4
\begin{problem}[2022 Spring \KAIST{} Introduction to Linear Algebra]
Let
\[
\mathbf A=
\begin{pmatrix}
0&0&0&1\\
1&0&0&0\\
0&1&0&1\\
0&0&2&1
\end{pmatrix}.
\]
Find $\mathbf A^{-1}$.
\end{problem}

% 2022 Spring KAIST Introduction to Linear Algebra Midterm, Problem 5
\begin{problem}[2022 Spring \KAIST{} Introduction to Linear Algebra]
Let $\mathbf A\in\R^{m\times n}$ have linearly independent columns. Show that if $\mathbf v_1,\dots,\mathbf v_p\in\R^n$ are linearly independent, then
\[
\mathbf A\mathbf v_1,\dots,\mathbf A\mathbf v_p
\]
are linearly independent.
\end{problem}

% 2022 Spring KAIST Introduction to Linear Algebra Midterm, Problem 6
\begin{problem}[2022 Spring \KAIST{} Introduction to Linear Algebra]
Let $T\colon\R^n\to\R^n$ be a linear operator satisfying
\[
\|T(\mathbf x)\|=\|\mathbf x\|
\qquad\text{for every }\mathbf x\in\R^n.
\]
Prove that
\[
T(\mathbf x)\cdot T(\mathbf y)=\mathbf x\cdot\mathbf y
\qquad\text{for all }\mathbf x,\mathbf y\in\R^n.
\]
\end{problem}

% 2022 Spring KAIST Introduction to Linear Algebra Midterm, Problem 7
\begin{problem}[2022 Spring \KAIST{} Introduction to Linear Algebra]
Show that
\[
\mathbf A=
\begin{pmatrix}
\frac{1}{\sqrt2}&0&\frac{1}{\sqrt2}\\
0&-1&0\\
-\frac{1}{\sqrt2}&0&\frac{1}{\sqrt2}
\end{pmatrix}
\]
represents a rotation $R_\theta$ about a line $\ell$ through the origin, followed by a reflection $H_\ell$ about the plane through the origin perpendicular to $\ell$. Find the standard matrices of $R_\theta$ and $H_\ell$.
\end{problem}

% 2022 Spring KAIST Introduction to Linear Algebra Midterm, Problem 8
\begin{problem}[2022 Spring \KAIST{} Introduction to Linear Algebra]
For $n\times n$ matrices $\mathbf A$ and $\mathbf B$, one has
\[
\tr(\mathbf A\mathbf B)=\tr(\mathbf B\mathbf A).
\]
Prove or disprove the following statement: for $n\times n$ matrices $\mathbf A,\mathbf B,\mathbf C$,
\[
\tr(\mathbf A\mathbf B\mathbf C)=\tr(\mathbf B\mathbf A\mathbf C).
\]
\end{problem}

% 2022 Spring KAIST Introduction to Linear Algebra Midterm, Problem 9
\begin{problem}[2022 Spring \KAIST{} Introduction to Linear Algebra]
Suppose $p(x)$ and $q(x)$ are polynomials with real coefficients, and let $r(x)=p(x)q(x)$. Show that for every $n\times n$ matrix $\mathbf A$,
\[
r(\mathbf A)=p(\mathbf A)q(\mathbf A).
\]
\end{problem}

% 2022 Spring KAIST Introduction to Linear Algebra Midterm, Problem 10
\begin{problem}[2022 Spring \KAIST{} Introduction to Linear Algebra]
\begin{problemsubparts}
\item Find all $2\times2$ orthogonal matrices that do not admit an $LU$ decomposition.
\item For each such orthogonal matrix, find a $PLU$ decomposition.
\end{problemsubparts}
\end{problem}

% 2022 Spring KAIST Introduction to Linear Algebra Midterm, Problem 11
\begin{problem}[2022 Spring \KAIST{} Introduction to Linear Algebra]
Let $\mathbf A=(a_{ij})\in\R^{n\times n}$ be triangular, with pairwise distinct diagonal entries:
\[
a_{ii}\ne a_{jj}
\qquad\text{whenever }i\ne j.
\]
Find all eigenvalues of $\mathbf A$ and determine the algebraic multiplicity of each eigenvalue.
\end{problem}

% 2022 Spring KAIST Introduction to Linear Algebra Midterm, Problem 12
\begin{problem}[2022 Spring \KAIST{} Introduction to Linear Algebra]
Prove or disprove the following statement. For a consistent linear system
\[
\mathbf A\mathbf x=\mathbf b
\]
with $n$ unknowns and $m$ equations, the solution set has the form
\[
\mathbf x_0+\Null(\mathbf A)
=
\{\mathbf x_0+\mathbf v\mid \mathbf v\in\Null(\mathbf A)\}
\]
for some solution $\mathbf x_0$.
\end{problem}

% 2022 Spring KAIST Introduction to Linear Algebra Midterm, Problem 13
\begin{problem}[2022 Spring \KAIST{} Introduction to Linear Algebra]
Let $\mathbf A\in\R^{n\times n}$ be nilpotent; that is, $\mathbf A^k=\mathbf0$ for some positive integer $k$. Show that the matrix transformation
\[
T_{\mathbf A}\colon\R^n\to\R^n,
\qquad
T_{\mathbf A}(\mathbf x)=\mathbf A\mathbf x,
\]
is not injective.
\end{problem}

% 2022 Spring KAIST Introduction to Linear Algebra Midterm, Problem 14
\begin{problem}[2022 Spring \KAIST{} Introduction to Linear Algebra]
Let $\mathbf A=(a_{ij})\in\R^{3\times3}$ be skew-symmetric, so that
\[
\mathbf A^{\mathsf T}=-\mathbf A.
\]
\begin{problemsubparts}
\item Find the characteristic polynomial of $\mathbf A$ and all real eigenvalues of $\mathbf A$.
\item Find the full eigenspace corresponding to each real eigenvalue.
\end{problemsubparts}
\end{problem}

% 2022 Spring KAIST Introduction to Linear Algebra Midterm, Problem 15
\begin{problem}[2022 Spring \KAIST{} Introduction to Linear Algebra]
Let
\[
\mathbf A=\frac16
\begin{pmatrix}
1&1&2\\
1&1&2\\
2&2&4
\end{pmatrix}.
\]
Draw $\Col(\mathbf A)$ and $\Null(\mathbf A)$ in $\R^3$. It is known that $\mathbf A$ is the standard matrix of the orthogonal projection onto $\Col(\mathbf A)$. For $n$ randomly generated points, transform them as specified in the MATLAB code below and scatter-plot the original and transformed points. Fill in blanks \textup{(1)}--\textup{(10)}.

{\small\ttfamily
\begin{tabbing}
\hspace{2em}\=\hspace{2em}\=\kill
\%\%\% plot the column space and the nullspace of A \%\%\%\\
A = 1/6 * [1, 1, 2; 1, 1, 2; 2, 2, 4];\\
x = \_\_\_(1)\_\_\_; \% generate a vector that contains\\
\>\% 150 equally spaced points, from -1 to 1.\\
y = x;\\
{}[X, Y] = \_\_\_(2)\_\_\_(x, y); \% construct a grid over the domain\\[0.15em]
a = [\_\_\_(3)\_\_\_]; \% choose a column vector ``a'' in the column space of A\\
plot3(x, \_\_\_(4)\_\_\_, \_\_\_(5)\_\_\_) \% draw the column space of A\\
hold on;\\[0.15em]
Z = \_\_\_(6)\_\_\_; \% compute the z-components of the nullspace (plane) of A\\
\_\_\_(7)\_\_\_; \% draw the surface of the nullspace of A\\
hold on; \% keep using the same figure window for future plots\\[0.15em]
\%\%\% scatter plot n randomly generated points and their transformed points \%\%\%\\
n = 5;\\
b = -1 + 2*rand(3,n); \% generate n random points in [-1,1]\textasciicircum 3\\
p = zeros(3,n);\\
for i=1:n\\
\>\% if norm(b(:,i)) is greater than or equal to 1/2,\\
\>\% project b(:,i) onto the column space of A;\\
\>\% otherwise stretch its entries by 1, 0.5, and 2, respectively\\
\>\_\_\_(8)\_\_\_\\
\>\>p(:,i) = \_\_\_(9)\_\_\_;\\
\>else\\
\>\>p(:,i) = \_\_\_(10)\_\_\_;\\
\>end\\
end\\[0.15em]
scatter3(b(1,:),b(2,:),b(3,:),'*');\\
scatter3(p(1,:),p(2,:),p(3,:));\\
axis equal;
\end{tabbing}
}
\end{problem}

% 2022 Spring KAIST Introduction to Linear Algebra Final, Problem 1
\begin{problem}[2022 Spring \KAIST{} Introduction to Linear Algebra]
Let $\mathbf A\in\R^{n\times n}$. Determine whether each of the following statements is true or false, and justify your answer.
\begin{problemsubparts}
\item Elementary row operations do not change $\Row(\mathbf A)$.
\item Elementary row operations do not change $\Col(\mathbf A)$.
\item Elementary row operations do not change $\dim\Null(\mathbf A^{\mathsf T})$.
\item The nonzero rows of any row echelon form $\mathbf R$ of $\mathbf A$ form a basis for $\Row(\mathbf A)$, even when $\mathbf R$ is not in reduced row echelon form.
\item The nonzero columns of any row echelon form $\mathbf R$ of $\mathbf A$ form a basis for $\Col(\mathbf A)$, even when $\mathbf R$ is not in reduced row echelon form.
\end{problemsubparts}
\end{problem}

% 2022 Spring KAIST Introduction to Linear Algebra Final, Problem 2
\begin{problem}[2022 Spring \KAIST{} Introduction to Linear Algebra]
Given $\mathbf A\in\R^{m\times n}$, let $\mathcal B$ be a basis for $\Row(\mathbf A)$ and let $\mathcal B'$ be a basis for $\Null(\mathbf A)$. Prove that
\[
\mathcal B\cup\mathcal B'
\]
is a basis for $\R^n$.
\end{problem}

% 2022 Spring KAIST Introduction to Linear Algebra Final, Problem 3
\begin{problem}[2022 Spring \KAIST{} Introduction to Linear Algebra]
Let
\[
\mathbf A=
\begin{pmatrix}
-1&0\\
0&1\\
0&1
\end{pmatrix},
\qquad
\mathbf b=
\begin{pmatrix}
0\\0\\1
\end{pmatrix}.
\]
\begin{problemsubparts}
\item Write the normal equations associated with $\mathbf A\mathbf x=\mathbf b$.
\item Find all least-squares solutions of $\mathbf A\mathbf x=\mathbf b$. If the least-squares solution is unique, state that it is unique.
\end{problemsubparts}
\end{problem}

% 2022 Spring KAIST Introduction to Linear Algebra Final, Problem 4
\begin{problem}[2022 Spring \KAIST{} Introduction to Linear Algebra]
Prove or disprove the following statement: if $\mathbf P$ is the standard matrix of an orthogonal projection of $\R^n$ onto a subspace of $\R^n$, then
\[
\tr(\mathbf P)=\rk(\mathbf P).
\]
\end{problem}

% 2022 Spring KAIST Introduction to Linear Algebra Final, Problem 5
\begin{problem}[2022 Spring \KAIST{} Introduction to Linear Algebra]
Let
\[
\mathbf A=
\begin{pmatrix}
1&1&1\\
1&2&2^2\\
1&3&3^2\\
1&4&4^2\\
1&5&5^2
\end{pmatrix}.
\]
Does $\mathbf A$ have full column rank? If not, find a basis for $\Col(\mathbf A)$ consisting of columns of $\mathbf A$.
\end{problem}

% 2022 Spring KAIST Introduction to Linear Algebra Final, Problem 6
\begin{problem}[2022 Spring \KAIST{} Introduction to Linear Algebra]
Let
\[
\mathbf A=
\begin{pmatrix}
2&0&0&0\\
0&-2&0&2\\
0&0&-3&0\\
0&2&0&1
\end{pmatrix}.
\]
\begin{problemsubparts}
\item Find all eigenvalues of $\mathbf A$ together with their algebraic multiplicities.
\item Determine whether $\mathbf A$ is orthogonally diagonalizable. If it is, find an orthonormal basis of $\R^4$ consisting of eigenvectors of $\mathbf A$; otherwise, explain why not.
\item Find the solution set of
\[
\mathbf A\mathbf x=
\begin{pmatrix}
200\\150\\300\\-75
\end{pmatrix}.
\]
\end{problemsubparts}
\end{problem}

% 2022 Spring KAIST Introduction to Linear Algebra Final, Problem 7
\begin{problem}[2022 Spring \KAIST{} Introduction to Linear Algebra]
Show that a square matrix is nilpotent if and only if $0$ is its only eigenvalue.
\end{problem}

% 2022 Spring KAIST Introduction to Linear Algebra Final, Problem 8
\begin{problem}[2022 Spring \KAIST{} Introduction to Linear Algebra]
Find a column-row factorization (CR decomposition) of each matrix.
\begin{problemsubparts}
\item
\[
\mathbf A=
\begin{pmatrix}
1&2&3&4\\
2&4&6&8\\
3&6&9&12\\
4&8&12&16
\end{pmatrix}.
\]
\item
\[
\mathbf B=
\begin{pmatrix}
1&2&3\\
0&2&4\\
0&0&-1
\end{pmatrix}.
\]
\end{problemsubparts}
\end{problem}

% 2022 Spring KAIST Introduction to Linear Algebra Final, Problem 9
\begin{problem}[2022 Spring \KAIST{} Introduction to Linear Algebra]
Let $W$ be a plane through the origin in $\R^3$.
\begin{problemsubparts}
\item Show that the standard matrix of the orthogonal projection of $\R^3$ onto $W$ is symmetric.
\item Find the standard matrix $\mathbf P$ of the orthogonal projection onto the plane
\[
x-y+z=0.
\]
\end{problemsubparts}
\end{problem}

% 2022 Spring KAIST Introduction to Linear Algebra Final, Problem 10
\begin{problem}[2022 Spring \KAIST{} Introduction to Linear Algebra]
Let
\[
\mathbf A=
\begin{pmatrix}
1&0&1\\
1&0&-1\\
1&1&0\\
-1&1&1
\end{pmatrix}.
\]
\begin{problemsubparts}
\item Find a $QR$ decomposition of $\mathbf A$.
\item Find the least-squares solution $\widehat{\mathbf x}$ of
\[
\mathbf A\mathbf x=
\begin{pmatrix}
1\\0\\0\\1
\end{pmatrix}.
\]
\end{problemsubparts}
\end{problem}

% 2022 Spring KAIST Introduction to Linear Algebra Final, Problem 11
\begin{problem}[2022 Spring \KAIST{} Introduction to Linear Algebra]
Prove the following statements.
\begin{problemsubparts}
\item Let $S\subseteq\R^n$ be a set of nonzero, pairwise orthogonal vectors. Show that $S$ is linearly independent.
\item If $\mathbf u_1,\mathbf u_2,\dots,\mathbf u_m$ are eigenvectors of a matrix $\mathbf A$ corresponding to distinct eigenvalues $\lambda_1,\lambda_2,\dots,\lambda_m$, show that
\[
\{\mathbf u_1,\mathbf u_2,\dots,\mathbf u_m\}
\]
is linearly independent.
\end{problemsubparts}
\end{problem}

% 2022 Spring KAIST Introduction to Linear Algebra Final, Problem 12
\begin{problem}[2022 Spring \KAIST{} Introduction to Linear Algebra]
\begin{problemsubparts}
\item Let $W$ be a subspace of $\R^n$, and define
\[
U=\{\mathbf x\in\R^n\mid \mathbf x\perp\mathbf w\text{ for every }\mathbf w\in W\}.
\]
Show that $U$ is a subspace of $\R^n$.
\item Let $W\subseteq\R^5$ be spanned by
\[
\mathbf u_1=
\begin{pmatrix}2\\2\\-1\\0\\1\end{pmatrix},
\qquad
\mathbf u_2=
\begin{pmatrix}1\\1\\-2\\0\\-1\end{pmatrix},
\qquad
\mathbf u_3=
\begin{pmatrix}0\\0\\1\\1\\1\end{pmatrix}.
\]
Find a basis for the subspace $U$ defined in part \textup{(a)}.
\end{problemsubparts}
\end{problem}

% 2022 Spring KAIST Introduction to Linear Algebra Final, Problem 13
\begin{problem}[2022 Spring \KAIST{} Introduction to Linear Algebra]
Suppose an invertible square matrix $\mathbf A$ has a singular value decomposition
\[
\mathbf A=\mathbf U\boldsymbol\Sigma\mathbf V^{\mathsf T}.
\]
Give singular value decompositions of $\mathbf A^{\mathsf T}$ and $\mathbf A^{-1}$ in terms of $\mathbf U$, $\boldsymbol\Sigma$, and $\mathbf V$.
\end{problem}

% 2022 Spring KAIST Introduction to Linear Algebra Final, Problem 14
\begin{problem}[2022 Spring \KAIST{} Introduction to Linear Algebra]
Find a singular value decomposition of
\[
\mathbf A=
\begin{pmatrix}
-1&1\\
-1&-1
\end{pmatrix}.
\]
\end{problem}

% 2022 Spring KAIST Introduction to Linear Algebra Final, Problem 15
\begin{problem}[2022 Spring \KAIST{} Introduction to Linear Algebra]
Let $\mathbf A$ be a singular symmetric $n\times n$ matrix, and let $\mathbf B$ be a symmetric perturbation of $\mathbf A$. An eigenvector corresponding to the largest eigenvalue is called a dominating eigenvector. The MATLAB code below computes the angle between dominating eigenvectors of $\mathbf A$ and $\mathbf B$, and the angle between $\Col(\mathbf A)$ and the dominating eigenvector of $\mathbf B$. Fill in blanks \textup{(1)}--\textup{(9)}.

{\small\ttfamily
\begin{tabbing}
\hspace{2em}\=\hspace{2em}\=\kill
\% Input: a (singular) symmetric n by n matrix A.\\
n = size(A,1);\\
E = \_\_\_(1)\_\_\_; \% Generate an n by n matrix with random numbers\\
\>\% from a standard normal distribution.\\
epsilon = 0.001;\\
B = A + epsilon*1/2*(E+E');\\[0.15em]
\% Compute vA and vB that are eigenvectors\\
\% corresponding to the largest eigenvalues of A and B, respectively.\\
vA = final\_MAS109(A, 1); vB = final\_MAS109(B, 1);\\[0.15em]
\% Compute cos(theta), where theta is an angle between vA and vB.\\
cos\_theta = \_\_\_(2)\_\_\_;\\[0.15em]
\% Compute CA whose columns are orthonormal eigenvectors of A\\
\% that span the column space of A.\\
CA = \_\_\_(3)\_\_\_;\\[0.15em]
\% Express vB as a sum of two orthogonal terms vB1 and vB2.\\
vB1 = \_\_\_(4)\_\_\_; \% projection of vB onto the column space of A\\
vB2 = \_\_\_(5)\_\_\_; \% projection of vB onto the null space of A\\[0.15em]
\% Compute cos(phi), where phi is an angle between vB and the column space of A.\\
cos\_phi = \_\_\_(6)\_\_\_;\\[0.35em]
\% Define final\_MAS109, which outputs CA, whose columns are orthonormal\\
\% eigenvectors corresponding to the k largest eigenvalues of A.\\
function \_\_\_(7)\_\_\_\\
\>[V, D] = eig(A); \% for symmetric A, eigenvectors are orthonormal\\
\>lambdas = \_\_\_(8)\_\_\_; \% eigenvalues in vector form\\
\>[sorted\_lambdas index] = sort(lambdas, 'descend');\\
\>\% Example: if lambdas = [10 -5 4], then index = [1 3 2].\\
\>CA = \_\_\_(9)\_\_\_;\\
end
\end{tabbing}
}
\end{problem}

% 2022 Fall KAIST Introduction to Linear Algebra Midterm, Problem 1
\begin{problem}[2022 Fall \KAIST{} Introduction to Linear Algebra]
\begin{problemsubparts}
\item Let $P_1$ be the plane through the points
\[
(1,3,-3),\qquad(8,1,6),\qquad(0,5,-6).
\]
Find an equation for $P_1$.
\item Let $P_2$ be the plane with normal vector $(11,-3,5)$ passing through $(0,12,11)$. Show that $P_1\cap P_2$ is a line, and give a vector equation for this line.
\end{problemsubparts}
\end{problem}

% 2022 Fall KAIST Introduction to Linear Algebra Midterm, Problem 2
\begin{problem}[2022 Fall \KAIST{} Introduction to Linear Algebra]
\begin{problemsubparts}
\item Find an $LU$ decomposition of
\[
\mathbf A=
\begin{pmatrix}
2&8&2\\
-4&-17&-2\\
6&23&11
\end{pmatrix}.
\]
\item Use the $LU$ decomposition from part \textup{(a)} to solve
\[
\mathbf A\mathbf x=
\begin{pmatrix}
1\\1\\0
\end{pmatrix}.
\]
\end{problemsubparts}
\end{problem}

% 2022 Fall KAIST Introduction to Linear Algebra Midterm, Problem 3
\begin{problem}[2022 Fall \KAIST{} Introduction to Linear Algebra]
For which values of $a$ does the system
\[
\begin{aligned}
3x_1-8x_2-5x_3+10x_4&=9,\\
-x_1+x_2-3x_3-6x_4&=-2,\\
-2x_1+8x_2+20x_3+16x_4&=-3,\\
x_1-x_2+3x_3+a(a-1)x_4&=a-1
\end{aligned}
\]
have no solution, exactly one solution, or infinitely many solutions?
\end{problem}

% 2022 Fall KAIST Introduction to Linear Algebra Midterm, Problem 4
\begin{problem}[2022 Fall \KAIST{} Introduction to Linear Algebra]
Use elementary row operations to find the inverse of
\[
\mathbf A=
\begin{pmatrix}
6&18&3&0\\
16&33&8&-2\\
1&3&0&0\\
-8&-16&-4&1
\end{pmatrix}.
\]
\end{problem}

% 2022 Fall KAIST Introduction to Linear Algebra Midterm, Problem 5
\begin{problem}[2022 Fall \KAIST{} Introduction to Linear Algebra]
Let $\mathbf A\in\R^{n\times n}$ be nilpotent; that is, $\mathbf A^k=\mathbf0$ for some integer $k>0$.
\begin{problemsubparts}
\item Show that the characteristic polynomial of $\mathbf A$ is $\lambda^n$; equivalently,
\[
\det(\lambda\mathbf I_n-\mathbf A)=\lambda^n.
\]
\item Prove or disprove that $3\mathbf I_n+\mathbf A^2$ is invertible. If it is invertible, find its inverse.
\end{problemsubparts}
\end{problem}

% 2022 Fall KAIST Introduction to Linear Algebra Midterm, Problem 6
\begin{problem}[2022 Fall \KAIST{} Introduction to Linear Algebra]
Let $\mathbf A\in\R^{200\times100}$ satisfy
\[
\|\mathbf A\mathbf x\|=\|\mathbf x\|
\qquad\text{for every }\mathbf x\in\R^{100}.
\]
Using direct computation, show that the columns of $\mathbf A$ are orthonormal.
\end{problem}

% 2022 Fall KAIST Introduction to Linear Algebra Midterm, Problem 7
\begin{problem}[2022 Fall \KAIST{} Introduction to Linear Algebra]
\begin{problemsubparts}
\item Let $\mathbf A\in\R^{9\times9}$ satisfy $\det(\mathbf A)=9$. Find
\[
\det\bigl((\adj(\mathbf A))^{-1}\bigr).
\]
\item Let $\mathbf B\in\R^{n\times n}$ be invertible and have integer entries. Prove that every entry of $\mathbf B^{-1}$ is an integer if and only if
\[
\det(\mathbf B)=\pm1.
\]
\end{problemsubparts}
\end{problem}

% 2022 Fall KAIST Introduction to Linear Algebra Midterm, Problem 8
\begin{problem}[2022 Fall \KAIST{} Introduction to Linear Algebra]
Let
\[
\mathbf A=
\begin{pmatrix}
0&3&0\\
0&0&3\\
3&0&0
\end{pmatrix}.
\]
\begin{problemsubparts}
\item Find all complex eigenvalues and eigenvectors of $\mathbf A$.
\item Express the operator
\[
T_{\mathbf A}\colon\R^3\to\R^3,
\qquad
T_{\mathbf A}(\mathbf x)=\mathbf A\mathbf x,
\]
as a composition of special operators. Describe all relevant properties using, as applicable, the terms contraction, dilation, expansion, orthogonal, rotation, reflection, and shearing. For a rotation, state the axis and angle; for a reflection, state which vectors are reflected and which are fixed.
\end{problemsubparts}
\end{problem}

% 2022 Fall KAIST Introduction to Linear Algebra Midterm, Problem 9
\begin{problem}[2022 Fall \KAIST{} Introduction to Linear Algebra]
Let $T\colon\R^4\to\R^4$ be a linear operator, and let $\mathbf v_1,\dots,\mathbf v_4,\mathbf w_1,\dots,\mathbf w_4\in\R^4$. Define
\[
\mathbf u_i=
\begin{pmatrix}
T(\mathbf v_i)\cdot\mathbf w_1\\
T(\mathbf v_i)\cdot\mathbf w_2\\
T(\mathbf v_i)\cdot\mathbf w_3\\
T(\mathbf v_i)\cdot\mathbf w_4
\end{pmatrix},
\qquad i=1,2,3,4.
\]
Show that $\mathbf u_1,\mathbf u_2,\mathbf u_3,\mathbf u_4$ are linearly independent if and only if $T$ is injective and both
\[
\{\mathbf v_1,\mathbf v_2,\mathbf v_3,\mathbf v_4\}
\qquad\text{and}\qquad
\{\mathbf w_1,\mathbf w_2,\mathbf w_3,\mathbf w_4\}
\]
are linearly independent.

\textit{Hint.} Use the matrices
\[
\mathbf V=\begin{pmatrix}\mathbf v_1&\mathbf v_2&\mathbf v_3&\mathbf v_4\end{pmatrix},
\qquad
\mathbf W=\begin{pmatrix}\mathbf w_1&\mathbf w_2&\mathbf w_3&\mathbf w_4\end{pmatrix},
\]
and the matrix representation $[T]$.
\end{problem}

% 2022 Fall KAIST Introduction to Linear Algebra Midterm, Problem 10
\begin{problem}[2022 Fall \KAIST{} Introduction to Linear Algebra]
Let $W$ be the plane through the origin in $\R^3$ perpendicular to
\[
\mathbf u=
\begin{pmatrix}
1\\2\\3
\end{pmatrix},
\]
and let $\mathbf A$ be the standard matrix of reflection about $W$.
\begin{problemsubparts}
\item Find linearly independent vectors $\mathbf v,\mathbf w\in\R^3$ such that
\[
\mathbf A\mathbf v=\mathbf v,
\qquad
\mathbf A\mathbf w=\mathbf w.
\]
\item Find $\mathbf A$.
\end{problemsubparts}
\end{problem}

% 2022 Fall KAIST Introduction to Linear Algebra Midterm, Problem 11
\begin{problem}[2022 Fall \KAIST{} Introduction to Linear Algebra]
Consider the MATLAB code below.

{\small\ttfamily
\begin{tabbing}
\hspace{2em}\=\hspace{2em}\=\kill
I = 1:0.5:2.3;\\
b = zeros(10, 1);\\
for i = 2*I\\
\>b(i) = i\textasciicircum 2;\\
end\\[0.15em]
u = [2 1 5 3];\\
v = [1 2];\\[0.15em]
n = numel(u);\\
w = zeros(n, 1);\\
j = 0;\\
for i = 1:n\\
\>if u(i) > 2\\
\>\>w(j+1) = u(i);\\
\>\>j = j+1;\\
\>end\\
end\\
w = w(1:j);
\end{tabbing}
}

What are the outputs of the following commands?
\begin{problemsubparts}
\item \texttt{disp(b(1)+b(3)+b(5))}
\item \texttt{disp(u'*v)}
\item \texttt{disp(u.*[v v])}
\item \texttt{disp(w)}
\end{problemsubparts}
\end{problem}

% 2023 Spring POSTECH Applied Linear Algebra Midterm, Problem 1
\begin{problem}[2023 Spring \POSTECH{} Applied Linear Algebra]
Determine whether each of the following statements is true or false. You do not need to justify your answers.
\begin{problemchoices}
\item Let $\mathbf A$ and $\mathbf B$ be matrices. If both $\mathbf A\mathbf B$ and $\mathbf B\mathbf A$ are defined, then both products are square matrices.
\item Let $\mathbf A$ be a square matrix. Suppose that
\[
\mathbf A=\mathbf L\mathbf D\mathbf U,
\]
where $\mathbf L$ is lower triangular with all diagonal entries equal to $1$, $\mathbf U$ is upper triangular with all diagonal entries equal to $1$, and $\mathbf D$ is diagonal with no zero diagonal entries. Then such a decomposition is unique.
\item If $\mathbf A$ is a square matrix and $\mathbf A^2$ is nonsingular, then $\mathbf A$ is nonsingular.
\item The set of all $n\times n$ skew-symmetric matrices,
\[
\mathbf A^{\mathsf T}=-\mathbf A,
\]
forms a subspace of the vector space of $n\times n$ matrices.
\item Every subspace of $\R^4$ is the null space of some matrix.
\item Let $H_\theta\colon \R^2\to\R^2$ denote reflection across the line through the origin making angle $\theta$ with the positive $x$-axis. Then $H_{\theta_1}\circ H_{\theta_2}$ is counterclockwise rotation by $2(\theta_1-\theta_2)$.
\item There exists a nonsquare matrix that has both a left inverse and a right inverse, possibly different from each other.
\item Let $V$ and $W$ be finite-dimensional vector spaces. If a linear transformation $T\colon V\to W$ is injective, then
\[
\dim V\le \dim W.
\]
\item Let $T\colon \R^n\to\R^m$ be linear. If the values $T(\mathbf v)$ are known for $n$ distinct nonzero vectors $\mathbf v\in\R^n$, then $T(\mathbf v)$ is determined for every $\mathbf v\in\R^n$.
\item For a square matrix $\mathbf A$, if the columns of $\mathbf A$ are linearly independent, then the rows of $\mathbf A$ are linearly independent.
\end{problemchoices}
\end{problem}

% 2023 Spring POSTECH Applied Linear Algebra Midterm, Problem 2
\Needspace{11\baselineskip}
\begin{problem}[2023 Spring \POSTECH{} Applied Linear Algebra]
Let $\mathbf A=(a_{ij})_{n\times n}$ be the matrix defined by
\[
a_{ij}=
\begin{cases}
1, & j-i=1,\\
0, & \text{otherwise}.
\end{cases}
\]
\begin{problemsubparts}
\item Evaluate the sum of all entries of $\mathbf A$,
\[
\sum_{1\le i,j\le n} a_{ij}.
\]
\item Let $\mathbf B=(b_{ij})$ be
\[
\mathbf B=\sum_{k=0}^{n-1}\mathbf A^k
=\mathbf I_n+\mathbf A+\mathbf A^2+\cdots+\mathbf A^{n-1}.
\]
Evaluate the sum of all entries of $\mathbf B$,
\[
\sum_{1\le i,j\le n} b_{ij}.
\]
\end{problemsubparts}
\end{problem}

% 2023 Spring POSTECH Applied Linear Algebra Midterm, Problem 3
\begin{problem}[2023 Spring \POSTECH{} Applied Linear Algebra]
Consider
\[
\mathbf A=
\begin{pmatrix}
1&-1&0\\
2&0&-2\\
0&a&4
\end{pmatrix}.
\]
\begin{problemsubparts}
\item Find all $a\in\R$ for which $\mathbf A$ is invertible.
\item Using the Gauss--Jordan method, find $\mathbf A^{-1}$ when $a=2$.
\end{problemsubparts}
\end{problem}

% 2023 Spring POSTECH Applied Linear Algebra Midterm, Problem 4
\begin{problem}[2023 Spring \POSTECH{} Applied Linear Algebra]
Consider
\[
\mathbf A=
\begin{pmatrix}
1&2&3&2\\
2&4&8&6\\
-1&-2&1&2
\end{pmatrix},
\qquad
\mathbf b=
\begin{pmatrix}
b_1\\b_2\\b_3
\end{pmatrix}.
\]
In solving $\mathbf A\mathbf x=\mathbf b$, apply Gaussian elimination to obtain
\[
\mathbf U\mathbf x=\mathbf c,
\]
and then reduce further to
\[
\mathbf R\mathbf x=\mathbf d,
\]
where $\mathbf U$ is the echelon matrix obtained from $\mathbf A$ and $\mathbf R$ is the reduced row echelon form of $\mathbf A$. Compute $\mathbf c$ and $\mathbf d$ in terms of $b_1,b_2,b_3$.
\end{problem}

% 2023 Spring POSTECH Applied Linear Algebra Midterm, Problem 5
\begin{problem}[2023 Spring \POSTECH{} Applied Linear Algebra]
Let $V=M_2(\R)$ be the vector space of $2\times2$ real matrices, and let $P_3$ be the vector space of real polynomials of degree at most $3$. Consider the linear map $T\colon P_3\to V$ defined by
\[
T(a_0+a_1x+a_2x^2+a_3x^3)
=
\begin{pmatrix}
a_0+a_1&-a_0-a_2\\
-a_1+a_3&a_2-a_3
\end{pmatrix}.
\]
\begin{problemsubparts}
\item Find the matrix $\mathbf A$ representing $T$ with respect to the ordered basis
\[
(1,x,x^2,x^3)
\]
of $P_3$ and the ordered basis $(\mathbf E_1,\mathbf E_2,\mathbf E_3,\mathbf E_4)$ of $V$, where
\[
\mathbf E_1=
\begin{pmatrix}1&0\\0&0\end{pmatrix},\quad
\mathbf E_2=
\begin{pmatrix}0&1\\0&0\end{pmatrix},\quad
\mathbf E_3=
\begin{pmatrix}0&0\\1&0\end{pmatrix},\quad
\mathbf E_4=
\begin{pmatrix}0&0\\0&1\end{pmatrix}.
\]
\item Find $\rk(\mathbf A)$ and $\dim\Null(\mathbf A)$.
\end{problemsubparts}
\end{problem}

% 2023 Spring POSTECH Applied Linear Algebra Midterm, Problem 6
\begin{problem}[2023 Spring \POSTECH{} Applied Linear Algebra]
Let
\[
\mathbf A=
\begin{pmatrix}
1&-2\\
1&-1\\
1&0\\
1&2
\end{pmatrix},
\qquad
\mathbf b=
\begin{pmatrix}
1\\2\\1\\2
\end{pmatrix}.
\]
Decompose
\[
\mathbf b=\mathbf b_c+\mathbf b_n,
\]
where $\mathbf b_c\in\Col(\mathbf A)$ and $\mathbf b_n\in\Null(\mathbf A^{\mathsf T})$, in the following two ways.
\begin{problemsubparts}
\item Solve the normal equation
\[
\mathbf A^{\mathsf T}\mathbf A\widehat{\mathbf x}=\mathbf A^{\mathsf T}\mathbf b
\]
and compute $\mathbf b_c=\mathbf A\widehat{\mathbf x}$.
\item Find an orthonormal basis $\{\mathbf q_1,\mathbf q_2\}$ for $\Col(\mathbf A)$ and express $\mathbf b_c$ as a linear combination of $\mathbf q_1$ and $\mathbf q_2$.
\end{problemsubparts}
\end{problem}

% 2023 Spring POSTECH Applied Linear Algebra Midterm, Problem 7
\begin{problem}[2023 Spring \POSTECH{} Applied Linear Algebra]
Suppose that data are expected to conform to the quadratic relation
\[
y(t)=A+Bt+Ct^2.
\]
The observed data are
\[
\begin{array}{c|ccccc}
t&-2&-1&0&1&2\\ \hline
y(t)&-4&1&3&4&6
\end{array}.
\]
Find the parameters $A,B,C$ that minimize the sum of squared errors.
\end{problem}

% 2023 Spring POSTECH Applied Linear Algebra Midterm, Problem 8
\begin{problem}[2023 Spring \POSTECH{} Applied Linear Algebra]
Let $V$ be a subspace of $\R^d$. Prove that if
\[
\{\mathbf v_1,\dots,\mathbf v_m\}
\qquad\text{and}\qquad
\{\mathbf w_1,\dots,\mathbf w_n\}
\]
are both bases for $V$, then $m=n$.
\end{problem}

% 2023 Spring POSTECH Applied Linear Algebra Midterm, Problem 9
\begin{problem}[2023 Spring \POSTECH{} Applied Linear Algebra]
Let $\mathbf A$ be an $m\times n$ real matrix. Prove that
\[
\rk(\mathbf A)=\rk(\mathbf A^{\mathsf T}\mathbf A)=\rk(\mathbf A\mathbf A^{\mathsf T}).
\]
\end{problem}

% 2023 Spring POSTECH Applied Linear Algebra Final, Problem 1
\begin{problem}[2023 Spring \POSTECH{} Applied Linear Algebra]
Determine whether each of the following statements is true or false. You do not need to justify your answers.
\begin{problemchoices}
\item There exist $3\times3$ matrices $\mathbf A$ and $\mathbf B$ such that
\[
\mathbf A\mathbf B-\mathbf B\mathbf A=\mathbf I_3.
\]
\item If $\mathbf A$ and $\mathbf B$ are simultaneously diagonalizable, and $\mathbf B$ and $\mathbf C$ are simultaneously diagonalizable, then $\mathbf A$ and $\mathbf C$ are simultaneously diagonalizable.
\item There exists a real $2\times2$ matrix $\mathbf A\ne\mathbf I_2$ such that
\[
\mathbf A^3=\mathbf I_2.
\]
\item If a $3\times3$ matrix $\mathbf A$ has eigenvalues $1,2,3$, then
\[
(\mathbf A-\mathbf I_3)(\mathbf A-2\mathbf I_3)(\mathbf A-3\mathbf I_3)
\]
cannot have a nonzero eigenvalue.
\item If a real matrix $\mathbf A$ is normal, then
\[
(\mathbf A^k)^+=(\mathbf A^+)^k
\]
for every positive integer $k$.
\item If $\det(\mathbf A)>0$ and $\det(\mathbf B)>0$, then
\[
\det(\mathbf A+\mathbf B)>0.
\]
\item If $\mathbf A=(a_{ij})$ is an $n\times n$ matrix satisfying
\[
\sum_{i=1}^n a_{ij}=5
\]
for every $j$, then $5$ is an eigenvalue of $\mathbf A$.
\item If $\mathbf A$ is an $n\times n$ symmetric matrix, then $e^{\mathbf A}$ is positive definite.
\item Let $\mathbf A$ and $\mathbf B$ be $n\times n$ matrices. If $\lambda$ is an eigenvalue of $\mathbf A$ and $\mu$ is an eigenvalue of $\mathbf B$, then $\lambda+\mu$ is an eigenvalue of $\mathbf A+\mathbf B$.
\item If a real symmetric matrix $\mathbf A$ is positive definite, then it is nonsingular.
\end{problemchoices}
\end{problem}

% 2023 Spring POSTECH Applied Linear Algebra Final, Problem 2
\begin{problem}[2023 Spring \POSTECH{} Applied Linear Algebra]
\begin{problemsubparts}
\item Find the volume of the parallelepiped determined by the edge vectors from $(0,0,0)$ to
\[
(-1,2,2),\qquad (2,-1,2),\qquad (2,2,-1).
\]
\item Compute
\[
\det
\begin{pmatrix}
1&1&1&1\\
1&-1&1&-1\\
1&1&-1&-1\\
1&-1&-1&1
\end{pmatrix}.
\]
\end{problemsubparts}
\end{problem}

% 2023 Spring POSTECH Applied Linear Algebra Final, Problem 3
\begin{problem}[2023 Spring \POSTECH{} Applied Linear Algebra]
Let $\mathbf A,\mathbf B,\mathbf C$ be $n\times n$ matrices and let $\mathbf 0$ be the $n\times n$ zero matrix. Show that
\[
\det
\begin{pmatrix}
\mathbf A&\mathbf B\\
\mathbf 0&\mathbf C
\end{pmatrix}
=\det(\mathbf A)\det(\mathbf C).
\]
\textit{Hint.} One possible approach is to use the Leibniz formula for the determinant.
\end{problem}

% 2023 Spring POSTECH Applied Linear Algebra Final, Problem 4
\begin{problem}[2023 Spring \POSTECH{} Applied Linear Algebra]
Let $(a_n)_{n=0}^{\infty}$ and $(b_n)_{n=0}^{\infty}$ be sequences satisfying
\[
a_{n+1}=3a_n+2b_n,
\qquad
b_{n+1}=a_n+4b_n,
\qquad
a_0=1,
\qquad
b_0=0.
\]
Evaluate
\[
\lim_{n\to\infty}\frac{a_n-b_n}{a_n+b_n}.
\]
\end{problem}

% 2023 Spring POSTECH Applied Linear Algebra Final, Problem 5
\begin{problem}[2023 Spring \POSTECH{} Applied Linear Algebra]
Let
\[
\mathbf A=
\begin{pmatrix}
1&1&2\\
0&1&1\\
0&0&1
\end{pmatrix},
\qquad
\mathbf B=
\begin{pmatrix}
1&0&0\\
1&1&0\\
2&1&1
\end{pmatrix},
\qquad
\mathbf C=
\begin{pmatrix}
1&0&0\\
1&1&1\\
0&0&1
\end{pmatrix}.
\]
Find two matrices among $\mathbf A,\mathbf B,\mathbf C$ that are similar to each other. Justify your answer.
\end{problem}

% 2023 Spring POSTECH Applied Linear Algebra Final, Problem 6
\begin{problem}[2023 Spring \POSTECH{} Applied Linear Algebra]
Let
\[
\mathbf A=
\begin{pmatrix}
3&2&2\\
2&3&-2
\end{pmatrix}.
\]
\begin{problemsubparts}
\item Find the singular value decomposition of $\mathbf A$.
\item Find the pseudoinverse $\mathbf A^+$ of $\mathbf A$.
\end{problemsubparts}
\end{problem}

% 2023 Spring POSTECH Applied Linear Algebra Final, Problem 7
\begin{problem}[2023 Spring \POSTECH{} Applied Linear Algebra]
\begin{problemsubparts}
\item Find all $\alpha\in\R$ such that
\[
\mathbf A=
\begin{pmatrix}
4&\alpha\\
\alpha&5
\end{pmatrix}
\]
is positive definite.
\item For a real number $\alpha$ such that $\mathbf A$ is positive definite, let $d_1(\alpha)$ and $d_2(\alpha)$ be the distances from the origin to, respectively, the closest and farthest points on
\[
4x^2+2\alpha xy+5y^2=1.
\]
Find the minimum value of
\[
f(\alpha)=d_1(\alpha)^2+d_2(\alpha)^2.
\]
\end{problemsubparts}
\end{problem}

% 2023 Spring POSTECH Applied Linear Algebra Final, Problem 8
\begin{problem}[2023 Spring \POSTECH{} Applied Linear Algebra]
Let $\mathbf A$ be a square matrix with real entries. Show that
\[
\tr(\mathbf A^{\mathsf T}\mathbf A)=0
\]
if and only if $\mathbf A=\mathbf 0$.
\end{problem}

% 2023 Spring POSTECH Applied Linear Algebra Final, Problem 9
\begin{problem}[2023 Spring \POSTECH{} Applied Linear Algebra]
For an $n\times n$ matrix $\mathbf A$, show that
\[
\det(e^{\mathbf A})=e^{\tr(\mathbf A)}.
\]
\textit{Hint.} By Schur's triangularization theorem over $\C$, there exists an invertible matrix $\mathbf M$ such that $\mathbf M^{-1}\mathbf A\mathbf M$ is upper triangular.
\end{problem}

% 2023 Spring KAIST Introduction to Linear Algebra Midterm, Problem 1
\begin{problem}[2023 Spring \KAIST{} Introduction to Linear Algebra]
Let $L$ be the line through $(1,1,1)$ and $(2,4,3)$.
\begin{problemsubparts}
\item Find a parametric equation of the plane $P$ through $Q=(-1,2,3)$ that is perpendicular to $L$.
\item Find the intersection point $B$ of $P$ and $L$, and find the distance between $Q$ and $B$.
\end{problemsubparts}
\end{problem}

% 2023 Spring KAIST Introduction to Linear Algebra Midterm, Problem 2
\begin{problem}[2023 Spring \KAIST{} Introduction to Linear Algebra]
For $a\in\R$, consider the linear system
\[
\begin{pmatrix}
1&2&6a\\
4&1&4\\
7&0&-a^2
\end{pmatrix}
\begin{pmatrix}x\\y\\z\end{pmatrix}
=
\begin{pmatrix}-a\\a^2-1\\-a+10\end{pmatrix}.
\]
For each value of $a$, determine the number of solutions of the system.
\end{problem}

% 2023 Spring KAIST Introduction to Linear Algebra Midterm, Problem 3
\begin{problem}[2023 Spring \KAIST{} Introduction to Linear Algebra]
Let
\[
\mathbf u=
\begin{pmatrix}1\\2\\10\\100\end{pmatrix},
\qquad
\mathbf v=
\begin{pmatrix}0\\1\\0\\-1\end{pmatrix},
\]
and let $\mathbf e_i^{\mathsf T}$ be the $i$th standard basis row vector of $\R^4$. Compute each of the following whenever it is defined.
\begin{problemchoices}
\item $\tr(\mathbf u\mathbf v^{\mathsf T})$.
\item $\tr(\mathbf u^{\mathsf T}\mathbf v)$.
\item $\mathbf e_2^{\mathsf T}\mathbf u\mathbf v^{\mathsf T}$.
\item $\mathbf e_2\mathbf u\mathbf v^{\mathsf T}$.
\item $\mathbf u\mathbf v^{\mathsf T}\mathbf e_2^{\mathsf T}$.
\item $\mathbf u\mathbf v^{\mathsf T}\mathbf e_2$.
\end{problemchoices}
\end{problem}

% 2023 Spring KAIST Introduction to Linear Algebra Midterm, Problem 4
\begin{problem}[2023 Spring \KAIST{} Introduction to Linear Algebra]
Let $\mathbf A\in\R^{3\times3}$ be invertible and satisfy
\[
\mathbf A^2=2\mathbf A.
\]
Find all possible values of $\det(\mathbf A)$, and justify both the possible and impossible values.
\end{problem}

% 2023 Spring KAIST Introduction to Linear Algebra Midterm, Problem 5
\begin{problem}[2023 Spring \KAIST{} Introduction to Linear Algebra]
Let $\mathbf e_i$ be the $i$th standard basis vector of $\R^n$. For arbitrary linear operators $T_1,T_2\colon\R^n\to\R^n$, determine which of the following functions are always linear. Justify your answers.
\begin{problemchoices}
\item $f_1(\mathbf x)=T_1(\mathbf x)+T_2(\mathbf x)$.
\item $f_2(\mathbf x)=T_1(\mathbf x)+\mathbf e_1$.
\item $f_3(\mathbf x)=T_1(\mathbf x)\cdot T_2(\mathbf x)$.
\item $f_4(\mathbf x)=\mathbf e_1\cdot T_1(\mathbf x)$.
\end{problemchoices}
\end{problem}

% 2023 Spring KAIST Introduction to Linear Algebra Midterm, Problem 6
\begin{problem}[2023 Spring \KAIST{} Introduction to Linear Algebra]
Let $T\colon\R^n\to\R^n$ be a linear operator. Prove that the following statements are equivalent:
\begin{problemclauses}
\item $\|T(\mathbf x)\|=\|\mathbf x\|$ for every $\mathbf x\in\R^n$.
\item $T(\mathbf x)\cdot T(\mathbf y)=\mathbf x\cdot\mathbf y$ for every $\mathbf x,\mathbf y\in\R^n$.
\end{problemclauses}
Then use this result to show that every length-preserving linear operator $T\colon\R^n\to\R^n$ preserves angles.
\end{problem}

% 2023 Spring KAIST Introduction to Linear Algebra Midterm, Problem 7
\begin{problem}[2023 Spring \KAIST{} Introduction to Linear Algebra]
Determine whether each statement is true or false. A false statement should be disproved by a counterexample.
\begin{problemchoices}
\item An $n\times n$ matrix $\mathbf A$ is invertible if and only if a row echelon form of $\mathbf A$ obtained by Gaussian elimination is upper triangular with every diagonal entry equal to $1$.
\item Every subspace $W\subseteq\R^n$ contains the zero vector.
\item If $\mathbf A$ and $\mathbf B$ are symmetric $n\times n$ matrices, then $\mathbf A\mathbf B$ is symmetric.
\item For a fixed nonzero vector $\mathbf v\in\R^n$, the function $T\colon\R^n\to\R^n$ defined by $T(\mathbf x)=\mathbf x+\mathbf v$ is a linear operator.
\item If $\mathbf A\in\R^{m\times n}$ and $n>m$, then $\Null(\mathbf A)$ contains a nonzero vector.
\end{problemchoices}
\end{problem}

% 2023 Spring KAIST Introduction to Linear Algebra Midterm, Problem 8
\begin{problem}[2023 Spring \KAIST{} Introduction to Linear Algebra]
Let
\[
\mathbf A=
\begin{pmatrix}
1&2&3\\
2&3&4\\
3&4&5
\end{pmatrix}.
\]
\begin{problemsubparts}
\item Determine whether $\mathbf A$ admits an $LU$ decomposition. If it does, find one; otherwise, find a $PLU$ decomposition.
\item Compute $\det(\mathbf A)$.
\end{problemsubparts}
\end{problem}

% 2023 Spring KAIST Introduction to Linear Algebra Midterm, Problem 9
\begin{problem}[2023 Spring \KAIST{} Introduction to Linear Algebra]
Let $\mathbf A\in\R^{n\times n}$ with $n\ge2$. Express
\[
\det\bigl(\adj(\mathbf A)\bigr)
\]
in terms of $n$ and $\det(\mathbf A)$.
\end{problem}

% 2023 Spring KAIST Introduction to Linear Algebra Midterm, Problem 10
\begin{problem}[2023 Spring \KAIST{} Introduction to Linear Algebra]
Let
\[
\mathbf A=
\begin{pmatrix}
0&2&1\\
2&-3&-2\\
-1&2&2
\end{pmatrix}.
\]
\begin{problemsubparts}
\item Find all eigenvalues of $\mathbf A$, their algebraic multiplicities, and the corresponding eigenvectors.
\item Find $3\mathbf A^{-1}$ if it exists.
\end{problemsubparts}
\end{problem}

% 2023 Spring KAIST Introduction to Linear Algebra Midterm, Problem 11
\begin{problem}[2023 Spring \KAIST{} Introduction to Linear Algebra]
A square matrix $\mathbf A$ is called idempotent if $\mathbf A^2=\mathbf A$.
\begin{problemsubparts}
\item Determine all possible eigenvalues of an idempotent matrix. For each possible eigenvalue, give an example of an idempotent matrix having that eigenvalue.
\item Can a single idempotent matrix have all of the possible eigenvalues found in part \textup{(a)}? If so, give an example.
\end{problemsubparts}
\end{problem}

% 2023 Spring KAIST Introduction to Linear Algebra Midterm, Problem 12
\begin{problem}[2023 Spring \KAIST{} Introduction to Linear Algebra]
Let
\[
\mathbf A=
\frac13
\begin{pmatrix}
1&-2&-2\\
-2&1&-2\\
-2&-2&1
\end{pmatrix}.
\]
\begin{problemsubparts}
\item Show that $\mathbf A$ is orthogonal.
\item Find all eigenvalues of $\mathbf A$ and the eigenspace corresponding to each eigenvalue.
\item Determine whether the orthogonal linear operator represented by $\mathbf A$ is a rotation about a line through the origin, a reflection about a plane through the origin, or neither. In the rotation case, find a direction vector of the axis and the rotation angle. In the reflection case, find a normal vector of the reflection plane.
\end{problemsubparts}
\end{problem}

% 2023 Spring KAIST Introduction to Linear Algebra Midterm, Problem 13
\begin{problem}[2023 Spring \KAIST{} Introduction to Linear Algebra]
Let $\mathbf X,\mathbf Y\in\R^{3\times3}$ be invertible.
\begin{problemsubparts}
\item Determine all $c\in\R$ for which
\[
\mathbf X\mathbf Y-\mathbf Y\mathbf X=c\mathbf I_3
\]
can hold. For each possible value, give an example of $\mathbf X$ and $\mathbf Y$; for each impossible value, explain why it is impossible.
\item Determine all $c\in\R$ for which
\[
\mathbf X\mathbf Y\mathbf X^{-1}\mathbf Y^{-1}=c\mathbf I_3
\]
can hold. For each possible value, give an example of $\mathbf X$ and $\mathbf Y$; for each impossible value, explain why it is impossible.
\end{problemsubparts}
\end{problem}

% 2023 Spring KAIST Introduction to Linear Algebra Midterm, Problem 14
\begin{problem}[2023 Spring \KAIST{} Introduction to Linear Algebra]
The following commands are entered in the MATLAB Command Window.

{\small\ttfamily
\begin{tabbing}
A = magic(4);\\
A\_times\_ones = A * [1; 1; 1; 1];\\
sum\_row2 = A\_times\_ones(2);\\
disp(sum\_row2);
\end{tabbing}
}

The MATLAB function \texttt{magic} takes a positive integer $n$ and returns an $n\times n$ matrix containing the integers $1,2,\dots,n^2$ such that every row sum and every column sum is the same.

\begin{problemsubparts}
\item What is the output of the code? Explain why.
\item Assuming the code above has already been executed, write MATLAB code that creates a variable \texttt{B} whose value is the product of \texttt{A} with itself.
\item Assuming the code from part \textup{(b)} has also been executed, the following commands are run:

{\small\ttfamily
\begin{tabbing}
B\_times\_ones = B * [1, 1, 1, 1]';\\
disp(B\_times\_ones);
\end{tabbing}
}

A column vector whose entries are all $1156$ is displayed. Explain why. \textit{Hint.} $1156=34^2$.
\end{problemsubparts}

Now suppose the following MATLAB function is available:

{\small\ttfamily
\begin{tabbing}
\hspace{2em}\=\hspace{2em}\=\kill
function y = my\_func(x)\\
\>if x > 0\\
\>\>y = sqrt(x);\\
\>elseif x < 0\\
\>\>y = x\string^2;\\
\>else\\
\>\>y = 0;\\
\>end\\
end
\end{tabbing}
}

\begin{problemsubparts}[resume]
\item Assuming all code through part \textup{(c)} has already been executed, what is the output of

{\small\ttfamily
\begin{tabbing}
x = [1, 1, 1, 1] * B\_times\_ones;\\
disp(my\_func(-3) + my\_func(x));
\end{tabbing}
}
\end{problemsubparts}
\end{problem}

% 2023 Spring KAIST Introduction to Linear Algebra Final, Problem 1
\begin{problem}[2023 Spring \KAIST{} Introduction to Linear Algebra]
Let
\[
\mathbf p_1=
\begin{pmatrix}1\\1\\1\end{pmatrix},
\qquad
\mathbf p_2=
\begin{pmatrix}2\\0\\1\end{pmatrix},
\qquad
\mathbf p_3=
\begin{pmatrix}1\\5\\x\end{pmatrix}
\]
be three points in $\R^3$. For which values of $x$ does there exist a two-dimensional linear subspace containing all three points? Justify your answer.
\end{problem}

% 2023 Spring KAIST Introduction to Linear Algebra Final, Problem 2
\begin{problem}[2023 Spring \KAIST{} Introduction to Linear Algebra]
Let $T\colon\R^n\to\R^n$ be a linear operator, and let
\[
\{\mathbf v_1,\dots,\mathbf v_k\}\subseteq\R^n
\]
be linearly independent. Show that if
\[
\{T(\mathbf v_1),\dots,T(\mathbf v_k)\}
\]
is not linearly independent, then $T$ is not injective. Equivalently, show that there exist distinct $\mathbf x,\mathbf y\in\R^n$ such that
\[
T(\mathbf x)=T(\mathbf y).
\]
\end{problem}

% 2023 Spring KAIST Introduction to Linear Algebra Final, Problem 3
\begin{problem}[2023 Spring \KAIST{} Introduction to Linear Algebra]
Let $\mathbf A$ be row equivalent to
\[
\mathbf R=
\begin{pmatrix}
0&1&0&0&4&0\\
0&0&1&1&7&0\\
0&0&0&0&0&1\\
0&0&0&0&0&0\\
0&0&0&0&0&0
\end{pmatrix}.
\]
For each item, either answer the question or state that there is not enough information.
\begin{problemsubparts}
\item Find $\dim\Null(\mathbf A)$.
\item Find a basis for $\Null(\mathbf A)$.
\item Find $\dim\Col(\mathbf A)$.
\item Find a basis for $\Col(\mathbf A)$.
\item Find $\dim\Row(\mathbf A)$.
\item Find a basis for $\Row(\mathbf A)$.
\end{problemsubparts}
\end{problem}

% 2023 Spring KAIST Introduction to Linear Algebra Final, Problem 4
\begin{problem}[2023 Spring \KAIST{} Introduction to Linear Algebra]
Let
\[
\mathcal B=
\left\{
\mathbf b_1=
\begin{pmatrix}1\\1\\1\\1\end{pmatrix},
\mathbf b_2=
\begin{pmatrix}2\\3\\1\\2\end{pmatrix},
\mathbf b_3=
\begin{pmatrix}5\\2\\4\\1\end{pmatrix}
\right\}
\]
be a basis for a subspace $W\subseteq\R^4$.
\begin{problemsubparts}
\item Find an orthogonal basis for $W$ that contains $\mathbf b_1$ and whose vectors all have integer coordinates.
\item Find a basis for $W^\perp$.
\item Find the orthogonal projection onto $W$ of
\[
\mathbf x=
\begin{pmatrix}2\\3\\5\\7\end{pmatrix}.
\]
\end{problemsubparts}
\end{problem}

% 2023 Spring KAIST Introduction to Linear Algebra Final, Problem 5
\begin{problem}[2023 Spring \KAIST{} Introduction to Linear Algebra]
Let
\[
V=\Null(\mathbf A)\subseteq\R^n
\]
for some matrix $\mathbf A$, and suppose $\mathbf v_1,\dots,\mathbf v_k\in V$ are linearly independent. Form a matrix $\mathbf B$ by appending the rows
\[
\mathbf v_1^{\mathsf T},\dots,\mathbf v_k^{\mathsf T}
\]
to the bottom of $\mathbf A$. If
\[
\{\mathbf w_1,\dots,\mathbf w_m\}
\]
is a basis for $\Null(\mathbf B)$, prove that
\[
\{\mathbf v_1,\dots,\mathbf v_k,\mathbf w_1,\dots,\mathbf w_m\}
\]
is a basis for $V$.
\end{problem}

% 2023 Spring KAIST Introduction to Linear Algebra Final, Problem 6
\begin{problem}[2023 Spring \KAIST{} Introduction to Linear Algebra]
Find an orthogonal diagonalization of
\[
\mathbf A=
\begin{pmatrix}
2&1&-1\\
1&2&1\\
-1&1&2
\end{pmatrix}.
\]
\end{problem}

% 2023 Spring KAIST Introduction to Linear Algebra Final, Problem 7
\begin{problem}[2023 Spring \KAIST{} Introduction to Linear Algebra]
Let
\[
\mathbf A=
\begin{pmatrix}
1&3\\
2&6
\end{pmatrix},
\qquad
\mathbf b=
\begin{pmatrix}
6\\7
\end{pmatrix}.
\]
\begin{problemsubparts}
\item Find all least-squares solutions of $\mathbf A\mathbf x=\mathbf b$.
\item Find the least-squares solution $\mathbf y$ of $\mathbf A\mathbf x=\mathbf b$ with minimum norm $\|\mathbf y\|$.
\end{problemsubparts}
\end{problem}

% 2023 Spring KAIST Introduction to Linear Algebra Final, Problem 8
\begin{problem}[2023 Spring \KAIST{} Introduction to Linear Algebra]
Find a singular value decomposition of
\[
\mathbf A=
\begin{pmatrix}
2&3\\
0&2
\end{pmatrix}.
\]
\end{problem}

% 2023 Spring KAIST Introduction to Linear Algebra Final, Problem 9
\begin{problem}[2023 Spring \KAIST{} Introduction to Linear Algebra]
Consider
\[
\mathbf A=
\begin{pmatrix}
a&b&c\\
\frac1{\sqrt2}&-\frac1{\sqrt2}&0\\
\frac1{\sqrt3}&\frac1{\sqrt3}&\frac1{\sqrt3}
\end{pmatrix},
\qquad
\mathbf B=
\begin{pmatrix}
d&e&f\\
\frac1{\sqrt2}&-\frac1{\sqrt2}&0\\
0&0&1
\end{pmatrix}.
\]
\begin{problemsubparts}
\item For which values of $a,b,c,d,e,f$ are $\mathbf A$ and $\mathbf B$ orthogonal matrices representing rotations?
\item When both matrices are orthogonal and represent rotations, let
\[
\mathbf C=\mathbf A^{\mathsf T}\mathbf B\mathbf A.
\]
Show that $\mathbf C$ represents a rotation about a line through the origin in $\R^3$, and find its axis and angle. You need not determine the orientation of the axis or the direction of rotation.
\end{problemsubparts}
\end{problem}

% 2023 Spring KAIST Introduction to Linear Algebra Final, Problem 10
\begin{problem}[2023 Spring \KAIST{} Introduction to Linear Algebra]
Let $\mathbf A\in\R^{n\times n}$ have rank $1$, where $n>1$.
\begin{problemsubparts}
\item Show directly, without invoking the corresponding rank-factorization theorem, that there exist nonzero vectors $\mathbf u,\mathbf v\in\R^n$ such that
\[
\mathbf A=\mathbf u\mathbf v^{\mathsf T}.
\]
\item Find all eigenvalues of $\mathbf A$, together with their algebraic multiplicities, in terms of $\mathbf u$ and $\mathbf v$. You may need to consider two cases.
\item Give a necessary and sufficient condition, in terms of $\mathbf u$ and $\mathbf v$, for $\mathbf A=\mathbf u\mathbf v^{\mathsf T}$ to be diagonalizable.
\end{problemsubparts}
\end{problem}

% 2023 Spring KAIST Introduction to Linear Algebra Final, Problem 11
\begin{problem}[2023 Spring \KAIST{} Introduction to Linear Algebra]
Let $\mathbf A,\mathbf B\in\R^{n\times n}$. Suppose $\mathcal B_1$ and $\mathcal B_2$ are orthonormal bases for $\R^n$, where $\mathcal B_1$ consists of eigenvectors of $\mathbf A$ and $\mathcal B_2$ consists of eigenvectors of $\mathbf B$. Show that there exists an orthonormal basis for $\R^n$ consisting of eigenvectors of $\mathbf A+\mathbf B$.
\end{problem}

% 2023 Spring KAIST Introduction to Linear Algebra Final, Problem 12
\begin{problem}[2023 Spring \KAIST{} Introduction to Linear Algebra]
Let $Q(\mathbf x)$ be a quadratic form on $\R^3$, where
\[
\mathbf x=
\begin{pmatrix}
x_1\\x_2\\x_3
\end{pmatrix}.
\]
\begin{problemsubparts}
\item Show that there exists a rotation matrix $\mathbf R=\mathbf R_{\mathbf u,\theta}$, depending on $Q$, such that $Q(\mathbf R\mathbf x)$ is diagonal; that is,
\[
Q(\mathbf R\mathbf x)=\lambda_1x_1^2+\lambda_2x_2^2+\lambda_3x_3^2
\]
for some $\lambda_1,\lambda_2,\lambda_3\in\R$.
\item Let
\[
Q(\mathbf x)=-x_1^2+x_2^2+4x_1x_3+4x_2x_3.
\]
Find a rotation matrix $\mathbf R$ such that $Q(\mathbf R\mathbf x)$ is diagonal. You need not determine the axis or angle of rotation.
\end{problemsubparts}
\end{problem}

% 2023 Spring KAIST Introduction to Linear Algebra Final, Problem 13
\begin{problem}[2023 Spring \KAIST{} Introduction to Linear Algebra]
Find all $n\times n$ upper triangular orthogonal matrices.
\end{problem}

% 2023 Spring KAIST Introduction to Linear Algebra Final, Problem 14
\begin{problem}[2023 Spring \KAIST{} Introduction to Linear Algebra]
Let $\mathbf A\in\R^{n\times n}$ be invertible. Show that there exist a unique orthogonal matrix $\mathbf Q$ and a unique upper triangular matrix $\mathbf R$ such that
\[
\mathbf A=\mathbf Q\mathbf R
\]
and every diagonal entry of $\mathbf R$ is positive.
\end{problem}

% 2023 Spring KAIST Introduction to Linear Algebra Final, Problem 15
\begin{problem}[2023 Spring \KAIST{} Introduction to Linear Algebra]
The following commands are entered in the MATLAB Command Window.

{\small\ttfamily
\begin{tabbing}
\hspace{2em}\=\hspace{2em}\=\kill
A = zeros(7);\\
for i = [1, 3, 5, 7]\\
\>for j = [1, 3, 5, 7]\\
\>\>A(i, j) = 1;\\
\>end\\
end\\
for i = [2, 4, 6]\\
\>for j = [2, 4, 6]\\
\>\>A(i, j) = 1;\\
\>end\\
end
\end{tabbing}
}

\begin{problemsubparts}
\item Some columns of \texttt{A} are eigenvectors of $\mathbf A$. Identify those columns and compute the corresponding eigenvalues.
\item Suppose the following commands are then executed:

{\small\ttfamily
\begin{tabbing}
eigvals = eig(A);\\
disp(max(eigvals));
\end{tabbing}
}

Here \texttt{max} returns the maximum element of its input. What is displayed? Explain why.
\end{problemsubparts}

Suppose the MATLAB function \texttt{householder}, which takes a column vector \texttt{x}, is defined by

{\small\ttfamily
\begin{tabbing}
\hspace{2em}\=\hspace{2em}\=\kill
function Q = householder(x)\\
\>[n, m] = size(x);\\
\>if x' * x == 0\\
\>\>Q = eye(n);\\
\>else\\
\>\>Q = eye(n) - 2 * (x * x') / (x' * x);\\
\>end\\
end
\end{tabbing}
}

Assuming all previous code has already been executed, the following commands are then run:

{\small\ttfamily
\begin{tabbing}
H = householder([-1; 0; 1; 0; -1; 0; 1]);\\
B = H * A * H;\\
syms t;\\
y = subs(det(t * eye(7) - B), t, 1);\\
{}[L, U, P] = lu(A);\\
{}[Q, R] = qr(U);\\
v = diag(inv(P) * L * inv(Q) * R);
\end{tabbing}
}

\begin{problemsubparts}[resume]
\item What is the value of \texttt{y} after executing the code?
\item What is the value of \texttt{v} after executing the code?
\end{problemsubparts}
\end{problem}


% 2024 Spring KAIST Introduction to Linear Algebra Midterm, Problem 1
\begin{problem}[2024 Spring \KAIST{} Introduction to Linear Algebra]
Determine whether each of the following sets is a subspace of $\R^4$. No justification is required.
\begin{problemchoices}
\item $\Span((1,2,3,4),(0,2,0,1))\cup\Span((1,0,3,0))$.
\item $\Span((1,1,3,-2),(2,3,1,-1))\cup\Span((1,3,-7,4))$.
\item The set of vectors in $\R^4$ whose entries sum to $0$.
\item The set of vectors in $\Span((1,1,1,-1),(-1,1,1,1))$ whose entries sum to $100$.
\item The set of vectors in $\R^4$ that are orthogonal to $(1,2,3,-4)$ and whose entries sum to $0$.
\end{problemchoices}
\end{problem}

% 2024 Spring KAIST Introduction to Linear Algebra Midterm, Problem 2
\begin{problem}[2024 Spring \KAIST{} Introduction to Linear Algebra]
Let $\mathbf A\in\R^{3\times3}$ satisfy
\[
\bigl(2\mathbf A^{\mathsf T}\bigr)^{-1}=
\begin{pmatrix}
3&0&2\\
-1&1&-1\\
-1&0&0
\end{pmatrix}.
\]
\begin{problemsubparts}
\item Find $\mathbf A$.
\item Find all eigenvalues of $\mathbf A$, their algebraic multiplicities, and the corresponding eigenspaces.
\end{problemsubparts}
\end{problem}

% 2024 Spring KAIST Introduction to Linear Algebra Midterm, Problem 3
\begin{problem}[2024 Spring \KAIST{} Introduction to Linear Algebra]
For $a\in\R$, consider the linear system
\[
\begin{pmatrix}
1&-2&4\\
-2&2&-3a\\
a&-8&a^2+2
\end{pmatrix}
\begin{pmatrix}x\\y\\z\end{pmatrix}
=
\begin{pmatrix}a-1\\-7\\a\end{pmatrix}.
\]
\begin{problemsubparts}
\item Determine the number of solutions for each $a\in\R$.
\item Use Cramer's rule to solve the system when $a=0$. Do not use another method.
\end{problemsubparts}
\end{problem}

% 2024 Spring KAIST Introduction to Linear Algebra Midterm, Problem 4
\begin{problem}[2024 Spring \KAIST{} Introduction to Linear Algebra]
Define $T\colon\R^4\to\R^4$ by
\[
T(a,b,c,d)=(y_1,y_2,y_3,y_4),
\]
where
\[
\begin{pmatrix}y_1&y_2\\y_3&y_4\end{pmatrix}
=
\begin{pmatrix}1&2\\0&1\end{pmatrix}
\begin{pmatrix}a&b\\c&d\end{pmatrix}
-
\begin{pmatrix}a&b\\c&d\end{pmatrix}
\begin{pmatrix}1&2\\0&1\end{pmatrix}.
\]
\begin{problemsubparts}
\item Show that $T$ is a linear transformation.
\item Find $\ker T$ and $\det([T]_{\mathcal E})$, where $\mathcal E$ is the standard ordered basis of $\R^4$.
\end{problemsubparts}
\end{problem}

% 2024 Spring KAIST Introduction to Linear Algebra Midterm, Problem 5
\begin{problem}[2024 Spring \KAIST{} Introduction to Linear Algebra]
Let
\[
\mathbf A=
\begin{pmatrix}
1&1&1\\
2&3&4\\
4&5&6
\end{pmatrix},
\qquad
\mathbf v_1=\begin{pmatrix}1\\1\\1\end{pmatrix},\quad
\mathbf v_2=\begin{pmatrix}2\\3\\4\end{pmatrix},\quad
\mathbf v_3=\begin{pmatrix}4\\5\\6\end{pmatrix}.
\]
\begin{problemsubparts}
\item Find an $LU$ decomposition of $\mathbf A$.
\item Find all $\mathbf v\in\R^3$ such that $\mathbf v-(3,4,7)$ lies in
\[
\Span(\mathbf v_1)^\perp\cap\Span(\mathbf v_2)^\perp\cap\Span(\mathbf v_3)^\perp.
\]
\end{problemsubparts}
\end{problem}

% 2024 Spring KAIST Introduction to Linear Algebra Midterm, Problem 6
\begin{problem}[2024 Spring \KAIST{} Introduction to Linear Algebra]
Suppose $\mathbf A,\mathbf B,\mathbf C,\mathbf D\in\R^{3\times3}$ satisfy
\[
\mathbf A\mathbf D=
\begin{pmatrix}
-1&-5&-8\\
-1&-5&-8\\
-2&-10&75
\end{pmatrix},
\qquad
\mathbf B\mathbf D=
\begin{pmatrix}
0&0&0\\
4&20&33\\
0&0&-46
\end{pmatrix},
\]
and
\[
(\mathbf A+\mathbf B)\mathbf C=\mathbf I_3.
\]
Find a matrix in reduced row echelon form that is row equivalent to $\mathbf D$.
\end{problem}

% 2024 Spring KAIST Introduction to Linear Algebra Midterm, Problem 7
\begin{problem}[2024 Spring \KAIST{} Introduction to Linear Algebra]
Let $\mathbf A\in\R^{3\times3}$ be orthogonal. Determine whether each statement is true or false, and justify your answer by proof or counterexample.
\begin{problemchoices}
\item $\mathbf A\mathbf B$ is orthogonal, where
\[
\mathbf B=
\begin{pmatrix}
0&0&1\\
0&1&0\\
1&0&0
\end{pmatrix}.
\]
\item $|\lambda|=1$ for every eigenvalue $\lambda$ of $\mathbf A$.
\end{problemchoices}
\end{problem}

% 2024 Spring KAIST Introduction to Linear Algebra Midterm, Problem 8
\begin{problem}[2024 Spring \KAIST{} Introduction to Linear Algebra]
On $\R^2$, let $T_1$ be reflection across the $x$-axis and let $T_2$ be reflection across the line $y=x$.
\begin{problemsubparts}
\item Show, using their standard matrices, that $T_1\circ T_2\ne T_2\circ T_1$.
\item Describe $T_1\circ T_2$ and $T_2\circ T_1$ as rotations about the origin.
\end{problemsubparts}
\end{problem}

% 2024 Spring KAIST Introduction to Linear Algebra Midterm, Problem 9
\begin{problem}[2024 Spring \KAIST{} Introduction to Linear Algebra]
Let $\mathbf A\in\R^{2\times2}$ be symmetric. Determine whether each statement is true or false, and justify your answer by proof or counterexample.
\begin{problemchoices}
\item $\mathbf A$ has a real eigenvalue.
\item If $\mathbf A$ has distinct eigenvalues $\lambda_1,\lambda_2$ and $\mathbf x_i$ is an eigenvector corresponding to $\lambda_i$ for $i=1,2$, then $\mathbf x_1\perp\mathbf x_2$.
\end{problemchoices}
\end{problem}

% 2024 Spring KAIST Introduction to Linear Algebra Midterm, Problem 10
\begin{problem}[2024 Spring \KAIST{} Introduction to Linear Algebra]
Let $\mathbf A,\mathbf B\in\R^{n\times n}$ and suppose that an invertible matrix $\mathbf Q\in\R^{n\times n}$ satisfies
\[
\mathbf A=\mathbf Q^{-1}\mathbf B\mathbf Q.
\]
Determine whether each statement is true or false, and justify your answer by proof or counterexample.
\begin{problemchoices}
\item $\tr(\mathbf A)=\tr(\mathbf B)$.
\item If $\lambda$ is an eigenvalue of $\mathbf A$, then $\lambda$ is an eigenvalue of $\mathbf B$.
\item If $\mathbf x$ is an eigenvector of $\mathbf A$, then $\mathbf x$ is an eigenvector of $\mathbf B$.
\end{problemchoices}
\end{problem}

% 2024 Spring KAIST Introduction to Linear Algebra Midterm, Problem 11
\begin{problem}[2024 Spring \KAIST{} Introduction to Linear Algebra]
Let
\[
\mathbf A=
\begin{pmatrix}
-2&0&0&0\\
1&-2&0&0\\
-7&4&-1&0\\
1&-8&13&1
\end{pmatrix},
\qquad
\mathbf B=
\begin{pmatrix}
1&6&-2&7\\
0&-1&11&5\\
0&0&1&-4\\
0&0&0&-2
\end{pmatrix}.
\]
Find $\det\bigl(\adj(\mathbf A\mathbf B)\bigr)$.
\end{problem}

% 2024 Spring KAIST Introduction to Linear Algebra Midterm, Problem 12
% Source correction: the printed statement in part (a) asks for an invertible 3x3 skew-symmetric matrix, which cannot exist over R. The evident missing word "no" is restored.
\begin{problem}[2024 Spring \KAIST{} Introduction to Linear Algebra]
\begin{problemsubparts}
\item Show that there is no invertible skew-symmetric $3\times3$ real matrix.
\item Suppose $\mathbf A\in\R^{n\times n}$ is skew-symmetric and $\mathbf I_n$ is the identity matrix. Show that $\mathbf I_n-\mathbf A$ is invertible.
\end{problemsubparts}
\end{problem}

% 2024 Spring KAIST Introduction to Linear Algebra Final, Problem 1
\begin{problem}[2024 Spring \KAIST{} Introduction to Linear Algebra]
Let
\[
\begin{aligned}
\mathbf v_1&=(1,1,-1),&
\mathbf v_2&=(2,2,-2),&
\mathbf v_3&=(2,2,2),\\
\mathbf v_4&=(2,2,-6),&
\mathbf v_5&=(3,3,5),&
\mathbf v_6&=(3,4,2).
\end{aligned}
\]
Find all $i\in\{2,3,4,5,6\}$ such that $\mathbf v_i$ can be expressed as a linear combination of $\mathbf v_1,\dots,\mathbf v_{i-1}$. Justify your answer.
\end{problem}

% 2024 Spring KAIST Introduction to Linear Algebra Final, Problem 2
\begin{problem}[2024 Spring \KAIST{} Introduction to Linear Algebra]
Let $\mathbf A\in\R^{2\times2}$ have characteristic polynomial
\[
p(x)=x^2-x+1.
\]
Compute
\[
\tr\left(\mathbf A^{109}+\bigl(\mathbf A^{-1}\bigr)^{109}\right).
\]
\end{problem}

% 2024 Spring KAIST Introduction to Linear Algebra Final, Problem 3
\begin{problem}[2024 Spring \KAIST{} Introduction to Linear Algebra]
Let
\[
\mathbf A=
\begin{pmatrix}
1&2&3\\
1&2&4\\
3&6&1
\end{pmatrix}.
\]
\begin{problemsubparts}
\item Find a basis for $\Col(\mathbf A)$.
\item Write
\[
\mathbf A=\mathbf c_1\mathbf r_1+\mathbf c_2\mathbf r_2+\cdots+\mathbf c_k\mathbf r_k,
\]
where $k=\rk(\mathbf A)$, each $\mathbf c_i$ is a column of $\mathbf A$, and each $\mathbf r_i$ is a $1\times3$ row vector.
\end{problemsubparts}
\end{problem}

% 2024 Spring KAIST Introduction to Linear Algebra Final, Problem 4
\begin{problem}[2024 Spring \KAIST{} Introduction to Linear Algebra]
Let
\[
\mathbf A=
\begin{pmatrix}
0&1&0\\
1&0&0\\
0&1&1
\end{pmatrix},
\qquad
\mathbf b=
\begin{pmatrix}1\\1\\1\end{pmatrix}.
\]
\begin{problemsubparts}
\item Find a $QR$ decomposition of $\mathbf A$.
\item Find the least-squares solution of $\mathbf A\mathbf x=\mathbf b$.
\end{problemsubparts}
\end{problem}

% 2024 Spring KAIST Introduction to Linear Algebra Final, Problem 5
\begin{problem}[2024 Spring \KAIST{} Introduction to Linear Algebra]
Let
\[
W=\Span\bigl((1,-1,1,0),(0,1,1,-1)\bigr)\subseteq\R^4,
\qquad
\mathbf b=
\begin{pmatrix}1\\0\\0\\1\end{pmatrix}.
\]
Find $\widehat{\mathbf w}\in W$ such that
\[
\|\mathbf b-\widehat{\mathbf w}\|\le\|\mathbf b-\mathbf w\|
\]
for every $\mathbf w\in W$.
\end{problem}

% 2024 Spring KAIST Introduction to Linear Algebra Final, Problem 6
\begin{problem}[2024 Spring \KAIST{} Introduction to Linear Algebra]
Find the least-squares line of best fit to the four points
\[
(0,-1),\quad(-1,-1),\quad(1,1),\quad(2,2).
\]
\end{problem}

% 2024 Spring KAIST Introduction to Linear Algebra Final, Problem 7
\begin{problem}[2024 Spring \KAIST{} Introduction to Linear Algebra]
Find a singular value decomposition of
\[
\mathbf A=
\begin{pmatrix}
1&-2&0\\
0&-2&1
\end{pmatrix}.
\]
\end{problem}

% 2024 Spring KAIST Introduction to Linear Algebra Final, Problem 8
\begin{problem}[2024 Spring \KAIST{} Introduction to Linear Algebra]
Let the ordered bases $\mathcal B_1$ and $\mathcal B_2$ of $\R^3$ be
\[
\mathcal B_1=
\left(
\begin{pmatrix}1\\0\\0\end{pmatrix},
\begin{pmatrix}1\\0\\-1\end{pmatrix},
\begin{pmatrix}0\\1\\1\end{pmatrix}
\right),
\]
\[
\mathcal B_2=
\left(
\begin{pmatrix}0\\1\\1\end{pmatrix},
\begin{pmatrix}-1\\3\\2\end{pmatrix},
\begin{pmatrix}2\\-3\\-2\end{pmatrix}
\right).
\]
Suppose $\mathcal B_3$ is an ordered basis of $\R^3$ satisfying
\[
\mathbf P_{\mathcal B_2\to\mathcal B_3}=
\begin{pmatrix}
1&1&0\\
0&-1&0\\
2&0&1
\end{pmatrix}.
\]
\begin{problemsubparts}
\item Find $\mathbf P_{\mathcal B_3\to\mathcal B_1}$.
\item Let $T\colon\R^4\to\R^3$ be defined by
\[
T(a,b,c,d)=(a+b,c-a,b+d).
\]
Find $[T]_{\mathcal B_1,\mathcal E}$, where $\mathcal E$ is the standard ordered basis of $\R^4$.
\end{problemsubparts}
\end{problem}

% 2024 Spring KAIST Introduction to Linear Algebra Final, Problem 9
\begin{problem}[2024 Spring \KAIST{} Introduction to Linear Algebra]
Let
\[
\mathbf A=
\begin{pmatrix}
1&0&y\\
0&1&y\\
2&2&y-3
\end{pmatrix},
\qquad y\in\R.
\]
\begin{problemsubparts}
\item Show that $1$ is an eigenvalue of $\mathbf A$ for every $y\in\R$.
\item Find all $y\in\R$ for which $\mathbf A$ is not diagonalizable. Justify your answer.
\item Find all $y\in\R$ for which $\mathbf A$ is orthogonally diagonalizable. Justify your answer and, in each case, orthogonally diagonalize $\mathbf A$.
\end{problemsubparts}
\end{problem}

% 2024 Spring KAIST Introduction to Linear Algebra Final, Problem 10
\begin{problem}[2024 Spring \KAIST{} Introduction to Linear Algebra]
Suppose $\mathbf B\in\R^{4\times4}$ satisfies
\[
\mathbf A\mathbf B=
\begin{pmatrix}
1&2&-1&2\\
3&6&-2&5\\
0&0&0&0\\
0&0&0&0
\end{pmatrix},
\qquad
\mathbf A=
\begin{pmatrix}
2&-1&1&1\\
3&2&-1&1\\
-1&3&2&-1\\
1&-1&3&2
\end{pmatrix}.
\]
Find a basis for $\Row(\mathbf B)$ and a basis for $\Null(\mathbf B^{\mathsf T})$. You may use $\det(\mathbf A)=45$.
\end{problem}

% 2024 Spring KAIST Introduction to Linear Algebra Final, Problem 11
\begin{problem}[2024 Spring \KAIST{} Introduction to Linear Algebra]
For
\[
\mathbf A=
\begin{pmatrix}
3&4&5&-2\\
6&8&10&-4\\
1&0&2&2\\
2&4&3&-4
\end{pmatrix},
\]
find matrices $\mathbf D$ and $\mathbf S$ such that
\[
\mathbf A=\mathbf S\mathbf D,
\]
the rows of $\mathbf D$ form a linearly independent subset of the rows of $\mathbf A$, and $\mathbf S^{\mathsf T}$ is in reduced row echelon form.
\end{problem}

% 2024 Spring KAIST Introduction to Linear Algebra Final, Problem 12
\begin{problem}[2024 Spring \KAIST{} Introduction to Linear Algebra]
Let $\mathbf P\in\R^{n\times n}$ satisfy $\mathbf P^2=\mathbf P$. Determine whether each statement is true or false, and justify your answer by proof or counterexample.
\begin{problemsubparts}
\item $\mathbf P$ is diagonalizable.
\item $\tr(\mathbf P)=\rk(\mathbf P)$.
\item If $\mathbf P$ is symmetric, then $\mathbf P$ is the standard matrix of the orthogonal projection of $\R^n$ onto $\Col(\mathbf P)$.
\end{problemsubparts}
\end{problem}

% 2024 Spring KAIST Introduction to Linear Algebra Final, Problem 13
\begin{problem}[2024 Spring \KAIST{} Introduction to Linear Algebra]
Let $\mathbf A\in\R^{m\times n}$ and $\mathbf B\in\R^{n\times k}$, and let $\mathbf R$ and $\mathbf S$ be their reduced row echelon forms, respectively. Determine whether each statement is true or false, and justify your answer by proof or counterexample.
\begin{problemsubparts}
\item $\rk(\mathbf A\mathbf B)=\rk(\mathbf R\mathbf B)$.
\item $\rk(\mathbf A\mathbf B)=\rk(\mathbf A\mathbf S)$.
\end{problemsubparts}
\end{problem}

% 2024 Spring KAIST Introduction to Linear Algebra Final, Problem 14
\begin{problem}[2024 Spring \KAIST{} Introduction to Linear Algebra]
Let $\mathbf A\in\R^{4\times5}$, $\mathbf v_1\in\R^5$, and $\mathbf v_2\in\R^4$. Suppose there exist linearly independent vectors $\mathbf x_1,\mathbf x_2,\mathbf x_3,\mathbf x_4\in\R^4$ such that
\[
\mathbf A\mathbf v_1=\mathbf x_1,
\qquad
\mathbf A\mathbf A^{\mathsf T}\mathbf v_2=\mathbf x_2,
\qquad
\mathbf A^{\mathsf T}\mathbf x_3=\mathbf 0,
\qquad
\mathbf A\mathbf A^{\mathsf T}\mathbf x_4=\mathbf 0.
\]
Compute $\rk(\mathbf A)$ and justify your answer.
\end{problem}

% 2024 Fall KAIST Introduction to Linear Algebra Midterm, Problem 1
\begin{problem}[2024 Fall \KAIST{} Introduction to Linear Algebra]
Find all $a,b,c\in\R$ for which
\[
\mathbf A=
\begin{pmatrix}
1&a-b+2c&2a-b+c\\
3&-1&a+c\\
0&-2&1
\end{pmatrix}
\]
is symmetric.
\end{problem}

% 2024 Fall KAIST Introduction to Linear Algebra Midterm, Problem 2
\begin{problem}[2024 Fall \KAIST{} Introduction to Linear Algebra]
Determine whether each function is a linear transformation. If it is, find its standard matrix.
\begin{problemsubparts}
\item $T_1\colon\R^3\to\R^3$ is reflection about the origin $(0,0,0)$.
\item $T_2\colon\R^3\to\R^3$ is orthogonal projection onto the plane
\[
2x+3y+4z=1.
\]
\end{problemsubparts}
\end{problem}

% 2024 Fall KAIST Introduction to Linear Algebra Midterm, Problem 3
\begin{problem}[2024 Fall \KAIST{} Introduction to Linear Algebra]
Let
\[
\mathbf v_1=(1,4,4),\qquad
\mathbf v_2=(-1,3,6),\qquad
\mathbf v_3=(-7,0,12)
\]
in $\R^3$.
\begin{problemsubparts}
\item Determine whether $\mathbf v_1,\mathbf v_2,\mathbf v_3$ are linearly independent. If they are, prove it. Otherwise, express one vector as a linear combination of the other two.
\item Find all $a\in\R$ for which $\mathbf w=(3,5,a)\in\Span(\mathbf v_1,\mathbf v_2,\mathbf v_3)$.
\end{problemsubparts}
\end{problem}

% 2024 Fall KAIST Introduction to Linear Algebra Midterm, Problem 4
\begin{problem}[2024 Fall \KAIST{} Introduction to Linear Algebra]
Let $\mathbf A$ be invertible.
\begin{problemsubparts}
\item Show that $\mathbf A^{\mathsf T}$ is invertible and
\[
\bigl(\mathbf A^{\mathsf T}\bigr)^{-1}=\bigl(\mathbf A^{-1}\bigr)^{\mathsf T}.
\]
\item Show that if $\mathbf A$ is skew-symmetric, that is, $\mathbf A^{\mathsf T}=-\mathbf A$, then $\mathbf A^{-1}$ is also skew-symmetric.
\end{problemsubparts}
\end{problem}

% 2024 Fall KAIST Introduction to Linear Algebra Midterm, Problem 5
\begin{problem}[2024 Fall \KAIST{} Introduction to Linear Algebra]
Find all $b\in\R$ for which the system
\[
\begin{cases}
x_1-2x_2-x_3=b^2,\\
-2x_1+3x_2+2x_3=b,\\
-4x_1+7x_2+4x_3=-1
\end{cases}
\]
is consistent. For each such $b$, find the general solution.
\end{problem}

% 2024 Fall KAIST Introduction to Linear Algebra Midterm, Problem 6
\begin{problem}[2024 Fall \KAIST{} Introduction to Linear Algebra]
Let
\[
\mathbf A=\frac13
\begin{pmatrix}
1&2&-2\\
2&1&2\\
-2&2&1
\end{pmatrix}.
\]
\begin{problemsubparts}
\item Determine whether $\mathbf A$ is orthogonal. Justify your answer.
\item Find all real eigenvalues of $\mathbf A$ and the eigenspace corresponding to each one.
\end{problemsubparts}
\end{problem}

% 2024 Fall KAIST Introduction to Linear Algebra Midterm, Problem 7
\begin{problem}[2024 Fall \KAIST{} Introduction to Linear Algebra]
Let
\[
\mathbf A=
\begin{pmatrix}
1&1&2\\
6&7&-9\\
0&3&-11
\end{pmatrix},
\qquad
\mathbf b=
\begin{pmatrix}0\\1\\107\end{pmatrix}.
\]
Find an $LU$ decomposition of $\mathbf A$, and use it to solve $\mathbf A\mathbf x=\mathbf b$ by forward substitution followed by backward substitution. Do not use inverses of the factors obtained from the decomposition.
\end{problem}

% 2024 Fall KAIST Introduction to Linear Algebra Midterm, Problem 8
\begin{problem}[2024 Fall \KAIST{} Introduction to Linear Algebra]
Let $\mathbf A,\mathbf B,\mathbf C,\mathbf D$ be matrices such that $\mathbf A\mathbf B$, $\mathbf A\mathbf C$, $\mathbf B\mathbf D$, and $\mathbf C\mathbf D$ are defined. Determine whether each statement is true or false.
\begin{problemsubparts}
\item If one column of $\mathbf B$ equals one column of $\mathbf C$, then one column of $\mathbf A\mathbf B$ equals one column of $\mathbf A\mathbf C$.
\item If $\mathbf B$ and $\mathbf C$ are row equivalent, then $\mathbf A\mathbf B$ and $\mathbf A\mathbf C$ are row equivalent.
\item If $\mathbf B$ and $\mathbf C$ are row equivalent, then $\mathbf B\mathbf D$ and $\mathbf C\mathbf D$ are row equivalent.
\end{problemsubparts}
\end{problem}

% 2024 Fall KAIST Introduction to Linear Algebra Midterm, Problem 9
\begin{problem}[2024 Fall \KAIST{} Introduction to Linear Algebra]
Let $T_1,T_2\colon\R^n\to\R^m$ be linear transformations and let $c_1,c_2\in\R$. Define
\[
(c_1T_1+c_2T_2)(\mathbf u)=c_1T_1(\mathbf u)+c_2T_2(\mathbf u),
\qquad \mathbf u\in\R^n.
\]
\begin{problemsubparts}
\item Show that $c_1T_1+c_2T_2$ is a linear transformation.
\item Suppose now that $m=n$. A scalar $\lambda\in\R$ is called an eigenvalue of a linear operator $S\colon\R^n\to\R^n$ if there exists a nonzero vector $\mathbf x\in\R^n$ such that $S(\mathbf x)=\lambda\mathbf x$. Determine whether the following statement is true or false: if $\lambda_i$ is an eigenvalue of $T_i$ for $i=1,2$, then $c_1\lambda_1+c_2\lambda_2$ is an eigenvalue of $c_1T_1+c_2T_2$.
\end{problemsubparts}
\end{problem}

% 2024 Fall KAIST Introduction to Linear Algebra Midterm, Problem 10
\begin{problem}[2024 Fall \KAIST{} Introduction to Linear Algebra]
For an integer $n>2$, define $\mathbf A=(a_{ij})\in\R^{n\times n}$ by
\[
a_{ij}=
\begin{cases}
r,&i=j,\\
1,&i+1=j\text{ and }i\ne2,\\
0,&\text{otherwise},
\end{cases}
\]
where $r\in\R$ is fixed.
\begin{problemsubparts}
\item Find the characteristic polynomial of $\mathbf A$.
\item Find all real eigenvalues of $\mathbf A$.
\item Find the eigenspace $E_\lambda(\mathbf A)$ corresponding to each real eigenvalue $\lambda$.
\end{problemsubparts}
\end{problem}

% 2024 Fall KAIST Introduction to Linear Algebra Midterm, Problem 11
\begin{problem}[2024 Fall \KAIST{} Introduction to Linear Algebra]
Let
\[
\mathbf u=(\sqrt2,\sqrt2,0),
\qquad
\mathbf v=(0,2,0).
\]
For $\theta\ge0$, let $T_\theta\colon\R^3\to\R^3$ be rotation about the positive $x$-axis through angle $\theta$.
\begin{problemsubparts}
\item Find $\theta\in[0,\pi]$ for which the area of the parallelogram spanned by $\mathbf u$ and $T_\theta(\mathbf v)$ is maximized.
\item For the angle found in part \textup{(a)}, find the standard matrix of $T_\theta$.
\item For the angle found in part \textup{(a)}, find a vector perpendicular to the parallelogram spanned by $\mathbf u$ and $T_\theta(\mathbf v)$.
\end{problemsubparts}
\end{problem}

% 2024 Fall KAIST Introduction to Linear Algebra Midterm, Problem 12
\begin{problem}[2024 Fall \KAIST{} Introduction to Linear Algebra]
Let $\mathbf A=(a_{ij})\in\R^{n\times n}$ and let $i,j\in\{1,\dots,n\}$.
\begin{problemsubparts}
\item State the definition of the $(i,j)$-cofactor of $\mathbf A$ and of the adjugate matrix $\adj(\mathbf A)$.
\item State the cofactor expansions of $\det(\mathbf A)$ along the $i$th row and the $j$th column.
\item Compute $\mathbf A\adj(\mathbf A)$ and justify your answer using suitable cofactor expansions.
\end{problemsubparts}
\end{problem}

% 2024 Fall KAIST Introduction to Linear Algebra Midterm, Problem 13
\begin{problem}[2024 Fall \KAIST{} Introduction to Linear Algebra]
Let $n>1$, let $\mathbf A\in\R^{n\times n}$, and let $\mathbf I_n$ be the identity matrix.
\begin{problemsubparts}
\item Show that if $\mathbf A$ is nilpotent, then $\mathbf I_n-\mathbf A$ is invertible.
\item For the matrix $\mathbf A=(a_{ij})$ defined by
\[
a_{ij}=
\begin{cases}
1,&i+1=j,\\
0,&\text{otherwise},
\end{cases}
\]
find the index of nilpotency of $\mathbf A$, that is, the smallest positive integer $k$ such that $\mathbf A^k=\mathbf0$, and find $(\mathbf I_n-\mathbf A)^{-1}$. No proof is required for this part.
\item For the matrix in part \textup{(b)}, find all fixed points of $\mathbf A$, that is, all $\mathbf x\in\R^n$ satisfying $\mathbf A\mathbf x=\mathbf x$.
\end{problemsubparts}
\end{problem}

% 2024 Fall KAIST Introduction to Linear Algebra Midterm, Problem 14
\begin{problem}[2024 Fall \KAIST{} Introduction to Linear Algebra]
Suppose $\mathbf u_1,\dots,\mathbf u_4$ span $\R^4$, and suppose there exists a linear transformation $T\colon\R^4\to\R^6$ such that
\[
\begin{aligned}
T(\mathbf u_1)&=(a_1,b_1,0,0,0,0),\\
T(\mathbf u_2)&=(0,a_2,b_2,0,0,0),\\
T(\mathbf u_3)&=(0,0,a_3,b_3,0,0),\\
T(\mathbf u_4)&=(b_4,0,0,a_4,0,0),
\end{aligned}
\]
where $a_1,a_2,a_3,a_4$ are even integers and $b_1,b_2,b_3,b_4$ are odd integers.
\begin{problemsubparts}
\item Find $\ker T$.
\item Find $\im T$.
\item Find $c_1,c_2,c_3,c_4\in\R$ such that
\[
T\left(\sum_{i=1}^4 c_i\mathbf u_i\right)=(1,0,0,0,0,0).
\]
\textit{Hint.} Each $c_i$ can be expressed as a rational function of
$a_1,a_2,a_3,a_4,b_1,b_2,b_3,b_4$.
\end{problemsubparts}
\end{problem}

% 2024 Fall KAIST Introduction to Linear Algebra Midterm, Problem 15
\begin{problem}[2024 Fall \KAIST{} Introduction to Linear Algebra]
\begin{problemsubparts}
\item For a positive odd integer $m$, compute $\det(\mathbf A)$ for an arbitrary $m\times m$ skew-symmetric matrix $\mathbf A$. Your answer should not depend on the specific matrix.
\item For a positive integer $n$, let $\mathbf B\in\R^{n\times n}$ be orthogonal with $\det(\mathbf B)<0$. Compute
\[
\det(\mathbf B+\mathbf I_n).
\]
Your answer should not depend on the specific matrix $\mathbf B$.
\end{problemsubparts}
\end{problem}

% 2024 Fall KAIST Introduction to Linear Algebra Final, Problem 1
\begin{problem}[2024 Fall \KAIST{} Introduction to Linear Algebra]
\begin{problemsubparts}
\item Find the standard matrix $\mathbf P$ for the orthogonal projection of $\R^3$ onto the plane
\[
x-5y+10z=0.
\]
\item Use $\mathbf P$ to find the orthogonal projection of $(5,2,1)\in\R^3$ onto the plane.
\end{problemsubparts}
\end{problem}

% 2024 Fall KAIST Introduction to Linear Algebra Final, Problem 2
\begin{problem}[2024 Fall \KAIST{} Introduction to Linear Algebra]
Let
\[
W=\Span\bigl((1,1,0),(1,2,-1),(0,1,1)\bigr)\subseteq\R^3.
\]
\begin{problemsubparts}
\item Find an orthonormal basis for $W$.
\item Find a $QR$ decomposition of
\[
\mathbf A=
\begin{pmatrix}
1&1\\
1&2\\
0&-1
\end{pmatrix}.
\]
\end{problemsubparts}
\end{problem}

% 2024 Fall KAIST Introduction to Linear Algebra Final, Problem 3
\begin{problem}[2024 Fall \KAIST{} Introduction to Linear Algebra]
Let $\mathbf A,\mathbf B\in\R^{m\times n}$ be row equivalent, and let $\mathbf a_1,\dots,\mathbf a_n$ and $\mathbf b_1,\dots,\mathbf b_n$ be their column vectors, respectively. Prove that for any $c_1,\dots,c_n\in\R$,
\[
\sum_{i=1}^n c_i\mathbf a_i=\mathbf0
\quad\Longleftrightarrow\quad
\sum_{i=1}^n c_i\mathbf b_i=\mathbf0.
\]
\end{problem}

% 2024 Fall KAIST Introduction to Linear Algebra Final, Problem 4
\begin{problem}[2024 Fall \KAIST{} Introduction to Linear Algebra]
Let $\mathcal B$ be any ordered basis of $\R^n$. Prove that vectors $\mathbf v_1,\dots,\mathbf v_k\in\R^n$ are linearly independent if and only if their coordinate vectors
\[
[\mathbf v_1]_{\mathcal B},\dots,[\mathbf v_k]_{\mathcal B}
\]
are linearly independent in $\R^n$.
\end{problem}

% 2024 Fall KAIST Introduction to Linear Algebra Final, Problem 5
\begin{problem}[2024 Fall \KAIST{} Introduction to Linear Algebra]
Let
\[
\mathbf P=
\begin{pmatrix}
2&1&0\\
0&2&1\\
1&0&2
\end{pmatrix}.
\]
Let $\mathcal E$ denote the standard ordered basis of $\R^3$.
\begin{problemsubparts}
\item Find the ordered basis $\mathcal B$ of $\R^3$ such that $\mathbf P$ is the transition matrix from $\mathcal B$ to $\mathcal E$.
\item Find the ordered basis $\mathcal C$ of $\R^3$ such that $\mathbf P$ is the transition matrix from $\mathcal E$ to $\mathcal C$.
\end{problemsubparts}
\end{problem}

% 2024 Fall KAIST Introduction to Linear Algebra Final, Problem 6
\begin{problem}[2024 Fall \KAIST{} Introduction to Linear Algebra]
Let
\[
W=\Span(\mathbf v_1,\mathbf v_2,\mathbf v_3,\mathbf v_4,\mathbf v_5)\subseteq\R^4,
\]
where
\[
\begin{aligned}
\mathbf v_1&=(2,7,3,1),&
\mathbf v_2&=(1,0,4,3),&
\mathbf v_3&=(0,2,3,0),\\
\mathbf v_4&=(3,6,1,2),&
\mathbf v_5&=(1,0,0,9).
\end{aligned}
\]
\begin{problemsubparts}
\item Find a subset of $\{\mathbf v_1,\mathbf v_2,\mathbf v_3,\mathbf v_4,\mathbf v_5\}$ that forms a basis for $W$.
\item Express each vector not chosen for the basis as a linear combination of the basis vectors.
\end{problemsubparts}
\end{problem}

% 2024 Fall KAIST Introduction to Linear Algebra Final, Problem 7
\begin{problem}[2024 Fall \KAIST{} Introduction to Linear Algebra]
Let
\[
\mathbf A=
\begin{pmatrix}
-1&-3\\
1&3\\
2&6
\end{pmatrix},
\qquad
\mathbf b=
\begin{pmatrix}1\\0\\1\end{pmatrix}.
\]
\begin{problemsubparts}
\item Find all least-squares solutions of $\mathbf A\mathbf x=\mathbf b$.
\item Find the least-squares solution with minimum norm $\|\mathbf x\|$.
\end{problemsubparts}
\end{problem}

% 2024 Fall KAIST Introduction to Linear Algebra Final, Problem 8
\begin{problem}[2024 Fall \KAIST{} Introduction to Linear Algebra]
\begin{problemsubparts}
\item Let $\mathbf A\in\R^{n\times n}$ have singular value decomposition
\[
\mathbf A=\mathbf U\boldsymbol\Sigma\mathbf V^{\mathsf T}.
\]
For $i=1,\dots,n$, let $\mathbf u_i$ and $\mathbf v_i$ be the $i$th columns of $\mathbf U$ and $\mathbf V$, respectively, and let $\sigma_i$ be the $(i,i)$-entry of $\boldsymbol\Sigma$. Find an orthonormal basis for $\R^n$ consisting of eigenvectors of $\mathbf A^{\mathsf T}\mathbf A$, and an orthonormal basis for $\R^n$ consisting of eigenvectors of $\mathbf A\mathbf A^{\mathsf T}$.
\item Find a singular value decomposition of
\[
\mathbf B=
\begin{pmatrix}
1&0&0\\
1&1&0\\
-1&1&0
\end{pmatrix}.
\]
\end{problemsubparts}
\end{problem}

% 2024 Fall KAIST Introduction to Linear Algebra Final, Problem 9
% Source correction: the printed problem has b in R^4, but A^T A x lies in R^5 and x is stated to be in R^5. The ambient space is therefore minimally corrected to R^5.
\begin{problem}[2024 Fall \KAIST{} Introduction to Linear Algebra]
Let
\[
\mathbf A=
\begin{pmatrix}
3&-5&-1&5&4\\
0&-3&2&7&6\\
4&-5&-9&2&-5\\
2&6&8&3&-6
\end{pmatrix}.
\]
\begin{problemsubparts}
\item Show that
\[
W=\left\{\mathbf b\in\R^5\mid \mathbf A^{\mathsf T}\mathbf A\mathbf x=\mathbf b
\text{ has at least one solution }\mathbf x\in\R^5\right\}
\]
is a subspace of $\R^5$. \textit{Hint.} Show that $W$ is one of the fundamental spaces of a matrix.
\item Find $\dim W$ and $\dim W^\perp$.
\end{problemsubparts}
\end{problem}

% 2024 Fall KAIST Introduction to Linear Algebra Final, Problem 10
\begin{problem}[2024 Fall \KAIST{} Introduction to Linear Algebra]
Let $W\subseteq\R^n$ be an $m$-dimensional subspace with $0<m<n$, and let $T\colon\R^n\to\R^n$ be the orthogonal projection onto $W$. Show that there exists an ordered basis $\mathcal B$ of $\R^n$ such that
\[
[T]_{\mathcal B}=
\begin{pmatrix}
\mathbf I_m&\mathbf A\\
\mathbf0&\mathbf0
\end{pmatrix}
\]
for some $m\times(n-m)$ matrix $\mathbf A$.
\end{problem}

% 2024 Fall KAIST Introduction to Linear Algebra Final, Problem 11
\begin{problem}[2024 Fall \KAIST{} Introduction to Linear Algebra]
Consider the proposition: if square matrices $\mathbf A$ and $\mathbf B$ are similar, then their characteristic polynomials are equal.
\begin{problemsubparts}
\item Prove the proposition.
\item Prove or disprove its converse.
\item Prove or disprove the following statement: if $\mathbf A$ and $\mathbf B$ are diagonalizable square matrices and have the same characteristic polynomial, then they are similar.
\end{problemsubparts}
\textit{Hint.} You may use without proof that diagonal matrices with the same characteristic polynomial are similar.
\end{problem}

% 2024 Fall KAIST Introduction to Linear Algebra Final, Problem 12
\begin{problem}[2024 Fall \KAIST{} Introduction to Linear Algebra]
Find vectors $\mathbf v_1,\mathbf v_2,\mathbf v_3,\mathbf v_4\in\R^3$ and $\mathbf w_1,\mathbf w_2,\mathbf w_3\in\R^4$ satisfying all of the following conditions.
\begin{problemclauses}
\item The entries of $\mathbf v_1,\mathbf v_2,\mathbf v_3,\mathbf v_4$ are exactly the twelve integers $1,2,\dots,12$, each used once.
\item The entries of $\mathbf w_1,\mathbf w_2,\mathbf w_3$ are exactly the twelve integers $1,2,\dots,12$, each used once.
\item For every $i\in\{1,2,3,4\}$ and $j\in\{1,2,3\}$, the vectors $\mathbf v_i$ and $\mathbf w_j$ share exactly one component.
\item The subspaces $\Span(\mathbf v_1,\mathbf v_2,\mathbf v_3,\mathbf v_4)\subseteq\R^3$ and $\Span(\mathbf w_1,\mathbf w_2,\mathbf w_3)\subseteq\R^4$ have the same dimension.
\end{problemclauses}
\end{problem}

% 2024 Fall KAIST Introduction to Linear Algebra Final, Problem 13
\begin{problem}[2024 Fall \KAIST{} Introduction to Linear Algebra]
Let $n$ be a positive integer and let $\mathbf I_n$ be the identity matrix.
\begin{problemsubparts}
\item Suppose $\mathbf A\in\R^{n\times n}$ satisfies $\mathbf A^2=\mathbf I_n$. Prove that:
\begin{problemclauses}
\item the eigenvalues of $\mathbf A$ are contained in $\{1,-1\}$;
\item $\Col(\mathbf A-\mathbf I_n)\subseteq\Null(\mathbf A+\mathbf I_n)$;
\item $\mathbf A$ is diagonalizable.
\end{problemclauses}
\item Suppose $\mathbf B\in\R^{n\times n}$ is orthogonal. Show that if $\mathbf B$ is diagonalizable over $\R$, then
\[
\mathbf B^2=\mathbf I_n.
\]
\end{problemsubparts}
\end{problem}

% 2025 Spring KAIST Introduction to Linear Algebra Midterm, Problem 1
\begin{problem}[2025 Spring \KAIST{} Introduction to Linear Algebra]
Consider a system of $m$ linear equations in $n$ unknowns.
\begin{problemsubparts}
\item If the system is homogeneous, explain why it has at least one solution.
\item If $m<n$, systems often have solutions. Give an example with $m=2$ and $n=3$ of the form
\[
\begin{cases}
a_{11}x+a_{12}y+a_{13}z=b_1,\\
a_{21}x+a_{22}y+a_{23}z=b_2,
\end{cases}
\]
in which the six coefficients $a_{11},a_{12},a_{13},a_{21},a_{22},a_{23}$ are pairwise distinct and the system has no solution. You need not solve the system.
\item If $m>n$, systems often have no solutions. Give an example with $m=3$ and $n=2$ of the form
\[
\begin{cases}
a_{11}x+a_{12}y=b_1,\\
a_{21}x+a_{22}y=b_2,\\
a_{31}x+a_{32}y=b_3,
\end{cases}
\]
in which the six coefficients $a_{11},a_{12},a_{21},a_{22},a_{31},a_{32}$ are pairwise distinct and the system has a unique solution. You need not solve the system.
\end{problemsubparts}
\end{problem}

% 2025 Spring KAIST Introduction to Linear Algebra Midterm, Problem 2
\begin{problem}[2025 Spring \KAIST{} Introduction to Linear Algebra]
Consider the system
\[
\begin{cases}
x+2y+3z=14,\\
2x-y+z=3,\\
3x+4y-z=8.
\end{cases}
\]
\begin{problemsubparts}
\item Write the augmented matrix and use elementary row operations to reduce it to reduced row echelon form.
\item Determine whether the system has a unique solution, infinitely many solutions, or no solution. Justify your answer and, if solutions exist, find all of them.
\end{problemsubparts}
\end{problem}

% 2025 Spring KAIST Introduction to Linear Algebra Midterm, Problem 3
\begin{problem}[2025 Spring \KAIST{} Introduction to Linear Algebra]
Let
\[
\mathbf A=
\begin{pmatrix}
1&2&3\\
4&5&6\\
7&8&9
\end{pmatrix}.
\]
\begin{problemsubparts}
\item Compute $\det(\mathbf A)$ and list all possible cases for the number of solutions of $\mathbf A\mathbf x=\mathbf b$ as $\mathbf b$ ranges over $\R^3$.
\item Determine the number of free variables in $\mathbf A\mathbf x=\mathbf0$ and find the general solution in parametric form.
\item For
\[
\mathbf x_1=\begin{pmatrix}-3\\-3\\-3\end{pmatrix},
\qquad
\mathbf x_2=\begin{pmatrix}0\\0\\2\end{pmatrix},
\]
determine whether each vector belongs to the range of $T_{\mathbf A}\colon\R^3\to\R^3$, $T_{\mathbf A}(\mathbf x)=\mathbf A\mathbf x$. Justify your answers.
\end{problemsubparts}
\end{problem}

% 2025 Spring KAIST Introduction to Linear Algebra Midterm, Problem 4
\begin{problem}[2025 Spring \KAIST{} Introduction to Linear Algebra]
\begin{problemsubparts}
\item State the Leibniz formula for the determinant of an $n\times n$ matrix $\mathbf A=(a_{ij})$ in terms of permutations.
\item Using the definition from part \textup{(a)}, compute
\[
\det\begin{pmatrix}
1&0&1\\
2&3&-1\\
-1&2&0
\end{pmatrix}.
\]
\item For $1\le i\le n$, state the cofactor expansion of $\det(\mathbf A)$ along the $i$th row.
\item Using cofactor expansion along suitable rows or columns, compute
\[
\det\begin{pmatrix}
0&1&0&0&0\\
2&0&-1&4&5\\
0&0&0&2&-3\\
1&0&0&1&2\\
0&1&0&0&-1
\end{pmatrix}.
\]
At each step, indicate the row or column used.
\item Using elementary row operations, find the inverse of
\[
\mathbf B=
\begin{pmatrix}
1&1&1&-1\\
1&1&-1&1\\
1&-1&1&1\\
-1&1&1&1
\end{pmatrix}.
\]
\end{problemsubparts}
\end{problem}

% 2025 Spring KAIST Introduction to Linear Algebra Midterm, Problem 5
\begin{problem}[2025 Spring \KAIST{} Introduction to Linear Algebra]
Recall that the adjugate matrix of a square matrix is the transpose of its cofactor matrix.
\begin{problemsubparts}
\item Find $\adj(\mathbf A)$ and $\mathbf A^{-1}$ for
\[
\mathbf A=
\begin{pmatrix}
2&0&3\\
0&3&2\\
-2&0&-4
\end{pmatrix}.
\]
\item Let $\mathbf B\in\R^{n\times n}$ with $n>1$. Prove that
\[
\det\bigl(\adj(\mathbf B)\bigr)=\det(\mathbf B)^{n-1}.
\]
\end{problemsubparts}
\end{problem}

% 2025 Spring KAIST Introduction to Linear Algebra Midterm, Problem 6
\begin{problem}[2025 Spring \KAIST{} Introduction to Linear Algebra]
\begin{problemsubparts}
\item Find the reduced row echelon form of the $n\times n$ matrix
\[
\begin{pmatrix}
0&1&1&\cdots&1&1\\
1&0&1&\cdots&1&1\\
1&1&0&\cdots&1&1\\
\vdots&\vdots&\vdots&\ddots&\vdots&\vdots\\
1&1&1&\cdots&0&1\\
1&1&1&\cdots&1&1
\end{pmatrix}.
\]
\item Find the minimum possible number of zero entries in an invertible $n\times n$ matrix whose entries belong to $\{0,1\}$. Justify your answer.
\end{problemsubparts}
\end{problem}

% 2025 Spring KAIST Introduction to Linear Algebra Midterm, Problem 7
\begin{problem}[2025 Spring \KAIST{} Introduction to Linear Algebra]
\begin{problemsubparts}
\item Find the characteristic polynomial of
\[
\mathbf A=
\begin{pmatrix}
0&0&0&0&-a_0\\
1&0&0&0&-a_1\\
0&1&0&0&-a_2\\
0&0&1&0&-a_3\\
0&0&0&1&-a_4
\end{pmatrix}.
\]
\item Find a matrix whose characteristic polynomial is
\[
2-3\lambda+4\lambda^2-\lambda^3+2\lambda^4+\lambda^5.
\]
\end{problemsubparts}
\end{problem}

% 2025 Spring KAIST Introduction to Linear Algebra Midterm, Problem 8
\begin{problem}[2025 Spring \KAIST{} Introduction to Linear Algebra]
Let $\lambda$ be an eigenvalue of $\mathbf A\in\R^{n\times n}$ with corresponding eigenvector $\mathbf x$.
\begin{problemsubparts}
\item If $\mathbf P\in\R^{n\times n}$ is invertible, prove that $\lambda$ is an eigenvalue of $\mathbf P^{-1}\mathbf A\mathbf P$ and find a corresponding eigenvector.
\item If $\mathbf A$ is invertible, prove that $\lambda^{-1}$ is an eigenvalue of $\mathbf A^{-1}$ and find a corresponding eigenvector.
\end{problemsubparts}
\end{problem}

% 2025 Spring KAIST Introduction to Linear Algebra Midterm, Problem 9
\begin{problem}[2025 Spring \KAIST{} Introduction to Linear Algebra]
Let $\mathbf A\in\R^{2\times2}$ be symmetric.
\begin{problemsubparts}
\item Prove that $\mathbf A$ has a real eigenvalue.
\item If $\mathbf A$ has one repeated eigenvalue, prove that $\mathbf A$ is a scalar multiple of $\mathbf I_2$.
\item If $\mathbf A$ has two distinct eigenvalues, prove that eigenvectors corresponding to distinct eigenvalues are orthogonal.
\end{problemsubparts}
\end{problem}

% 2025 Spring KAIST Introduction to Linear Algebra Midterm, Problem 10
\begin{problem}[2025 Spring \KAIST{} Introduction to Linear Algebra]
Let
\[
\mathbf A=
\begin{pmatrix}
-\frac1{\sqrt2}&-\frac1{\sqrt2}\\
\frac1{\sqrt2}&-\frac1{\sqrt2}
\end{pmatrix},
\qquad
\mathbf B=
\begin{pmatrix}
\frac1{\sqrt2}&\frac1{\sqrt2}\\
\frac1{\sqrt2}&-\frac1{\sqrt2}
\end{pmatrix}.
\]
\begin{problemsubparts}
\item Prove that $\mathbf A$ is the standard matrix of a rotation of $\R^2$, and find its rotation angle $\theta\in[0,\pi]$.
\item Prove that $\mathbf B$ is the standard matrix of a reflection of $\R^2$, and find its reflection line.
\item Determine whether $\mathbf A\mathbf B$ is the standard matrix of a rotation or a reflection. If it is a rotation, find its angle $\theta\in[0,\pi]$; if it is a reflection, find its reflection line.
\end{problemsubparts}
\end{problem}

% 2025 Spring KAIST Introduction to Linear Algebra Midterm, Problem 11
\begin{problem}[2025 Spring \KAIST{} Introduction to Linear Algebra]
Let $\mathbf A\in\R^{m\times n}$.
\begin{problemsubparts}
\item Prove that $\mathbf A^{\mathsf T}\mathbf A$ is symmetric.
\item If $\tr(\mathbf A^{\mathsf T}\mathbf A)=0$, prove that $\mathbf A=\mathbf0$.
\item Prove that
\[
\Null(\mathbf A)=\Null(\mathbf A^{\mathsf T}\mathbf A).
\]
\end{problemsubparts}
\end{problem}

% 2025 Spring KAIST Introduction to Linear Algebra Midterm, Problem 12
\begin{problem}[2025 Spring \KAIST{} Introduction to Linear Algebra]
Let
\[
\mathbf A=
\begin{pmatrix}
0&1&0\\
0&0&1\\
1&0&0
\end{pmatrix}.
\]
Prove or disprove that $\mathbf A$ is the standard matrix of a rotation of $\R^3$. If it is, find the rotation axis and rotation angle.
\end{problem}

% 2025 Spring KAIST Introduction to Linear Algebra Final, Problem 1
\begin{problem}[2025 Spring \KAIST{} Introduction to Linear Algebra]
Let
\[
\mathbf v_1=\begin{pmatrix}2\\5\\1\end{pmatrix},
\qquad
\mathbf v_2=\begin{pmatrix}1\\2\\3\end{pmatrix},
\qquad
V=\Span(\mathbf v_1,\mathbf v_2)\subseteq\R^3.
\]
\begin{problemsubparts}
\item Determine whether $\mathbf v_1$ and $\mathbf v_2$ are linearly independent, and find $\dim V$.
\item Find an orthonormal basis for $V$.
\item Extend the orthonormal basis found in part \textup{(b)} to an orthonormal basis for $\R^3$.
\end{problemsubparts}
\end{problem}

% 2025 Spring KAIST Introduction to Linear Algebra Final, Problem 2
\begin{problem}[2025 Spring \KAIST{} Introduction to Linear Algebra]
Let
\[
\mathbf A=
\begin{pmatrix}
1&2&1&0\\
0&1&-1&2\\
2&5&0&2
\end{pmatrix}.
\]
\begin{problemsubparts}
\item Compute $\rk(\mathbf A)$ and $\dim\Null(\mathbf A)$.
\item Find a basis for $\Col(\mathbf A)$ and a basis for $\Null(\mathbf A)$.
\item Compute $\rk(\mathbf A^{\mathsf T})$ and $\dim\Null(\mathbf A^{\mathsf T})$.
\item Find a basis for $\Col(\mathbf A^{\mathsf T})$ and a basis for $\Null(\mathbf A^{\mathsf T})$.
\end{problemsubparts}
\end{problem}

% 2025 Spring KAIST Introduction to Linear Algebra Final, Problem 3
\begin{problem}[2025 Spring \KAIST{} Introduction to Linear Algebra]
Let $U,W\subseteq\R^n$ be subspaces, and define
\[
U+W=\{\mathbf u+\mathbf w\mid \mathbf u\in U,\ \mathbf w\in W\}.
\]
\begin{problemsubparts}
\item Prove that $U+W$ and $U\cap W$ are subspaces of $\R^n$.
\item Prove that
\[
\dim(U+W)=\dim U+\dim W-\dim(U\cap W).
\]
\item For
\[
U=\Span\left(
\begin{pmatrix}1\\2\\0\\1\end{pmatrix},
\begin{pmatrix}0\\1\\1\\1\end{pmatrix}
\right),
\]
\[
W=\Span\left(
\begin{pmatrix}1\\0\\1\\0\end{pmatrix},
\begin{pmatrix}1\\4\\2\\3\end{pmatrix},
\begin{pmatrix}1\\8\\3\\6\end{pmatrix}
\right)\subseteq\R^4,
\]
compute $\dim(U+W)$ and find a basis for $U+W$.
\end{problemsubparts}
\end{problem}

% 2025 Spring KAIST Introduction to Linear Algebra Final, Problem 4
\begin{problem}[2025 Spring \KAIST{} Introduction to Linear Algebra]
Let $\mathbf A,\mathbf B\in\R^{n\times n}$.
\begin{problemsubparts}
\item Prove that
\[
\Null(\mathbf B)\subseteq\Null(\mathbf A\mathbf B).
\]
\item Suppose $(\mathbf v_1,\dots,\mathbf v_p)$ is a basis for $\Null(\mathbf B)$ and $(\mathbf v_1,\dots,\mathbf v_{p+q})$ is a basis for $\Null(\mathbf A\mathbf B)$. Prove that
\[
\mathbf B\mathbf v_{p+1},\dots,\mathbf B\mathbf v_{p+q}
\]
are linearly independent.
\item Prove that
\[
\dim\Null(\mathbf A\mathbf B)\le \dim\Null(\mathbf A)+\dim\Null(\mathbf B)
\]
and
\[
\rk(\mathbf A\mathbf B)\ge \rk(\mathbf A)+\rk(\mathbf B)-n.
\]
\item Prove that $\rk(\mathbf A)=k$ if and only if
\[
\mathbf A=\mathbf u_1\mathbf v_1^{\mathsf T}+\cdots+\mathbf u_k\mathbf v_k^{\mathsf T}
\]
for some linearly independent vectors $\mathbf u_1,\dots,\mathbf u_k\in\R^n$ and some linearly independent vectors $\mathbf v_1,\dots,\mathbf v_k\in\R^n$.
\item Prove that
\[
\rk(\mathbf A+\mathbf B)\le \rk(\mathbf A)+\rk(\mathbf B)
\]
and
\[
\dim\Null(\mathbf A+\mathbf B)\ge \dim\Null(\mathbf A)+\dim\Null(\mathbf B)-n.
\]
\end{problemsubparts}
\end{problem}

% 2025 Spring KAIST Introduction to Linear Algebra Final, Problem 5
\begin{problem}[2025 Spring \KAIST{} Introduction to Linear Algebra]
Let $\mathbf A\in\R^{m\times n}$ and $\mathbf b\in\R^m$.
\begin{problemsubparts}
\item A least-squares solution $\widehat{\mathbf x}$ of $\mathbf A\mathbf x=\mathbf b$ is an exact solution of $\mathbf A\widehat{\mathbf x}=\widehat{\mathbf b}$ for some $\widehat{\mathbf b}\in\R^m$. Express $\widehat{\mathbf b}$ using the orthogonal projection onto one of the fundamental spaces of $\mathbf A$.
\item If $\mathbf A$ has full column rank, determine the number of least-squares solutions of $\mathbf A\mathbf x=\mathbf b$ and justify your answer.
\item If $\mathbf A$ does not have full column rank, determine the number of least-squares solutions of $\mathbf A\mathbf x=\mathbf b$ and justify your answer.
\item Write the normal equations for $\mathbf A\mathbf x=\mathbf b$ and explain why every solution of the normal equations is a least-squares solution of $\mathbf A\mathbf x=\mathbf b$.
\item Find all least-squares solutions of $\mathbf A\mathbf x=\mathbf b$ for
\[
\mathbf A=
\begin{pmatrix}
2&1\\
4&2\\
-2&1
\end{pmatrix},
\qquad
\mathbf b=\begin{pmatrix}3\\2\\1\end{pmatrix}.
\]
\end{problemsubparts}
\end{problem}

% 2025 Spring KAIST Introduction to Linear Algebra Final, Problem 6
\begin{problem}[2025 Spring \KAIST{} Introduction to Linear Algebra]
\begin{problemsubparts}
\item Let $\mathbf P\in\R^{n\times n}$ satisfy $\rk(\mathbf P)=k$, $\mathbf P^2=\mathbf P$, and $\mathbf P^{\mathsf T}=\mathbf P$. Prove that $\mathbf P$ is the standard matrix of the orthogonal projection onto some $k$-dimensional subspace $W\subseteq\R^n$.
\item Classify $\mathbf P$ from part \textup{(a)} as positive semidefinite, negative semidefinite, or indefinite.
\item Let $\mathbf H\in\R^{n\times n}$ satisfy $\mathbf H^{\mathsf T}=\mathbf H$, and suppose $\mathbf I_n-\mathbf H$ has rank $1$ and has eigenvalue $2$. Prove that $\mathbf H$ is the standard matrix of reflection across the hyperplane $\mathbf x^\perp$ for some nonzero $\mathbf x\in\R^n$.
\item Classify $\mathbf H$ from part \textup{(c)} as positive semidefinite, negative semidefinite, or indefinite.
\end{problemsubparts}
\end{problem}

% 2025 Spring KAIST Introduction to Linear Algebra Final, Problem 7
\begin{problem}[2025 Spring \KAIST{} Introduction to Linear Algebra]
\begin{problemsubparts}
\item Using the Gram--Schmidt process, find a $QR$ decomposition of
\[
\mathbf A=
\begin{pmatrix}
1&2&1\\
-1&-2&3\\
0&4&5
\end{pmatrix}.
\]
\item Using reflections (Householder matrices), find a $QR$ decomposition of
\[
\mathbf B=
\begin{pmatrix}
0&3&1\\
0&0&5\\
1&1&1
\end{pmatrix}.
\]
\end{problemsubparts}
\end{problem}

% 2025 Spring KAIST Introduction to Linear Algebra Final, Problem 8
\begin{problem}[2025 Spring \KAIST{} Introduction to Linear Algebra]
Let
\[
\mathbf v_1=\begin{pmatrix}1\\0\\1\end{pmatrix},\qquad
\mathbf v_2=\begin{pmatrix}0\\1\\2\end{pmatrix},\qquad
\mathbf v_3=\begin{pmatrix}1\\1\\0\end{pmatrix},
\]
and let $\mathcal B=(\mathbf v_1,\mathbf v_2,\mathbf v_3)$.
\begin{problemsubparts}
\item Prove that $\mathcal B$ is an ordered basis for $\R^3$.
\item Find the change-of-basis matrix $\mathbf P_{\mathcal B\to\mathcal S}$ from $\mathcal B$ to the standard ordered basis $\mathcal S$ of $\R^3$.
\item For
\[
\mathbf v=\begin{pmatrix}3\\2\\1\end{pmatrix},
\]
compute the coordinate vector $[\mathbf v]_{\mathcal B}$.
\item Express $\mathbf v$ in terms of $\mathbf P_{\mathcal B\to\mathcal S}$ and $[\mathbf v]_{\mathcal B}$.
\end{problemsubparts}
\end{problem}

% 2025 Spring KAIST Introduction to Linear Algebra Final, Problem 9
\begin{problem}[2025 Spring \KAIST{} Introduction to Linear Algebra]
For $\mathbf x\in\R^3$, let $\mathbf x'$ be obtained by rotating $\mathbf x$ by $45^\circ$ counterclockwise about the $z$-axis, viewed from the positive $z$-axis toward the origin. Let $\mathbf x''$ be obtained by rotating $\mathbf x'$ by $30^\circ$ counterclockwise about the $y$-axis, viewed from the positive $y$-axis toward the origin.
\begin{problemsubparts}
\item Find the matrix $\mathbf A$ such that $\mathbf A\mathbf x=\mathbf x''$ for every $\mathbf x\in\R^3$.
\item Determine whether $\mathbf A$ is similar to a Householder matrix.
\end{problemsubparts}
\end{problem}

% 2025 Spring KAIST Introduction to Linear Algebra Final, Problem 10
\begin{problem}[2025 Spring \KAIST{} Introduction to Linear Algebra]
\begin{problemsubparts}
\item Find a singular value decomposition of
\[
\mathbf A=
\begin{pmatrix}
3&2&2\\
2&3&-2
\end{pmatrix}.
\]
\item Find a polar decomposition of
\[
\mathbf B=
\begin{pmatrix}
3&0\\
4&5
\end{pmatrix}.
\]
\end{problemsubparts}
\end{problem}

% 2025 Spring KAIST Introduction to Linear Algebra Final, Problem 11
\begin{problem}[2025 Spring \KAIST{} Introduction to Linear Algebra]
\begin{problemsubparts}
\item Let
\[
\mathbf A=
\begin{pmatrix}
-1&0&1\\
-3&2&4\\
1&-1&-2
\end{pmatrix},
\qquad
\mathbf B=
\begin{pmatrix}
1&2&-2\\
2&1&-2\\
2&2&-3
\end{pmatrix}.
\]
Determine whether $\mathbf A$ is similar to $\mathbf B$. If so, find an invertible matrix $\mathbf P$ such that
\[
\mathbf A=\mathbf P^{-1}\mathbf B\mathbf P.
\]
\item Let
\[
\mathbf C=
\begin{pmatrix}
1&3&-3\\
0&1&0\\
0&-3&4
\end{pmatrix},
\qquad
\mathbf D=
\begin{pmatrix}
2&1&1\\
1&2&1\\
1&1&2
\end{pmatrix}.
\]
Determine whether $\mathbf C$ is similar to $\mathbf D$. If so, find an invertible matrix $\mathbf Q$ such that
\[
\mathbf C=\mathbf Q^{-1}\mathbf D\mathbf Q.
\]
\end{problemsubparts}
\end{problem}
